\documentclass[../main.tex]{subfiles}
\graphicspath{{\subfix{../}}}

\begin{document}
\textit{Izotóptérkép, atommagok tömege, mérete, kötési energiája. Energia és tömeg ekvivalencia, tömegdefektus. A cseppmodell és a félempirikus kötési formula.
Maghasadás, magfúzió, radioaktivitás. Sugárzás és anyag kölcsönhatása: Bethe–Bloch-formula. Radioaktív bomlások, magátalakulások. Elemi részecskék és alapvető kölcsönhatások. Kísérleti eszközök.}

\subsection{Izotóptérkép, atommagok tömege, mérete, kötési energiája}
Az atommagnak két lényeges jellemzője van, amelyek mindegyike többé-kevésbé egy elemi érték egész számú többszöröseként változik: az elektromos töltés, és a tömeg. Az
atommag elektromos töltése pontosan egész számú (Z) többszöröse a proton töltésnek,
a tömeg pedig közelítőleg egész számú (A) többszöröse a proton tömegének. Az atomma-
gok tömegének mérésekor hamarosan rájöttek, hogy vannak azonos töltésű, de különböző
tömegű atommagok. HEVESY György (1885-1966, kémiai Nobel-díj 1944) mutatta meg,
hogy ezen atommagokkal alkotott atomok kémiailag teljesen azonos módon viselkednek,
és ezért az elemek periódusos rendszerének azonos helyén kell legyenek. Ezeket az atommagokat izotópoknak nevezték el (izo toposz - görögül azonos helyet jelent). Bár az
atomok kémiailag azonosak, atommagjaiknak mégis nagyon különböző tulajdonságaik
lehetnek: az egyik bomlik, a másik nem, az egyiknek hosszabb a felezési ideje a másiké rövidebb, az egyik nagyobb energiával bomlik, a másik kisebbel. A hidrogénatom
kivételével A > Z, azaz különböző számok. Az egyik legizgalmasabb kérdést az jelentette, hogy miből ered az atommag tömege, ill. az elektromos töltése? Az akkor ismert
építőelemekből Rutherford úgy képzelte el az atommagok összetételét, hogy az atommagokban van A db proton, és A - Z elektron, és ezeket a vonzó Coulomb-kölcsönhatás
tartja össze. Mivel az elektronok tömege sokkal kisebb a protonokénál, ezért az atommag
tömege jó közelítéssel az A proton tömegével egyezik meg. Az elektronok töltése pedig
éppen semlegesíti A - Z proton töltését, és így az atommag töltése pontosan Z proton
töltésével lesz egyenlő. Ez a modell már kezdetben is felvetett egy megválaszolatlan kérdést: mi dönti el, hogy egy elektron az atommag körül ”külső” elektronként keringjen,
vagy pedig ”bebújjon” az atommagba, és ott legyen?

Rutherford, és fiatal munkatársa James CHADWICK (1881-1974, fizikai Nobel-díj 1935) éveken keresztül keresték a ”kis
hidrogénatomot”, azaz egy olyan részecskét, amelyet ugyanúgy egy proton és egy elektron alkot, mint a hidrogénatomot, csakhogy ez az elektron most nem ”kint” kering, hanem a
proton mellé közel bújik. Ennek elektromosan semleges részecskének kellene lenni, tömege alig valamivel nagyobb, mint a proton tömege, és olyan kis mérete is kellene legyen,
mint egy protonnak. Rutherford és munkatársai felhagytak ezekkel a kutatásokkal, amikor az 1920-as években a kvantummechanika kimutatta, hogy a Coulomb-kölcsönhatás
képtelen lenne egy elektront egy atommag roppant kis térfogatában kötve tartani! Ilyen modellel nem lehetett az atommagok szerkezetét a kvantummechanikával összhangban
értelmezni. Ezt a nyugtalanító helyzetet oldotta meg 1932-ben a neutron felfedezése.

1932-ben Fréderic JOLIOT-CURIE (1900-1958, Nobel-díj 1935) feleségével Irène CURIE-vel
(1897-1956, Nobel-díj 1935) Franciaországban, valamint James CHADWICK (1891-1974, Nobel-díj 1935) Angliában az alfa-részecskék által berilliumból kiváltott különleges,
nagy áthatolóképességű sugárzás tulajdonságait tanulmányozták. A sugárzást sem elektromos, sem mágneses térrel nem lehetett eltéríteni, és az anyag elég vastag rétegein is át
tudott hatolni. Ezeknek alapján kézenfekvő volt az a feltételezés, hogy ez a sugárzás is valamilyen gamma-sugárzás, azaz elektromágneses természetű. Voltak azonban olyan jelek
is, amelyek ennek a sugárzásnak a különleges létére utaltak: a sugárzást az anyagok másképpen nyelték el, mint az addig ismert gamma-sugárzásokat. Ezért a három fenti kutató
úgy döntött, hogy e különleges sugárzásnak az anyaggal való kölcsönhatását ”közvetlenül” is tanulmányozzák. Joliot-Curiék figyeltek meg először, hogy ködkamrába helyezettI meant that animals have much bigger C vitamin reserves than we do but again, no shade throwing at all
gázok atommagjai néha erősen meglökődnek a berillium-sugárzás hatására. Joliot-Curiék
ezt a Compton-szóródáshoz hasonló jelenségnek tulajdonították. Ahogyan a Compton-jelenségnél egy gamma-foton ütközik egy elektronnal, és azt meglöki, ugyanúgy - gondolta
Joliot-Curie - egy nagy energiájú gamma-foton itt egy atommagot lök meg. A ködkamrában megfigyelt nyom adataiból a meglökött atommag mozgási energiája meghatározható
volt, ennek ismeretében pedig ki lehetett számítani annak a gamma-fotonnak az energiáját, amely ilyen mértékben meg tudta lökni az atommagot. A számítás
meglepő eredményt adott: a Be ”gamma-sugárzásának” energiájára a korábban ismert
gamma-energiák sokszorosa jött ki! Még meglepőbb volt viszont az, hogy ha kicserélték a
ködkamra töltőgázát, a meglökésekből kiszámított gamma-energia teljesen másnak adódott.
ami nyilvánvalóan képtelenség! A problémát Chadwick oldotta meg: feltételezte, hogy a
Be-sugárzás nem gamma-sugárzás, hanem egy addig ismeretlen, elektromosan semleges
részecske ütközik a gáz atommagjaival. Az alapvető különbség a kettő között az, hogy
a gamma-sugárzás fotonjainak a nyugalmi tömege nulla, Chadwick pedig megengedte,
hogy feltételezett részecskéjének legyen valamekkora nyugalmi tömege is. Emiatt az útközéseket jellemző mechanikai mennyiségek - a mozgási energia ($E$) és a lendület ($p$)
- a Chadwick-féle részecskénél $E = \frac{p^2}{2m}$ alakban függnek össze, míg a gamma-fotonnál
$E = pc$ alakban. Természetesen, a Chadwick-féle részecskénél két ismeretlen paramétert
kell meghatározni, az $E$-t és az $m$-et. Ezek viszont két mérésből - két különböző tömegű
gázzal ütköztetve - már meghatározhatók. Chadwick több különböző gázzal is ütköztette a Be-sugárzás részecskéit, és ellentmondásmentes eredményeket kapott:
az új semleges részecske tömege mindig a proton tömegével nagyjából egyezőnek adódott.
Az új részecske a neutron nevet kapta.

A neutron felfedezése után az atommagok összetétele is világossá vált: az atommagot protonok és neutronok alkotják. Ezeket közös névvel nukleonoknak nevezzük. A
protonok száma, amelyet szokásosan Z-vel jelölünk, az elem rendszáma. A protonok
és a neutronok számának az összegét az atommag tömegszámának nevezzük, és A-val jelöljük.

\subsubsection*{Az atommagok jelölése, izotóp, izóbár}
Egy atommagot (nuklidot) a benne lévő protonok Z számával (rendszám) és az összes
részecskék A számával (tömegszám), valamint az elem vegyjelével szokás jelölni. Például
az urán-atommagban 92 db proton, és összesen 238 db részecske (nukleon) van. Ezért a
jelölése: $^{238}_{92}U$. Nyilván a neutronok N száma ebből a két adatból kiszámítható:
\begin{equation*}
    N = A - Z
\end{equation*}
\textbf{Definíció:} \textit{
     Azonos Z rendszámú, de különböző A tömegszámú atommagokat izotópoknak, azonos A tömegszámú, de különböző Z rendszámú atommagokat pedig izobároknak nevezünk.
    Ritkábban használják még az azonos N neutronszámú, de különböző A tömegszámú atommagok elnevezésére az izotónok elnevezést.
}

Például az $^{238}_{92}U$ és a $^{235}_{92}U$ az urán két izotópja (azonos a rendszámuk), a $^{40}_{19}K$ és a
$^{40}_{20}Ca$ atommagok pedig izobárok (azonos a tömegszámuk).

\subsubsection{Atommagok mérete}
A Rutherford-kísérlet csak felső korlátot tudott adni az atommagok méretére; megállapította, hogy az atommagok legalább tízezerszer kisebb sugárúak, mint maga az atom. Az atommag sugarának pontos meghatározása azonban nem egyszerű feladat, ehhez a Rutherford által használt néhány MeV energiájú alfa-sugárzásnál nagyobb energiájú — pontosabban rövidebb hullámhosszúszú — részecskék szükségesek. A későbbi kísérletek alapján kiderült, hogy az atommag nem egy éles határvonalal rendelkező objektum, amelynek sugara pontosan definiálható lenne. Az atommagban lévő anyagsűrű eloszlására többé‑kevésbé jó modellt ad a két paraméteres Fermi‑függvény:
\begin{equation}
\rho(r)=\rho_0\,\frac{1}{1+\exp\!\big(\tfrac{r-R}{d}\big)}\label{eq:fermi}
\tag{1.8}
\end{equation}
Ennek az eloszlásnak az alakja az 1.4. ábrán látható. Az atommag „sugarán” tehát azt az $R$ sugárát érthetjük, amelynél a középponti sűrűség a felére csökken. A felület diffuzitását a $d$ paraméter mutatja meg.
\begin{figure}[H]
    \centering
    \includegraphics[width=0.6\textwidth]{imgs/13-tetel/nuclear_density_distribution.png}
    \caption{Az atommag sűrűségeloszlása a sugár függvényében a Fermi-eloszlás szerint.}
\end{figure}

\subsubsection*{Atommagok sugarának meghatározása}
Mivel az atommag roppant kicsiny, a sugár meghatározásának egyik módja az, hogy mikro-részecskéket (protonokat, elektronokat, neutronokat, alfa-részecskéket stb.) szóratunk az atommagon, és a szórásképből következtetünk az atommag méretére. Itt mindenképpen figyelembe kell vennünk a részecskék hullámtermészetét: mérésünk felbontóképességét az alkalmazott részecske de\,Broglie-hullámhossza korlátozza. A de\,Broglie-hullámhossz a jólismert összefüggéssel adható meg:
\[
\lambda=\frac{h}{p},
\]
ahol $h$ a Planck‑állandó, $p$ pedig a részecske lendülete (impulzusa). A de\,Broglie-hullámhossznál sokkal kisebb méretű objektumok már nem figyelhetők meg; ennek megfelelően a vizsgálandó tárgy felbontása csak akkor lehetséges, ha a részecske hullámhossza nagyjából akkora vagy kisebb, mint a tárgy jellemző mérete.

Itt megjegyzendő, hogy Rutherford kísérletei idején a kísérletek kiértékelésében a klasszikus mechanika eszközeit használta, hiszen a de\,Broglie-hullámhosszról és a kvantummechanikáról még nem rendelkezett megfelelő elméleti háttérrel. Ennek ellenére a természet „kegyes” volt abban az értelemben, hogy az általa használt alfa‑részecskék nagy tömege miatt hullámhosszuk az atommagok mérettartományába esett, így ezekkel már közvetlen információ nyerhető volt az atommag kiterjedéséről. Ugyanakkor a Rutherford‑féle klasszikus hatáskeresztmetszet‑számítás és a kvantummechanikai megközelítés nem mindig ad pontosan ugyanazt az eredményt; ez részben a Coulomb‑potenciál sajátosságaiból adódik.

További komplikációt okoz, hogy a különböző részecskék más és más módon lépnek kölcsönhatásba az atommaggal. Az elektronok nem vesznek részt az atommagot összetartó erős kölcsönhatásban, ezért az elektronszórás elsősorban az atommagon belüli töltéseloszlást (a protonok eloszlását) vizsgálja. Ha neutronokkal szóratunk, azok a Coulomb‑kölcsönhatásra nem érzékenyek, így információt adnak az atommag belső anyageloszlásáról (a nukleonok eloszlásáról). Ezek alapján külön beszélhetünk az atommag „töltés” sugaráról és a nukleáris (anyagi) sugaráról.

A következőkben röviden áttekintünk néhány mérési módszert és a belőlük levont következtetéseket.

\textbf{Elektronszóródás.} Az atommagokon belüli töltéseloszlást előszőr R. HOFSTÄDTER (1915-1990, Nobel-díj 1961) az 1950-es évek második felében.
Az ő módszere az volt, hogy nagy energiájú ($E > 100 \text{MeV}$
\footnote{
    A makroszkopikus világunkban használt SI-energiaegység ($J$) az atomok és a mikrorészecskék világában alkalmatlan, mivel az ott előforduló energiák ennek
töredékei, csak a J két számjegyű negatív kitevős hatványaival kifejezhetők. Ezért a
mikrofizikában az elektronvoltot ($eV$) használjuk energiaegységként. Egy elektronvolt
mozgási energiát kap egy elemi töltéssel ($1,6 \cdot 10^{-19} \text{C}$) rendelkező részecske (pl. elektron,
proton), ha 1 V feszültség gyorsítja. Ezért az átváltás a két energiaegység között:
\begin{equation*}
1 \text{eV} = 1,6 \cdot 10^{-19} \text{J}.
\end{equation*}
Természetesen, használjuk még az eV többszöröseit, amelyeket a szokásos SI előtagokkal
(kilo-, mega-, giga-, tera stb.) képezünk.
})
elektronokkal bombázta az atommagokat, és a szórt elektronok szögeloszlását mérte. Mivel az ilyen nagy energiájú elektronoknak a de\,Broglie-hullámhossza hullámhossza nagyon rövid,
ezért érzékeny az atommagon belüli töltéseloszlásra. Hofstädter ezekkel a kísérletekkel megállapította, hogy a stabil atommagok belsejében (a legkönnyebb magok kivételével)
a töltések eloszlása jó közelítéssel jól leírható a Fermi-eloszlással, amelynek az $R$ sugárparamétere (amit az atommag töltéssugarának is nevezünk) a tömegszámmal nagyjából az alábbi összefüggést követi:
\begin{equation}
R = r_0 \sqrt[3]{A} \quad \text{ahol} \quad r_0 \approx 1,2 \text{ fm} \quad (1 \text{ fm} = 10^{-15} \text{ m})
\end{equation}
A magok diffuzitása $d$ nagyjából független az atom méretétől ($d \approx 0,5 \text{ fm}$). Az $fm$ az SI rendszer szerint femtométert jelent, de atomfizikusok gyakran "Fermi"-ként hivatkoznak rá,
Enrico Fermi (1901-1954, Nobel-díj 1938) olasz fizikus után.
Későbbi kísérletek kimutatták, hogy ez csak közelítőleg igaz, ugyanis az atomok finomszerkezete miatt kisebb-nagyobb mértékben eltérhetnek a valós töltéssugarak a fenti összefüggéstől.

\textbf{Neutronszórás.} Az elektronokkal a "probléma", hogy van töltésük, így Coulomb-erő miatt csak bizonyos közelségre tudjuk az atommagot megközelíteni. Ez ad egy "effektív" méretet,
hisz a legtöbb interakció, amiben az atom részt vesz, ezzel a Coulomb-térrel történik. De ha a valós, fizikai kiterjedését akarjuk megállapítani az atommagnak, akkor egy olyan részecskét kell
használni, aminek nincs töltése. Erre kiváló választás az ekkor még nemrég felfedezett Neutron. Természetesen a de\,Broglie-hullámhosszra itt is figyelni kell,
ezért itt is nagy energiájú neutronokat használtak ($10 \text{MeV}$.től egészen $1.4 \text{GeV}$-ig). A kísérletek során az derült ki, hogy a neutronokkal mért atommag sugara
kissé nagyobb, mint az elektronszórással mért töltéssugár ($1.3 \text{fm} < r_0 < 1.4 \text{fm}$). Ez arra utal, hogy az atommagban a protonok (töltéssugár) és a neutronok eloszlása (nukleáris sugár) nem teljesen azonos.
Ez a különbség különösen nagy a neutron-gazdag atommagokban. A Nehéz atommagokban az atom felszínén egy úgynevezett neutron-bőr alakulhat ki, ahol a neutronok sűrűsége jelentősen nagyobb, mint a protonoké.
A könnyebb atommagokban pedig egy neutron udvar (neutron halo). Például a $^{181}_{55}Cs$ atommag felszínén kb. $2 \text{fm}$ vastag neutron bőr lehet ami 10 neutronból áll.
A $^6_2He$ atommagban pedig a neutronok olyan lazán kapcsolódnak egymáshoz, hogy az atommag akár a $^{48}Ca$ atommag méretét is elérheti, miközben csak 4 neutronból álló udvara van.

\textbf{Tükörmagok kötési energiája.} Az eddigiektől különböző elven kaphatunk becslést az atommagok sugarára a tükörmagok kötési energiájának összehasonlításával. Tükörmagoknak azokat az atommagokat nevezzük, amelyek egymásból a protonok és a neutronok számának felcserélésével adódnak (pl. \(\,^{13}_6\mathrm{C}\leftrightarrow{}^{16}_7\mathrm{N}\), vagy \(\,^{39}_{20}\mathrm{Ca}\leftrightarrow{}^{39}_{19}\mathrm{K}\) stb.). Ha a tükörmagok kötési energiáját a folyadékcsepp‑modell szerint számítjuk, akkor mindegyik energiaterületagban megegyezik, kivéve a Coulomb‑energia tagot. Egy \(R\) sugarú, egyenletesen töltött gömb Coulomb‑energiája
\begin{equation}
E_c=\frac{3}{5}\,\frac{1}{4\pi\varepsilon_0}\,\frac{(Ze)^2}{R}.
\tag{1.10}
\end{equation}

Foglalkozzunk csak azzal az esettel, amikor a két tükörmag rendszáma csak eggyel tér el! Ekkor az energiakülönbség a Coulomb‑tagból adódóan
\begin{equation}
\Delta E_c=\frac{3}{5}\,\frac{1}{4\pi\varepsilon_0}\,\frac{e^2}{R}\big(Z^2-(Z-1)^2\big)
=\frac{3}{5}\,\frac{1}{4\pi\varepsilon_0}\,\frac{e^2}{R}\,(2Z-1).
\tag{1.11}
\end{equation}

Nyilván
\begin{equation}
A = Z + N = Z + (Z-1) = 2Z-1,
\tag{1.12}
\end{equation}
ahol \(A\) a tömegszám. Tegyük fel továbbá, hogy a magsugár az
\[
R = r_0\,A^{1/3}
\]
alakot követi; behelyettesítve és az \(A\)-t a nevezőbe írva, akkor kapjuk:
\begin{equation}
\Delta E_c = \frac{3}{5}\,\frac{1}{4\pi\varepsilon_0}\,\frac{e^2}{r_0}\,A^{2/3}.
\tag{1.13}
\end{equation}

A kötési energia különbség kísérletileg több úton is megmérhető. Az egyik lehetőség az, hogy a tükörmagok egyike radioaktív és például béta‑bomlással elbomlik a másik magra; ilyenkor a bomláskor kiszabaduló részecske (elektron vagy pozitron) maximális energiájából a két atommag energiakülönbsége — azaz \(\Delta E_c\) — meghatározható. Az energiakülönbség meghatározásának másik módja a magreakciók kinyomozása. Például, ha a \(\,^{11}_5\mathrm{B}\)-et protonokkal bombázzuk, előfordulhat, hogy a proton kilök egy neutront a magból, miközben a maga fogva marad, $^{13}_6\mathrm{C}$ atommagot hozva létre.
Ez a reakció csak akkor mehet végbe, ha a bombázó proton energiának legalább $\Delta E_c$ energiájuk volt.
A reakció ”energiaküszöbét” meghatározva, $\Delta E_c$ megmérhető. Ha már $\Delta E_c$ megvan, $Z$ ismert, és így az $R$ magsugár
meghatározható. Ezzel a módszerrel is nyilván a mag töltéssugarát határozzuk meg.
Az ilyen típusú mérésekből $r_0 = 1.22 \text{ fm}$. Ez láthatóan szép összhangban van a más módszerekkel kapott értékekkel.

\subsubsection{Atomok töltése}
z atommagok töltése \(Ze\), ahol \(Z\) az atommagban lévő protonok száma (rendszám), \(e\) pedig az elemi töltés. A töltés meghatározása az atomok vonalas röntgenspektrum keletkezése alapján lehetséges. Gyors elektronokkal bombázva a vizsgálni kívánt anyagot a Coulomb‑taszítás következtében egy kötött elektron kilökődhet az atomból. Kilökés után egy külső pályáról elektron ugorhat a megürült állapotba (1.6. ábra), és ennek következtében jellegzetes röntgen‑foton keletkezik.

Eközben az atom elektromágneses sugárzást (fotont) bocsát ki, amelynek energiája:
\begin{equation}
h\nu = E_n - E_1 = Z^2\cdot I_H\cdot\left(1-\frac{1}{n^2}\right),
\tag{1.14}
\end{equation}
ahol \(I_H\) a hidrogénatom ionizációs energiája. A képlet felírásakor feltételezzük, hogy a kötött elektron az 1s állapotból távozott, és helyébe az \(n\) főkvantumszámú állapotból ugrott be egy másik elektron. A képlet \(Z>1\) elektronokat tartalmazó atomok esetén kissé módosulhat, mert az atommag töltését a többi elektron részben leárnyékolja. Ennek megfelelően az effektív töltés \(Z-z\) formában lép fel, ahol \(z\) az árnyékoló elektronok effektív száma. Így a kibocsátott foton energiája írható:
\begin{equation}
h\nu = (Z-z)^2\cdot I_H\cdot\left(1-\frac{1}{n^2}\right).
\tag{1.15}
\end{equation}

\begin{figure}[H]
    \centering
    \includegraphics[width=0.7\textwidth]{imgs/13-tetel/characteristic_xray.png}
    \caption{Az ábra szemlélteti a karakterisztikus röntgensugárzás keletkezésének lehetséges eseteit: (a) egy gyors elektron kiüti egy kötött elektronot, (b) mindkét elektron eltávozik és lyuk marad, (c) egy külső elektron ugrik be a lyukba és karakterisztikus röntgen‑fotont bocsát ki.}
\end{figure}

A speciális vonalak frekvenciája egyszerűsítve gyakran a következő formában szerepel:
\begin{equation}
\nu = C\cdot (Z - z)^2,
\tag{1.16}
\end{equation}
ahol \(C\) az átmenetre jellemző konstans. Ez a Moseley‑törvény egyszerű alakja; a K‑vonalakra vonatkozóan Henry Moseley (1887–1915) fedezte fel a kísérleti összefüggést, amely azt mutatja, hogy a 1s‑elektronok által kialakított vonalak frekvenciája közel négyzetesen nő a hatékony rendszámmal. A Moseley‑törvény tette lehetővé a rendszám kísérleti meghatározását és a periódusos rendszer rendezésének megerősítését.

\subsubsection{Atommagok tömege}
Az atommagok tömegének ismerete alapvető elméleti és kísérleti jelentőségű. Egy \((A,Z)\) összetételű atommag tömege kisebb, mint az őt alkotó protonok és neutronok tömegének összege:
\begin{equation}
M(A,Z) = Z\,m_{\text{proton}} + (A-Z)\,m_{\text{neutron}} - \Delta m.
\tag{1.17}
\end{equation}
A \(\Delta m\) különbség – a \emph{tömeghiány} – abból adódik, hogy az atommagban lévő részecskék kötött állapotban vannak, és onnan a részecskéket kiszabadítani csak a kötési energia befektetésével lehet. Az Einstein‑féle \(E = mc^2\) összefüggésnek megfelelően a tömeghiány szoros kapcsolatban van a kötési energiával:
\[
E_{\text{kötési}} = \Delta m\,c^2.
\]
Az atommagok tömegének pontos megmérésével tehát az atommagok kötési energiáját is meg lehet határozni.

A részecskék tömegének meghatározására szolgáló módszert tömegspektroszkópiának, a mérésre alkalmas berendezéseket pedig tömegspektrográfoknak, illetve tömegspektrométereknek nevezzük.
\begin{figure}[H]
    \centering
    \includegraphics[width=0.6\textwidth]{imgs/13-tetel/mass_spectrograph.png}
    \caption{A tömegspektrográf felépítése.}
\end{figure}
Az ionforrásban a vizsgálandó anyagot ionizáljuk, majd egy kihúzó feszültséggel az ionokat felgyorsítjuk, nyalábszerűen kilőjük. A nyaláb egy „sebességszelektorba” kerül, amely lényegében egymásra merőleges, homogén elektromos és mágneses térből áll. Az elektromosan töltött részecskékre a ($q\mathbf{vB}$) mágneses Lorentz‑erő, valamint a ($q\mathbf{E}$) elektromos erő hat, egymással ellentétes irányban. Azok a részecskék haladnak át irányváltoztatás nélkül, amelyekre ($qE = qvB$), amiből ($v = \dfrac{E}{B}$), tömegtől és töltéstől függetlenül. Ezért hívják ezt az elrendezést sebességszelektornak. Az így beállított sebességű (esetleg több fajta tömegű részecskét is tartalmazó) nyaláb egy szektor‑mágneses térbe kerül, ahol a Lorentz‑erő hatására a töltött részecskék körpályára állnak. A centripetális erőt a Lorentz‑erő biztosítja, azaz
$\frac{m v^2}{r} = q v B$ ,
amiből \begin{equation} r = \frac{mv}{qB} \tag{1.18} \end{equation} adódik.

Egyes elrendezésekben a sebességszelektor a szektor mágneses tér kezdeti szakaszán van, azaz a sebességszelektorban használt mágneses tér ugyanakkora, mint a részecskéket eltérítő tér. Ekkor kapjuk: 
\begin{equation} 
    r = \frac{m}{q},\frac{E}{B^2}, \qquad\text{illetve}\qquad m = \frac{rqB^2}{E}. \tag{1.19}
\end{equation} A részecskék töltése nyilván az elemi töltés (vagy annak egész számú többszöröse), így ($E$) és ($B$) ismeretében, a pályasugár ($r$) mérésével a tömeg meghatározható.

Bár a tömegspektrográf működési elve egyszerű, a gyakorlati megvalósításnál több fontos körülményt is figyelembe kell venni.
\begin{itemize}
    \item Ahhoz, hogy a detektorba különböző helyekre érkező részecskéket meg lehessen különböztetni egymástól, oda a szétvált nyaláboknak fókuszálva kell érkezni. Az elektromos és mágneses terek a töltött részecskékre elektromágneses lencseként hatnak. Ezeket gondosan kell beállítani, hogy a fókuszpontjuk a detektor síkjába essen. Még így is lesznek „lencsehibák”, amelyek az elérhető felbontást korlátozzák.
    \item A tömegspektrográf „felbontásán” azt a legkisebb ($\Delta m/m$) hányadost értjük, amely tömegkülönbséget a detektorsíkban még meg tudunk különböztetni. A jó tömegspektrográfok ($10^{-6}$) felbontást is el tudnak érni. Ez azt jelenti, hogy ha két tömeg egy milliomod résszel tér el csak egymástól, azt már különböző tömegként észleljük.
    \item Fontos paraméter a „fényerő”, amely lényegében az időegység alatt a tömegspektrográfon áthaladt részecskék számával arányos. Általában kompromisszumot kell kötni a felbontás és a fényerő között: ha a felbontást növeljük, a fényerő csökken, és fordítva.
\end{itemize}

\subsubsection*{Az atomi Tömegegység (Atomic Mass Unit, u)}
Ahhoz, hogy a tömeget milliomodrésznyi pontossággal meghatározzuk, a $1{,}19$ egyenletben szereplő mennyiségeket is legalább ilyen pontosan kellene ismerni. Ez, sajnos, általában nem keresztülvihető. A gyakorlatban a tömegspektrográfot egy adott tömeg kalibrálja, és utána relatív mérésekkel határozzák meg a többi atom tömegeit. Az egyik legpontosabban megmért tömegű atommag a $^{12}\mathrm{C}$. Ezért a nemzetközi megállapodások a $^{12}\mathrm{C}$ atom tömegének $1/12$‑ed részében határozzák meg az atomi tömegegységet (u, amu, atomic mass unit). A $^{12}\mathrm{C}$ atom tömegének tehát definíció szerint pontosan $12{,}00000000\ \mathrm{u}$ (atomi tömegegység). Hangsúlyozzuk, hogy ez a $^{12}\mathrm{C}$ atom (és nem atommag!) tömege, tehát a $6$ szénatom $6$ db elektronjának a tömegét is tartalmazza.  
Értéke a szokásos tömegegységben: $1\ \mathrm{u}=1{,}6605387310^{-27}\ \mathrm{kg}$.  

\textbf{Tömegduplett módszer}  
Tegyük fel, hogy a $^{1}\mathrm{H}$ atom tömegét szeretnénk megmérni. Hiába kalibráltuk a tömegspektroszkópunkat a $^{12}\mathrm{C}$ atomra, a $^{1}\mathrm{H}$ atomot nem tudjuk ugyanazzal a beállítással megmérni,
hiszen a két ion fajlagos töltése kb. tizedenkétes faktorra eltér. Ha pedig a mágneses vagy elektromos teret megváltoztatjuk, akkor ezek beállításának hibája elrontja a mérésünk pontosságát.
Ehelyett állítsuk be a spektroszkópunkat úgy, hogy a $128$ tömegszámú, egyszeresen ionizált részecskék jussanak el a detektorra. Az ionforrásunkban ionizáljunk egy keveréket, amely $\mathrm{C}_9\mathrm{H}_{20}$‑ból (nonán) és $\mathrm{C}_{10}\mathrm{H}_{8}$‑ból (naftalin) áll.
Mindkét molekula molekulatömege $128$. A pontos tömegek persze eltérnek egy kicsit, ezért a $10^{-6}$ különbséget is felbontani képes spektrográfunkban a két molekula‑ion különböző helyre csapódnak. Az eltérésből a tömegkülönbséget meg tudjuk mérni, az abszolút tömegeket viszont nem.
A tömegkülönbség mért értéke:  $\Delta m = 0,093 90032 \pm 0,000 00012 \text{u}$ . (A pontosság itt kb $10^{-6}$.) Másrészt
$\Delta m = m(\mathrm{C}9\mathrm{H}{20}) - m(\mathrm{C}_{10}\mathrm{H}_8) = 12\cdot m(^{1}\mathrm{H}) - m(^{12}\mathrm{C}),$

amiből
\begin{equation} m(^{1}\mathrm{H}) = \frac{m(^{12}\mathrm{C})}{12} + \frac{\Delta m}{12} = 1{,}000000 + \frac{\Delta m}{12} = 1{,}00782503 \pm 0{,}00000001\ \mathrm{u} \tag{1.21} \end{equation}

Vegyük észre, hogy a relatív tömeg-meghatározás miatt itt a hiba már csak $10^{-8}$ nagyságrendű! Ha már a hidrogénatom tömegét elegendően pontosan ismerjük, nagyon sok atom tömegét meg lehet határozni a fenti módszerrel, különböző összetételű etalon‑gázkeverékeket alkalmazva. Például a nitrogén tömegének meghatározásához állítsuk be úgy a tömegspektrométerünket, hogy a 28 ‑as tömegszámú, egyszeresen ionizált ionokat tudjuk detektálni. Az ionforrásba pedig tegyünk $N_2$ gáz és $\mathrm{C}_2\mathrm{H}_4$ (etén) keveréket. A tömegkülönbség mért értéke: $\Delta m{N} = 0{,}025152196 \pm 0{,}000000030\ \mathrm{u}.$ (A hiba itt is nagyságrendileg $10^{-6}$.) Másrészt:
\begin{equation}
    \Delta m{N} = m(\mathrm{C}_2\mathrm{H}_4) - m(\mathrm{N}_2) = 2\cdot m(^{12}\mathrm{C}) + 4\cdot m(^{1}\mathrm{H}) - 2\cdot m(^{14}\mathrm{N})
\end{equation}
amiből
\begin{equation} m(^{14}\mathrm{N}) = m(^{12}\mathrm{C}) + 2\cdot m(^{1}\mathrm{H}) - \frac{\Delta m_{N}}{2} = 14{,}00307396 \pm 0{,}00000002\ \mathrm{u} \tag{1.23} \end{equation}

Nem minden egyes atommag tömegét szükséges tömegspektrométerekkel pontosan megmérni. Ha vannak jól megmért viszonyítási pontjaink, akkor az egyes magreakciókban ill.\ magbomlásokban felszabaduló energia mérésével e viszonyítási pontok alapján nagyon sok egyéb tömege pontosan meghatározható. Tekintsük az A + B → C + D atommag‑reakciót, ahol az A és a B atommagok reakciójából C és D atommagok keletkeznek. A részecskék mozgási energiájának mérésével a reakcióban felszabaduló energiát ( Q ) meghatározhatjuk. Másrészt viszont
\begin{equation*}
Q = c^2 \cdot \left[ m ( A ) + m ( B ) - m ( C ) - m ( D ) \right]
\end{equation*}

Példaként vegyük a 1 H + 14 N → 12 C + 3 H reakciót! Tömegduplett mérésekből tudjuk, hogy m ( 1 H ) = 1,007 8252   u ,   m ( 14 N ) = 14,003 074   u és m ( 3 H ) = 3,016 049   u . A mért reakcióenergia pedig Q = 22,135 5 ± 0,001 0   MeV . Ezekből tehát kapjuk:
\begin{equation}
    m(^{12}\mathrm{N}) = m(^{1}\mathrm{H}) + m(^{14}\mathrm{N}) - m(^{3}\mathrm{H}) - \frac{Q}{c^2} = 12{,}018613 \pm 0{,}000010\ \mathrm{u}.    
\end{equation}
A mérési hiba legnagyobbrészt a $Q$ érték hibájából származik, a másik három tömeg
értéke sokkal nagyobb pontossággal ismert. A $^{12}\mathrm{N}$ atommag radioaktív, felezési ideje
0,01 s. Ez túl rövid ahhoz, hogy a tömegét közvetlen mérésekkel (tömegspektroszkóppal) meghatározhatnánk. A most leírt módon azonban a rövid felezési idejű radioaktív
izotópok tömegét is meghatározhatjuk.

\subsubsection{Atomok kötési energiája}
A jelenleg ismert több mint 2000 atommag két nagy csoportra osztható: a stabil atom-
magok és a radioaktív atommagok csoportjára. Stabilnak nevezzük azokat az atomma-
gokat, amelyek nem bomlanak el, pontosabban amelyek élettartama olyan hosszú, hogy
a jelenlegi eszközeinkkel nem tudjuk a bomlásukat kimutatni (például a proton élettar-
tamáról jelenleg annyi tudunk csak, hogy nagyobb, mint $10^{32}$ év); radioaktívak pedig
azok, amelyek rövidebb-hosszabb idő alatt kimutathatóan elbomlanak.

Az atommagok kötött nukleon-rendszerek, és csak bizonyos kötési energia befektetésével bonthatók szét. Mivel ez az energia nem áll rendelkezésre, ez akadályozza meg,
hogy az atommagok a létrejöttük után azonnal szétesnek összetevőikre.
\newline
\newline
\textbf{Definíció 1.:} \textit{
    Kötési energia: Egy részecske kötési energiáján azt az energiát értjük,
    amit ahhoz kell befektetni, hogy az adott részecskét a kötésből kiszakítsuk, és onnan végtelen távolra vigyük.
}
\newline
\textbf{Definíció 2.:} \textit{
    Az atommag kötési energiája az az energia, amit ahhoz kell befektetni,
    hogy az atommagot teljesen szétbontsuk alkotóelemeire (protonokra és neutronokra), és
    ezeket egymástól végtelen távolra elvigyük.
}
\newline
\newline
Fontos megkülönböztetnünk az energiát a kötési energiától. Egy rendszer teljes energiáját az Einstein-féle $E = m \cdot c^2$ képlet alapján a rendszer tömegének ismeretében lehet
megadni. Itt az energiaskála nullpontja a tömeg- és energia nélküli, ”anyagmentes” állapot. Ezen a skálán minden létező rendszer energiája pozitív.
Tekintsünk példaként egy egyszerű atommagot, a deuteront! Ennek a teljes energiája
a tömegével kifejezhető: $E_d = M_d \cdot c^2$. Bontsuk szét ezt egy protonra és egy neutronra,
és a két részecskét vigyük olyan távol, hogy már ne álljanak egymással kölcsönhatásban!
A második rendszer teljes energiája tehát: $E = m_p \cdot c^2 + m_n \cdot c^2$. A szétbontáshoz B
energiát kellett befektetni, tehát magasabb energiájú állapotot kapunk:
\begin{equation}
    m_p \cdot c^2 + m_n \cdot c^2 = M_d \cdot c^2 + B
\end{equation}
\begin{figure}[H]
    \centering
    \includegraphics[width=0.6\textwidth]{imgs/13-tetel/binding_energy.png}
    \caption{Az atommag energiája és kötési energiája.}
\end{figure}
Gyakran azonban csak az az energiaváltozás érdekel minket, amely egy kötött rendszer létrejöttekor történik. Ilyenkor - megállapodás szerint - ott vesszük fel a rendszer
energiájának a nulla szintjét, amikor a rendszer elemei még végtelen távol vannak egymástól, azaz nincsenek kötésben. Amikor létrejön a kötés, a
rendszer mélyebb energiájú állapotba kerül, tehát ilyen nullapont választása mellett az
energiája negatív lesz ($E < 0$). Ahhoz, hogy onnan kihozzuk – a kötést megszüntessük –,
pozitív energiát ($B$) kell befektessünk. Ezt a pozitív energiát nevezzük kötési energiának.
Más szóval $B = - E$. A kötési energia tehát pozitív, a kötött rendszer energiája viszont
negatív.

Ha az atommagok kötési energiáját úgy vizsgáljuk, hogy a kötési energiát elosztjuk a nukleonok számával, akkor megkapjuk az egy nukleonra jutó átlagos kötési energiát. Ez az érték
jó közelítéssel jellemzi az adott atommag stabilitását. Minél nagyobb ez az érték, annál stabilabb az atommag. Érdekesség, hogy a $Z = 26$ után, kb. azonos a kötési energia nukleononként.
Ez azért van, és a magerő telítődik, így a nagyobb atommagoknál már nem nő tovább a kötési energia nukleononként. A "telítés" azt jelenti, hogy
minden nukleonnak megvan a maximális számú közvetlen szomszédja, amivel kölcsönhatásba léphet. Azért csökken enyhén mégis, mert a proton-proton taszítás egyre nagyobb szerepet játszik a nagyobb atommagoknál.
\begin{figure}[H]
    \centering
    \includegraphics[width=0.8\textwidth]{imgs/13-tetel/binding_energy_per_nucleon.png}
    \caption{Az egy nukleonra jutó átlagos kötési energia az atommag tömegszámának függvényében. Látható, hogy a legstabilabb atommagok a közepes tömegűek, kb. $A= 56$ körül.}
\end{figure}
Az ábrán látható, hogy a $^{56}_{26}\mathrm{Fe}$ környékén van a legnagyobb stabilitás (bár a tényleges maximum a $^{62}_{28}\mathrm{Ni}$-nál van, ez egy ritka izotópja a Nikkelnek).
Ez azt jelenti, hogy a könnyű atommagok összeolvadásával (fúzió) és a nehéz atommagok hasadásával is energia szabadítható fel, mivel mindkét esetben a keletkező atommagok közelebb kerülnek a legstabilabb atommagokhoz.
Továbbá következmény, hogy a a csillagokban a vasnál nehezebb elemek már nem igazán keletkeznek fúzióval, a vasnál nehezebb atomok nagyrészt szupernovák által jönnek létre.

Itt fontos megjegyezni, hogy a "stabilitás" ebben a kontextusban a kötési energiára vonatkozik. A nehezebb atomokban "erősebben" kötődnek a nukleonok, de a proton-proton taszítás nem "telítődik",
az folyamatosan nő, nagyjából $Z^2$-tel, így a nehezebb atommagok egyre instabilabbak lesznek a radioaktív bomlás szempontjából (itt más effektusok is közrejátszanak, mint a magerő hatótávolsága, kvantumhatások, stb.).

\subsubsection{Izotóptérkép}
\begin{figure}[H]
    \centering
    \includegraphics[width=0.8\textwidth]{imgs/13-tetel/isotope_map.png}
    \caption{Az izotóptérkép a nukleonok számának (N) és a protonok számának (Z) függvényében ábrázolja az ismert atommagokat. A fekete pontok stabil magokat jelölnek, a színes pontok pedig a különböző radioaktív bomlásokban részt vevő magokat. A vízszintes irányban béta-bomlás, a függőleges irányban alfa-bomlás történik. \href{https://atomfizika.elte.hu/mrf/2019/vgmrf5_atommag.pdf}{forrás}}
\end{figure}

A Magerő szemléletes tulajdonságai:
\begin{itemize}
    \item A stabil atommagok (fekete) jobbra hajlik: A nagy atommagoknál több a neutron, proton-proton taszítás miatt.
    \item Csipkés a széle: A magok bizonyos számú neutron/proton után instabillá válnak.
    \item Csatornák fekete részek között: Páros neutronszámok stabilabbak - a magerő párkölcsönhatás jellegű.
    \item Spinfüggő (annak ellenére, hogy a $\pi$-mezon spinje 0)
\end{itemize}

Az izotóptérkép alapján látható, hogy a lehetséges atommagok közül csak relatíve kevés stabil. A legtöbb atommag előbb-utóbb valamilyen folyamat során elbomlik:
\begin{itemize}
    \item Radioaktív bomlás (alfa-, béta-, gamma-bomlás)
    \item Kötési energia még negatívabb állapot felé történő átalakulás (pl. hasadás, egyes magreakciók)
    \item Negatív energiájú (kötött) atommag se stabil
    \item Pozitív energiájű atommag nem is létezik (izotóptérkép széle, rezonanciák)
\end{itemize}

Érdekesség viszont, hogy a stabil atommagok száma bizonyos rnedszámok környékén nagyon magas. Pl. a $Z=2, 8, 20, 50, \ldots$ rendszámoknál kiugróan sok stabil atommag található. 
Ezeket a számokat nevezik \textit{mágikus számoknak} az atommagok szerkezetében. A mágikus számoknál a magok különösen stabilak, mivel ezeknél a protonok és neutronok héjai teljesen betöltöttek.
Emiatt a betöltöttség miatt pedig gömbszimmetrikusak és a kvadropol momentumuk közel nulla. Emiatt a magerő egy közel átlagteret alakít ki, ami miatt a magok stabilabbak lesznek.

\subsubsection*{Stabil atommagok völgye}
A stabil atommagokat az energia-felület alján kell keressük, hiszen ezek már nem tudnak alacsonyabb energiájú állapotba elbomlani. Az A tömegszámú atommagok közül a
$Z_{\min} (A)$ rendszám közelébe eső atommagok a stabilak, ill. a leghosszabb felezési idejűek.
Gyakran a Z,A változók helyett a $Z,N$ változópárt használják (itt $N$ a neutronok száma). Természetesen, ezek a változók egymással összefüggenek, hiszen $A = Z + N$.
A stabil atommagok völgyének két oldalán helyezkednek el a radioaktív magok. A
völgy egyik oldalán a negatív béta-bomlók (ott, ahol a rendszám növelése kedvezőbb energetikailag), a völgy másik oldalán az elektronbefogással, ill. pozitron-bomlással bomlók,
hiszen ott a rendszám csökkentésével haladhat az anyag a völgy mélypontja felé. Az
átalakulás után a nukleonok erősebben kötöttek az új atommagban, mint a régiben, az
energiájuk csökken. Ez az energiakülönbség felszabadul az átalakuláskor és a környezetbe távozik (a bomláskor keletkező részecskék elviszik). Az atommagok energia-felülete
segítségével tehát a radioaktív bomlások is értelmezhetővé válnak. A sokféle radioaktív
bomlás tehát egyetlen közös okra vezethető vissza: az anyag vándorlására a nukleáris
energia-felület legmélyebb energiájú állapota felé.

\subsection{Energia és tömeg ekvivalencia, tömegdefektus}
Láttuk korábban, hogy az atommagok tömege kisebb mint az alkotóelemeinek tömegének összege. A protonok és neutronok tömegét a tömeg-spektrométerrel külön-külön is meg tudjuk mérni,
és ezeket összerakva, és az Einstein-féle energia-tömeg ekvivalenciát $E = mc^2$, meg tudjuk határozni a közési energiát:
\begin{equation}
    \begin{aligned}
        &E = m c^2 \qquad \text{Rendszer teljes energiája} \\
        &M = Z \cdot m_{\text{proton}} + N \cdot m_{\text{neutron}} - \Delta m \qquad \text{Atommag tömege} \\
        &\rightarrow E_{\text{kötési}} = \Delta m \cdot c^2 \qquad \text{Atommag kötési energiája}
    \end{aligned}
\end{equation}
Ha pedig átrendezük a tömeghiányra vonatkozó egyenletet, akkor megkapjuk a tömegdefektus képletét:
\begin{equation}
    - \Delta m = M - \left( Z \cdot m_{\text{p}} + N \cdot m_{\text{n}} \right)
\end{equation}
Tehát a kötési energia:
\begin{equation}
    E_{\text{kötési}} = \left( M -  Z \cdot m_{\text{p}} - N \cdot m_{\text{n}} \right) \cdot c^2 < 0
\end{equation}
Mivel $M$ kisebb, mint az alkotóelemek tömegének összege, ezért a kötési energia negatív lesz, ami azt jelenti, hogy az atommag egy alacsonyabb energiájú állapot, mint a protonok és neutronok külön-külön.

\subsection{A cseppmodell és a félempirikus kötési formula}
Az atomok és az atomi elektronállapotok leírására a kvantummechanikai atommodell
nagyon sok sikert könyvelhet el. A sok elektront tartalmazó atomoknál problémát és
komplikációt okoz ugyan az elektronok egymás közötti kölcsönhatásának a figyelembevétele, ám a központi, (majdnem) pontszerűnek tekinthető atommag által létrehozott
centrális Coulomb-potenciál elsődleges hatása mellett az elektronok egymás közötti kölcsönhatása általában perturbációszámítással elég jó közelítéssel figyelembe vehető.

Egészen más a helyzet az atommagoknál. Itt nincsen egy központi mag, ami elsődleges rendezőként megszabná a nukleonok mozgását. Itt csak a nukleonok egymás közötti
kölcsönhatása hozza létre a kötést. Emiatt egy sok nukleont tartalmazó rendszer esetén
a kölcsönhatások száma rohamosan nő (még akkor is, ha csak két-test kölcsönhatásokat
feltételezünk), és ez roppant nehézzé teszi a nukleonok közötti kölcsönhatásokból kiinduló,
ab initio kvantummechanikai leírást. További problémát jelent, hogy a nukleonok
közötti kölcsönhatás sokkal bonyolultabb szerkezetű, és nem is ismert olyan pontossággal
és részletességgel, mint az atommag és az elektronok közötti Coulomb-kölcsönhatás.

Mindezek ahhoz vezetnek, hogy a magfizikusok az atommagok leírására - különböző
szempontok alapján - többféle egyszerűsítő modellt kénytelenek használni.

\textbf{Cseppmodell.} A cseppmodellben az atommagot elektromosan töltött folyadékcsepphez hasonlítjuk. táblázat alapján beláthatjuk, hogy melyek azok a tulajdonságok, amelyek a makrosz-
kopikus folyadékokat alkalmassá teszik arra, hogy az atommagokat modellezzük velük.
\begin{table}[H]
    \centering
    \begin{tabular}{|c|p{0.40\linewidth}|p{0.40\linewidth}|}
        \hline
        & \textbf{Folyadékcsepp} & \textbf{Atommag} \\
        \hline
        Alkotórészek & Molekulák & Nukleonok \\
        \hline
        Hatótávolság & Rövid (néhán molekulaméret $10^{-10}$ m) & Rövid (néhány nukleonméret $10^{-15}$ m)\\
        \hline
        Kölcsönhatás & Vonzó, és nem tesz különbséget a molekulák között (Van der Waals-erők), rövid távon taszítóvá válik (Pauli-elv) & Vonzó, nem tesz különbséget a nukleonok között (magerők), rövid távon taszítóvá válik (magerők tulajdonsága) \\
        \hline
    \end{tabular}
    \caption{A folyadékok és az atommagok közötti analógia.}
\end{table}

\textbf{Héjmodell.} (A héjmodell értelmezése inkább a Kvantummechanika domain-jába tartozik, ezért itt csak röviden említjük.) 
A folyadékcsepp-modell az atommagok sok tulajdonsgának az általános trendjét, szabályszerűségeit jól tudja értelmezni. Vannak azonban olyan atomma-
gok, amelyek nem követik ezeket az általános szabályokat, és kitűnnek a szomszédaikhoz
képest kiemelkedő stabilitásukkal. A megfigyelések szerint ezekben az atommagokban
a protonok vagy a neutronok száma (vagy mindkettő) a következő számok valamelyike:
2, 8, 20, 28, 50, 82, 126. Ezeket a számokat mágikus számoknak, az ilyen proton- vagy
neutronszámú atommagokat pedig mágikus atommagoknak nevezzük. Van néhány olyan
atommag is, amelyekben mind a proton- mind a neutronszám mágikus szám: $^4_2\mathrm{He}, ^{16}_8\mathrm{O},
^{40}_{20}\mathrm{Ca}, ^{56}_{28}\mathrm{Ni}, ^{208}_{82} \mathrm{Pb}$. Ezeket kétszer mágikus atommagoknak hívjuk.
A kiemelkedő stabilitás (nagy kötési energia) mellett több más tulajdonság is utal
ezeknek az atommagoknak a különleges viselkedésére. Például
\begin{itemize}
    \item azoknak az elemeknek van a legtöbb stabil izotópja, amelyekben a protonszám a mágikus számok valamelyike
    \item azoknak az atommagoknak a gerjesztéséhez szükséges kiemelkedően nagy energia, amelyekben akár a proton- akár a neutronszám (vagy mindkettő) a mágikus számok valamelyike
    \item a mágikus atommagok valamennyien gömbszimmetrikusak (kvadrupól-momentumuk nulla), és ez ugyancsak a kiemelkedő stabilitás és szimmetria egyik megnyilvánulási formája.
\end{itemize}

Az atomok tulajdonságainak és az elemek periódusos rendszerének értelmezésében bevált az elektronok héjmodellje. A Periódusos Rendszerben a nemesgázok felelnek meg
a ”mágikus” konfigurációnak, hiszen a nemesgázok (kémiai) stabilitása kiemelkedő! A
nemesgázoknak azonban a következő rendszámaik vannak: 2, 10, 18, 36, 54, 86. Az első
kivételével ezek teljesen más számok, mint az atommagoknál megfigyelt mágikus számok!
Nyilvánvaló, hogy a hasonlóság mellett lényeges különbségek is lesznek a két modell között. Hasonlítsuk tehát össze a két rendszert néhány tulajdonság szempontjából:
\begin{table}[H]
    \centering
    \begin{tabular}{|c|p{0.30\linewidth}|p{0.30\linewidth}|p{0.30\linewidth}|}
        \hline
        \textbf{Tulajdonság} & \textbf{Elektronburok} & \textbf{Atommag} \\
        \hline
        Alkotórészek & Elektronok & Nukleonok (n, p) \\
        \hline
        Alkotórészek tipikus tömege & 0.0005 u & 1 u \\
        \hline
        Tipikus méret (sugár) & 10$^{-10}$ m & 10$^{-14}$ m \\
        \hline
        Tipikus De\,Broglie hullámhossz & 10$^{-10}$ m & 10$^{-14}$ m \\
        \hline
        Pauli-elv & Igen & Igen \\
        \hline
        Tipikus kötési energia & $1-10$ eV & $4-8$ MeV \\
        \hline
        Potenciális energia & Coulomb potenciál & "Lapos" átlagpotenciál \\
        \hline
    \end{tabular}
    \caption{A folyadékok és az atommagok közötti analógia.}
\end{table}
Az elektronburok és az atommag tulajdonságainak összehasonlítása mutatja, hogy
az őket alkotó részecskék tömegében, energiájában és a rendszer geometriai méretében
fennálló több nagyságrendnyi különbség ellenére az elektronok rendszere és a nukleonok
rendszere hasonlít egymásra két dologban:
\begin{itemize}
    \item mindkettőben a részecskék tipikus hullámhossza a rendszer geometriai méreteinek nagyságrendjébe esik
    \item mind az elektronokra mind pedig a nukleonokra érvényes a Pauli-elv.
\end{itemize}
Az előbbi azt jelenti, hogy mindkettőre a részecskék kvantummechanikai modelljét kell
alkalmazzuk, az utóbbi viszont azt, hogy a protonok és a neutronok ugyanúgy kettesével
töltik be a rendelkezésre álló állapotokat, mint azt az elektronok teszik az atomburok
felépítésénél. Ebből azonnal következik, hogy az s állapotokba 2 proton és 2 neutron fér
(ellentétes spinnel), a p állapotokba 6 proton és 6 neutron, a d állapotokba 10 proton és
10 neutron, és így tovább.

\subsubsection*{Mit nevezhetünk egy héjnak?}
Akár az atomok, akár az atommagok kiemelkedő stabilitását szeretnénk értelmezni, meg
kell vizsgálnunk ennek a stabilitásnak a fizikai okát. Az ok az energiaszint-rendszerben rej-
lik. Egy kvantummechanikai rendszer (és ilyenek az atomok és az atommagok is) akkor
gerjeszthető könnyen, ha a legutolsó betöltött állapot ”fölött” (magasabb energián) közel
vannak üres állapotok. Ekkor ugyanis a betöltött állapotokban lévő valamelyik részecske
(elektron az atomban, ill. nukleon az atommagban) kis energiát felvéve, ”könnyen” ma-
gasabb energiájú állapotba kerülhet, és ezzel a rendszerünk máris gerjesztett állapotba
került. Ebből azonnal következik, hogy azok a rendszerek lesznek különösen stabilak,
ahol az utolsó betöltött állapot ”felett” nagy energiahézag van. Ekkor ugyanis csak nagy
energia befektetésével lehet a rendszert gerjeszteni. Ebből következik az energiahéj alábbi
definíciója:
\newline
\newline
\textbf{Definíció:} \textit{
    Energiahéjnak azon állapotok rendszerét nevezzük, amelyek energiában
közel vannak egymáshoz, de a többi állapottól a héjon belüli energiakülönbségeknél jóval
nagyobb energiahézag választja el őket.
}
\newline
\newline
Ennek alapján tehát a mágikus számok értelmezéséhez meg kell vizsgáljuk az atom-
magokban kialakuló nukleon-állapotok energiáit, és meg kell keressük az állapotrendszer-
ben található energia-ugrásokat.

\subsubsection{Cseppmodell}
A kölcsönhatás hasonlósága az oka az atommagok és a folyadékcseppek nagyfokú
hasonlóságának. A hasonlóságnak egyik érdekes következménye az atommag anyagának,
az ún. maganyagnak az atommag méretétől független, állandó sűrűsége. A maganyag
sűrűsége
\begin{equation}
    \rho = \frac{M}{V} = \frac{A \cdot m_{\text{nukleon}}}{\frac{4}{3} \pi R^3}
\end{equation}
Ha behelyettesítjük a maganyag sugárára vonatkozó $R = r_0 \cdot A^{1/3}$ összefüggést, akkor
\begin{equation}
    \rho = \frac{3 \cdot m_{\text{nukleon}}}{4 \pi r_0^3}
\end{equation}
hiszen az A tömegszám kiesett a képletből ($m_{\text{nukleon}}$ egyetlen nukleon tömegét jelenti). A
maganyag sűrűségére nagyságrendi becslést is adhatunk, hiszen a fenti képletben minden
mennyiség ismert.

A nukleon tömege kb. $m_{\text{nukleon}} \approx 1,67 \cdot 10^{-27}$ kg, és a maganyag sugara
közelítőleg $r_0 \approx 1,2 \cdot 10^{-15}$ m. Ezeket behelyettesítve a sűrűség
képletébe, megkapjuk a maganyag sűrűségének nagyságrendi értékét:
\begin{equation}
    \rho \approx 2,3 \cdot 10^{17} \ \mathrm{kg/m^3}
\end{equation}
Ez az érték hihetetlenül nagy! Összehasonlításképpen, a víz sűrűsége $10^3 \ \mathrm{kg/m^3}$, a vasé $7,9 \cdot 10^3 \ \mathrm{kg/m^3}$, a fehér törpe csillagok anyagának sűrűsége
kb. $10^9 \ \mathrm{kg/m^3}$. A maganyag sűrűsége tehát a hétköznapi anyagok sűrűségének több mint
$10^{14}$-szerese, és még a fehér törpe csillagok anyagának sűrűségét is több mint $10^8$-szor meghaladja! Ez az óriási sűrűség csakis annak köszönhető, hogy a maganyag
nukleonjai rendkívül kis térfogatban helyezkednek el.

\subsubsection{Az atommag energiája a cseppmodell szerint}
Az atommag energiáját a cseppmodell alapján határozzuk meg, annak figyelembevételével, hogy az atommag energiáját az egyes nukleonok energiáinak összege adja meg.
Vegyük tehát először szemügyre az egyes nukleonok kötési energiáját! Az energiák meghatározásánál a szabad részecskék energiájánál vesszük fel az energiaskála nullpontját,
tehát a kötött részecskék energiája negatív lesz.

\subsubsection*{A térfogati energia}
A csepp belsejében minden részecskének több szomszédja van, több részecskéhez tud
kötődni, ezért erősebben kötött (alacsonyabb energiájú) állapotban van, mint a felszínen
lévő részecskék. (Ilyen oka van a folyadékok felületi feszültségének is.) Ha minden nuk-
leon ”belső” nukleon lenne, akkor azonos energiával lenne kötve az atommagban, hiszen
a rövid hatótávolságú magerők miatt csak a közvetlen környezetével áll kölcsönhatásban
\begin{figure}[H]
    \centering
    \includegraphics[width=0.6\textwidth]{imgs/13-tetel/csepp_modell.png}
    \caption{A nukleonok kötöttsége bent és a felszínen.}
\end{figure}
eljöljük egy ilyen nukleon kötési energiáját $e_b$-vel! A nukleon energiája tehát $-e_b$.
Egy A nukleont tartalmazó atommag energiája tehát A-szor ekkora lenne, azaz $-e_b \cdot A$.
Az atommag energiájának ezt a részét térfogati energiának nevezzük, mivel a benne lévő
részecskék száma a mag térfogatával arányos. Azaz $E_V = -e_b \cdot A$. A negatív előjel
mutatja, hogy ez az energiatag erősíti a kötést.

\subsubsection*{A felületi energia}
Az atommag felszínén lévő nukleonok száma az atommag felszínével arányos, ezért a
felszín miatt bekövetkező energianövekedés nagysága: $E_F = b_F \cdot 4\pi R^2$. Ez az energiatag
pozitív előjelű, hiszen a térfogati energiában a felszíni nukleonokat is ”belső” nukleonként
vettük számba, ezért a kötésüket túlbecsültük.

\subsubsection*{A Coulomb-energia}
A protonok elektromos töltésüknél fogva taszítják egymást, és így az atommag energiájának meghatározásakor a nukleáris energiatagok mellett egy Coulomb-energia is fellép.
A Coulomb-energia $\frac{Q^2}{R}$-el arányos, ahol $Q$ az atommag töltése, és $R$ pedig a sugara. Felhasználva, hogy $Q = Ze$, kapjuk, hogy a Coulomb energia $E_C = b_C \frac{(Ze)^2}{R}$.
A Coulombenergia gyengíti a kötést, tehát pozitív előjelű.

\subsubsection*{Az asszimetriai energia}

Még a cseppmodellben is figyelembe kell venni, hogy az atommag nem klasszikus objektum, hiszen a benne lévő részecskék de Broglie-hullámhossza az objektum méretével
összemérhető, valamint, hogy a protonokra és a neutronokra érvényes a Pauli-elv. A
Pauli-elv miatt egy energiaszinten legfeljebb két részecske lehet, ezért a protonok és a
neutronok számának eltérése többletenergiát eredményez. Egy $N$ neutront és $Z$ protont tartalmazó atommagban a protonoknál magasabb szintre szorult neutronok száma
($N - Z$). De ugyancsak ($N - Z$)-vel arányos az egy neutronra jutó átlagos többletenergia is,
mivel ($N - Z$)-vel arányosan magas szintig töltődnek be a többlet energiaszintek. A Pauli-elvből adódó többletenergiát, - amit aszimmetria-energiának nevezünk - tehát a következőképpen írhatjuk fel:
\begin{equation}
    E_A = e_A \frac{(N-Z)^2}{A} = e_A \frac{(A - 2Z)^2}{A}
\end{equation}
A neutron-proton aszimmetria gyengíti a kötést, ezért az aszimmetria energia pozitív
előjelű.
\begin{figure}[H]
    \centering
    \includegraphics[width=0.6\textwidth]{imgs/13-tetel/csepp_modell2.png}
    \caption{Asszimetrikus energia ábrázolása.}
\end{figure}

\subsubsection*{A párenergia}
Az atommagok kötési energiájának tapasztalati (empirikus) vizsgálatából leszűrhető,
hogy azok az atommagok, amelyekben valamelyik nukleonfajta (proton vagy neutron)
páros számú, erősebben kötöttek, mint amelyekben páratlan számú proton és/vagy neutron van.
Párosság szempontjából három csoportba sorolhatjuk az atommagokat:
\begin{itemize}
    \item Páros-páros (ps-ps) atommagok: mindkét nukleonfajta páros számú
    \item Páros-páratlan (ps-ptl) vagy páratlan-páros (ptl-ps) atommagok: az egyik nukleonfajta páros, a másik páratlan számú
    \item Páratlan-páratlan (ptl-ptl) atommagok: mindkét nukleonfajta páratlan számú
\end{itemize}
A párenergia felléptének az az oka, hogy az atommagokban az azonos fajta nukleonok párokba kapcsolódnak. Ennek
a párkölcsönhatásnak az az eredménye, hogy a több száz stabil atommag közül mindössze
négy van olyan, amelyikben mindkét komponens páratlan számú: $^1_1H_1$, $^6_3Li_3$, $^{10}_5B_5$, $^{14}_7N_7$. 
A párenergia felléptet a cseppmodell nem tudja megmagyarázni, ennek kvantummechanikai
oka van. Az atommagok energiájának kiszámításakor egy tapasztalati képlettel mégis
figyelembe szoktuk venni:
\begin{equation}
    E_P = e_p \cdot \delta \cdot A^{-\frac{1}{2}} \quad \text{Ahol,} \quad \delta = \begin{cases}
        1, & \text{ps-ps magok esetén} \\
        0, & \text{ps-ptl vagy ptl-ps magok esetén} \\
        -1, & \text{ptl-ptl magok esetén}
    \end{cases}
\end{equation}
A tapasztalati $e_p$ konstans negatív, mert páros-páros atommagok esetén a kötést erősíteni, páratlan-páratlanok esetén pedig gyengíteni kell.

\subsubsection{A félempirikus kötési energia formula}
A fentieket összefoglalva megállapíthatjuk, hogy az atommagok teljes energiáját tehát öt
tag összegeként kaphatjuk meg, amelyek az atommagok összetételének (a $Z$ rendszámnak
és az $A$ tömegszámnak), valamint az atommag sugarának ($R$) a függvényei:
\begin{equation}
    E(Z, A, R) = -e_V \cdot A + b_f \cdot 4 \pi R^2 + b_C \frac{(Ze)^2}{R} + e_A \frac{(A - 2Z)^2}{A} + e_p \cdot \delta \cdot A^{-\frac{1}{2}}
\end{equation}
Megint felhasználva, hogy az atommag sugara $R = r_0 \cdot A^{1/3}$, az atommag teljes energiája a következőképpen írható fel:
\begin{equation}
    \boxed{
        E(Z, A) = -e_V \cdot A + e_F \cdot A^{2/3} +  e_C \frac{Z^2}{A^{1/3}} + e_A \frac{(A - 2Z)^2}{A} + e_p \cdot \delta \cdot A^{-\frac{1}{2}}
    }
\end{equation}
Ezt nevezik a Weizsäcker-féle \textbf{félempirikus kötési energia formulának}. A képletben szereplő egyes tagokat tapasztalati (empirikus) úton meghatározott konstansok.
Ezeknek az értékeit az irodalomban megtalálhatjuk (néhol kissé eltérő értékekkel):
\begin{itemize}
    \item Térfogati energia együtthatója: $e_V \approx 15.69$ MeV
    \item Felületi energia együtthatója: $e_F \approx 17.81$ MeV
    \item Coulomb energia együtthatója: $e_C \approx 0.69$ MeV
    \item Asszimetriai energia együtthatója: $e_A \approx 23.75$ MeV
    \item Párenergia együtthatója: $e_p \approx -9.31$ MeV
\end{itemize}
Az atommag kötési energiája tehát a fentiek alapján
\begin{equation}
    B(Z, A) = - E(Z, A)
\end{equation}
A fenti konstansokkal a fenti képlet az ismert több mint 2000 különböző (stabil és radioaktív)
atommag energiáját kb. 1-2\%-os pontossággal megadja. (A képlet által adott értéktől csak a mágikus atommagoknál vannak kisebb eltérések). Ebből is látszik, hogy a
cseppmodell az atommagok kötési energiájának lényeges vonásait jól ragadja meg. 

\subsection{Maghasadás, magfúzió, radioaktivitás}

\subsubsection{Maghasadás}
A maghasadás felfedezése (1938) O. HAHN (1879-1968, Kémiai Nobel-díj 1944), és F. STRASSMANN (1902-1980) által, alapjaiban megváltoztatta a XX. század arculatát,
hiszen ez tette lehetővé az atomenergia békés célú felhasználását és, sajnos, az atomfegyverek létrehozását is. Indokolt tehát, hogy kissé részletesebben foglalkozzunk ezzel a speciális atommagreakcióval.

Maghasadáskor egy nehéz atommag (tórium, urán, transzuránok) két közepes méretű atommagra hasad szét. Ez bekövetkezhet magától (ilyenkor spontán maghasadásról
beszélünk), vagy valamilyen külső behatásra (pl. neutron elnyelését követően). A folyamatot úgy képzelhetjük el, hogy a kezdetben kissé deformált nehéz atommag megnyúlik,
majd befűződik, és végül két részre szakad. A keletkezett részeket hasadványoknak nevezzük. Ezt a folyamatot nagy energia felszabadulása és néhány neutron kilépése szükségszerűen kíséri.

\subsubsection*{Maghasadás energiaviszonyai}
A mikrofizikában, és így a magfizikában is, a makroszkopikus világunkban általában megszokott energiáknál sok nagyságrenddel kisebb energiák fordulnak elő. Ezért ezekben az
energiatartományokban nem az SI-egységrendszerben szokásosan használt $J$ (Joule) egységet, hanem egy annál jóval kisebb egységet,
az elektronvoltot ($eV$), illetve annak többszöröseit szokás használni. $1 eV$ energiát kap egy elemi töltésű ($1.6 \cdot 10^{-19} C$) részecske, ha $1 V$ potenciálkülönbségen halad át. Ezért
\begin{equation}
    1 eV = 1.6 \cdot 10^{-19} J
\end{equation}
A nagyobb energiák leírására ennek többszöröseit használjuk: $1 keV = 10^3 eV$,
$1 MeV = 10^6 eV$, $1 GeV = 10^9 eV$ stb.

Az elemi folyamatokban felszabaduló energiák leírása ezekben az egységekben célszerűbb, míg a makroszkopikus, műszaki alkalmazásoknál, természetesen a $J$ és annak
többszörösei megfelelőbbek.
\begin{figure}[H]
    \centering
    \includegraphics[width=0.6\textwidth]{imgs/13-tetel/maghasadas.png}
    \caption{Maghasadás szemléltetése.}
\end{figure}

A maghasadás energiaviszonyainak megértéséhez elegendő, ha az atommagok energia-
felületére gondolunk. Amikor az anyag nagyot ”ugrik” az energiavölgyön, azaz a nehéz atommagokból maghasadással közepes tömegű atommagok keletkeznek,
akkor minden nukleon mélyebben fekvő, erősebben kötött állapotba jut! Bár a
völgynek ez a része ”lankás”, azaz egy-egy nukleonnak csak viszonylag kis (kb. $1 MeV$)
energianyeresége származik az ugrásból, de a több mint 200 nukleon együttesen mégis
jelentős energianyereséget jelent. A számításhoz az egyszerűség kedvéért a párenergiát
hanyagoljuk el.
Tegyük fel, hogy egy $Z$ rendszámú $A$ tömegszámú atommag két egyenlő részre hasad.
Ha $E(Z, A)$ jelenti a kiindulási atommag energiáját, akkor a maghasadás létrejöttének
energia-feltétele az, hogy: $E(Z, A) - 2E\left(\frac{Z}{2}, \frac{A}{2}\right) > 0$. Felírva ezt a két kifejezést a félempírikus formula alapján, azonnal látható, hogy csak a felületi energia és a Coulomb-energia ad járulékot
a különbségbe, a többi tag kiesik. Az eredmény:
\begin{equation}
    e_C \left(\frac{Z^2}{A^{1/3}}\right)\left(1-2^{-\frac{2}{3}}\right) - e_F A^{\frac{2}{3}} \left(1-2^{-\frac{1}{3}}\right) > 0
\end{equation}
Ezt átrendezve:
\begin{equation}
    \frac{Z^2}{A} > \frac{e_F}{e_C} \cdot \frac{1 - 2^{-\frac{1}{3}}}{1 - 2^{-\frac{2}{3}}}
\end{equation}
Vagy behelyettesítve az együtthatók értékeit:
\begin{equation}
    \frac{Z^2}{A} > 18
\end{equation}
Ahol $e_C$ és $e_F$ a Coulomb-energia, illetve a felületi energia együtthatói és a cseppmodellből származnak.

Mivel a könnyű atommagoknál $A = 2 \cdot Z$, ezért nagyjából a $Z > 36$ rendszámtól (ez a kripton) kezdve a maghasadás energia-felszabadulással jár.
A különbséget kiszámíthatjuk a $^{238}U$ esetére. Itt $Z = 92$, $A = 238$. Tehát
\begin{equation}
    0.11 \cdot 10^{-12} \frac{92^{2}}{238^{\frac{1}{3}}} \left(1 - 2^{-\frac{2}{3}}\right) - 2.85 \cdot 10^{-12} \cdot 238^{\frac{2}{3}} \left(2^{-\frac{1}{3} - 1}\right) = 27.1 \cdot 10^{-12} J = 169.4 \ MeV
\end{equation}

A kapott érték a kísérletileg mért $\approx 200 MeV$ értéknél mintegy 15\%-al kisebb. Ezen
azonban nem lehet csodálkozni, hiszen több közelítést is alkalmaztunk. Az egyik köze-
lítési például az volt, hogy azt feltételeztük, hogy az uránmag két egyenlő részre hasad.
Tudjuk, hogy ez általában nem így történik, ezért ettől a becsléstől nem is várhatunk 10-
20\%-nál pontosabb eredményt. Számításunk azt azért megmutatja, hogy maghasadáskor
mikrofizikai léptékkel mérve hatalmas, $150-200 MeV$ nagyságrendű energia szabadul fel!
Ez nagyjából százmilliószor akkora, mint a $C+O_2 \rightarrow CO_2$ égési (kémiai) reakcióban felszabaduló energia.

\subsubsection*{Hasadási gát}
Ha a maghasadás energianyereséges, miért nem hasadt már szét réges-régen minden uránatommag? A maghasadás létrejöttét megakadályozza egy potenciálgát, az ún. hasadási
gát. Ennek kialakulását a fordított folyamat segítségével érthetjük meg a legkönnyeb-
ben. A fordított folyamat során a széthasadt mag két része (a két hasadvány) egymás felé
közeledik, majd egyesül. Amíg a két mag távol van egymástól, addig lényegében csak
mozgási energiájuk van. Ahogyan közelednek egymáshoz, a közöttük lévő Coulomb-
taszítás mind erősebbé válik, ezért mozgási energiájuk csökken, a potenciális energia nő.
Ez addig folytatódik, amíg ”össze nem érnek”, és a rövid hatótávolságú, erősen vonzó
magerők működésbe nem lépnek. Innen kezdve ”összerántja” őket a magerőkbol adódó
nukleáris ”felületi feszültség”, azaz a potenciális energia ismét lecsökken. Így alakul ki az ábrán látható potenciál, ami tehát a két töredékmag (hasadványok) kölcsönhatását
írja le a közöttük lévő távolság függvényében. A potenciálgát ”teteje” ott van, ahol a
maghasadás során a két hasadványmag még éppen érintkezik. Ez a potenciálgát a hasa-
dási gát. Természetesen, a hasadási gát ”belülről” is gát, és ez az, ami megakadályozza
azt, hogy egy atommag spontán gyorsan széthasadjon, hiába járna a folyamat végül nettó
energiafelszabadulással.

Vizsgáljuk meg, hogy egy atommaggal milyen feltétel mellett alakulhat ki egyáltalán hasadási gát. A maghasadás nyilván az alapállapotú atommag kis deformációjával
kezdődik. Jelöljük a deformációs paramétert $\epsilon$-al, és szorítkozzunk $\epsilon \ll 1$ kis deformációkra. Feltételezzük, hogy az alakváltozás kis kvadrupólus deformációval indul, azaz a
kezdetben gömb alakú atommag enyhén megnyúlt, forgási ellipszoid alakot vesz fel. A mag három féltengelyének hossza a következőképpen fejezhető
\begin{figure}[H]
    \centering
    \includegraphics[width=0.6\textwidth]{imgs/13-tetel/hasadasi_gat.png}
    \caption{A maghasadás potenciálgátja.}
\end{figure}
ki az $\epsilon$ deformációs paraméterrel, és a mag eredeti $R_0$ sugarával:
\begin{equation}
    \begin{aligned}
        R_z &= R_0 \left(1 + \frac{2}{3}\epsilon \right) \\
        R_x = R_y &= R_0 \left(1 - \frac{1}{3}\epsilon \right)
    \end{aligned}
\end{equation}
Azaz a mag a $z$-tengely mentén megnyúlik, arra merőlegesen viszont mindkét irányban
lecsökken a mérete. Ilyen választás mellett a mag térfogata $\epsilon$-ban első rendig változatlan
marad (az atommag anyaga első közelítésben nem nyomódik össze).

A deformációhoz szükséges energiát a már megismert Weizsäcker-féle félempirikus kötési energia formula alapján becsülhetjük meg. Mivel az atommag anyagát
összenyomhatatlannak tételezzük fel, azaz a deformáció során a mag térfogata állandó
marad, ezért a kötési energia formulában a térfogati energia tag nem változik. A deformáció nyilván nem érinti a mag összetételét, ezért mind az aszimmetria-tag, mind pedig
a pár-energia tag állandó marad. A deformáció során azonban változik - megnő - a felületi energia tag, hiszen a deformált mag felülete nagyobb, és változik a Coulomb-energia
tag is - csökken -, hiszen a deformációval a protonok átlagos távolsága nagyobb lesz. Kis
deformációknál feltehetjük, hogy ezek a változások a deformációs paraméterrel lineáris
kapcsolat szerint változnak (a sorfejtésben megállunk a lineáris tagnál). Más szóval az
energiaváltozás a deformáció hatására:
\begin{equation}
    \Delta E = \left(e_F A^{\frac{2}{3}}\right)a_F \cdot \epsilon + \left(e_C \frac{Z^2}{A^{\frac{1}{3}}}\right)a_C \cdot \epsilon
\end{equation}
Itt $a_F$ , illetve $a_C$ az arányossági tényezők a sorfejtésekben. Pontos számítások szerint
$a_F = 0.025$ és $a_C = 0.012$. Érdekes, hogy az energiaváltozás két, egymással ellentétes
irányba ható hatás eredményeként adódik. Ez azt jelenti, hogy az egyes tagok abszolút
értékétől függően az energiaváltozás lehet pozitív, de lehet negatív is. Ha már kis deformációkra is $\Delta E < 0$, akkor a maghasadás azonnal el tud indulni, hiszen nincs hasadási gát.
ilyenkor azonnali \textit{spontán maghasadás} következik be. Ennek a feltétele tehát
\begin{equation}
    e_F A^{\frac{2}{3}} \cdot 0.025 < e_C \frac{Z^2}{A^{\frac{1}{3}}} \cdot 0.012
\end{equation}
Ebből átrendezéssel azt kapjuk, hogy:
\begin{equation}
    \frac{Z^2}{A} \leq \frac{0.025 \cdot e_F}{0.012 \cdot e_C} \approx 54
\end{equation}
A nagy atommagok tartományában $\frac{Z}{A} \approx 0.39$ (pl. az $^{235}U$-ra $92/235=0.391$), és ezt
behelyettesítve azt kapjuk, hogy $Z > 136$ esetén már nem lehet semmilyen hasadási
gát, az ilyen atommagok már létre sem jöhetnek. A Periódusos Rendszernek a spontán
maghasadás szabja meg a végét. Valójában a Periódusos Rendszernek már valahol $Z=110$
és $Z=120$ között vége szakad. Ekkor ugyan még létezik hasadási gát, de már olyan
alacsony, hogy az atommag alagúteffektussal a másodperc töredéke alatt el tud hasadni.
Az uránmag esetén a hasadási gát még elég magas, kb. $7-8 MeV$. Ennyit ”be kell
fektessünk”, hogy átemeljük az atommagot a magas potenciálgáton, de ezt a befektetett
energiát bőségesen visszakapjuk: a hasadás során kb. 25-30-szer ekkora energia válik
szabaddá.

Alagúteffektussal az uránmag ez alatt a magas potenciálgát alatt csak roppant kis
valószínűséggel tud áthaladni - az urán spontán maghasadása kétmilliószor ritkább, mint
alfa-bomlása, pedig ez utóbbi sem gyakori.

Vannak persze a 92-es rendszámú uránnál is nehezebb, mesterségesen előállított atommagok (transzuránok ). Minél nagyobb a rendszám és a tömegszám, annál alacsonyabb
lesz a hasadási gát. A 103 rendszámú kurcsatoviumnál a bomlás már elsősorban maghasadással történik, a felezési idő mindössze egy másodperc.

\subsubsection*{Neutronokkal létrehozott maghasadás}
Az előző szakaszban láttuk, hogy az urán környékén a hasadási gát elég magas, 7-8 MeV körül van. Felvetődik a kérdés, hogy hogyan tud akkor egy termikus neutron az
235U-ban maghasadást létrehozni, hiszen a termikus energia sok nagyságrenddel kisebb
(szobahőmérsékleten $\approx 0,025 eV$). A megoldás a neutron kötési energiájában rejlik . Látható, hogy amikor egy neutron befogódik az atommagba, a nukleáris
vonzás miatt a neutron kötési energiája ”felszabadul”, illetve átadódik a magnak. Ez sok
\begin{figure}[H]
    \centering
    \includegraphics[width=0.6\textwidth]{imgs/13-tetel/neutron_befogado.png}
    \caption{Neutron befogódás és maghasadás.}
\end{figure}

nagyságrenddel nagyobb lehet, mint a neutron (hőmozgásából eredő) mozgási energiája.
Ez a kötési energia már elegendő lehet arra, hogy az atommagot átsegítse a hasadási
gáton.

Felvetődik egy másik kérdés is: nem várható, hogy a hasadási gát lényegesen különböző magasságú legyen a 235U és a 238U esetén. Mi az oka mégis annak, hogy lassú,
termikus neutronok hatására a $^{235}U$ el tud hasadni, a $^{238}U$ pedig nem?
A válasz a pár-energiában rejlik. Az $^{235}_{92}U+n \leftarrow ^{236}_{92}U^*$ atommag-reakcióban páratlan-páros atommagból páros-páros atommag keletkezik, ami erősebben kötött. Ezért ebben
a reakcióban a pár-energia hozzáadódik a neutron pár-energia nélkül számított kötési
energiájához. A $^{238}_{92}U+n \rightarrow ^{239}_{92}U^*$ reakció esetén éppen fordított a helyzet. Itt erősebben
kötött, páros-páros atommagból keletkezik egy gyengébben kötött, páratlan-páros atommag. Itt tehát a pár-energia levonódik a neutron pár-energia nélkül számított kötési
energiájából. Az $^{235}U$-nál emiatt a pár-energia kétszeresével több energia jut az atommagnak a neutron befogódása után, mint az $^{238}U$-nál. Az energiaviszonyok olyanok, hogy
az egyik esetben elegendő az energia a hasadási gáton való átjuttatáshoz, a másik esetben viszont éppen nem. Az $^{238}U$-nál a maghasadás küszöbreakció, csak kb. 1 MeV-nél
nagyobb energiájú neutronok képesek maghasadást létrehozni.

Hasonló a helyzet a többi, urán környékén lévő nehéz elemnél is (pl. transzuránoknál):
a páratlan-páros atommagok (pl. $^{239}_{94}Pu$) lassú neutronok hatására is hasadóképesek, a
páros-páros atommagok viszont (pl. $^{240}_{94}Pu$) kisebb-nagyobb energiaküszöbbel rendelkeznek.

\subsubsection*{A maghasadásban felszabaduló energia térbeli és időbeli megoszlása}
A gyakorlati felhasználás szempontjából nagyon fontos, hogy a hasadáskor felszabaduló
energia milyen komponensekből áll, és térben hol, ill. időben mikor jelenik meg. A
táblázat összefoglalja a különböző komponensek arányát és tulajdonságait arra az esetre,
amikor a maghasadás egy atomreaktor üzemanyagpálcájában következik be, amelyet
hűtőközeg (víz) vesz körül.
\begin{table}[H]
    \centering
    \begin{tabular}{|p{0.30\linewidth}|c|c|c|c|}
        \hline
        Leírás & Értéke & Aránya & Térben & Időben \\
        \hline
        Hasadványok kinetikus energiája & 168 MeV & 82\% & Reaktor üzemanyag & Azonnal \\
        \hline
        Hasadási neutronok által bevitt energia & 5 MeV & 2.4\% & hosszú hatótávolság & Azonnal \\
        \hline
        Prompt $\gamma$-fotonok által elvitt energia & 7 MeV & 3.4\% & hosszú hatótávolság & Azonnal \\
        \hline
        Hasadási termékek $\beta$-részecskéi által bevitt energia & 8 MeV & 3.9\% & közepes hatótávolság & később \\
        \hline
        Hasadási termékek $\gamma$-részecskéi által bevitt energia & 7 MeV & 3.4\% & hosszú hatótávolság & később \\
        \hline
        Hasadási termékek $\beta$-bomlásakor keletkezett antineutrínók által bevitt energia & 10 MeV & 4.9\% & nagyon messze & később \\
        \hline
    \end{tabular}
\end{table}

Az antineutrínók által elvitt energiát nyugodtan elfelejthetjük, hiszen nincs arra mód,
hogy azokat bárhogyan is hasznosítani tudjuk. Nem ez a helyzet a hasadási termékek
béta-bomlásakor kibocsátott elektronok, illetve gamma-fotonok energiájával. Ezek ugyan
közepes, ill. hosszú hatótávolságúak, ám az energiájuk nagy részét a hűtőközegben,
és/vagy a reaktortartályban fogják leadni, azaz azt melegítik. Fontos, hogy ezek már korábban megtörtént maghasadások kísérőjelenségei, azaz időben jóval a maghasadás után
jelentkeznek. Összességében ezek a maghasadásokkor felszabaduló energiának $3.9+3.4 = 7,3\%$-át jelentik. Egy 3000 MW hőteljesítményű atomreaktorban ez a 7,3\% igen nagy,
$219 MW$ hőteljesítményt jelent! A korábbi maghasadások következtében létrejött radioaktivitás ennyivel fűti a reaktort még azután is, hogy a láncreakció (a maghasadás)
leállt! Emiatt roppant fontos, hogy egy atomreaktorban még a láncreakció leállítása után
se maradjon ki az üzemanyag hűtése, hiszen ezt a radioaktivitást nem lehet ”leállítani”.
Természetesen, a radioaktív bomlás miatt ez a fűtési teljesítmény csökken. Eleinte rohamosan (hiszen a rövid felezési idejű, legnagyobb aktivitású izotópok hamar elbomlanak),
később egyre lassabban. A kiégett fűtőelemeket ezért még éveig folyamatosan hűteni kell.
\newline
\newline
További infó: \href{https://mek.oszk.hu/16000/16052/16052.pdf#section.8.2}{Kisérleti magfizika - Hasadási neutron, atomreaktorok}

\subsubsection{Magfúzió}
A maghasadáson alapuló energiatermelés mellett a nukleáris energia felszabadításának
másik útja a magfúzió. A fúzió során a periódusos rendszer elején levő, kis rendszámú
elemek egyesülnek, ami - az atommagok energia-felülete alapján - 
energiafelszabadulással jár. Nagyobb tömegszámú, erősebben kötött atommag jön létre,
és ezáltal kötési energia szabadul fel. A fúziót az atommagok elektrosztatikus (Coulomb-)
taszítása gátolja, ezért csak valamekkora aktiválási energia befektetésével hozható létre.
A Coulomb-gát hidrogénizotópok esetén a lehető legkisebb, ezért ezeknek a fúziója
aktiválható a legkönnyebben. A fúzió létrejöttét segíti az alagúteffektus, és ezért kis valószínűséggel a Coulomb-gát alatti energiáknál is végbemehet.
A maghasadásnál a láncreakciót a neutronokkal tudtuk megvalósítani, mert a neutronok befogásakor a hasadás aktiválásához elegendő energia szabadult fel. Sajnos, a
fúziónál ilyen lehetőség nincs. A fúzióhoz ugyanis két atommag szükséges, amelyek elegendően nagy sebességgel mozognak egymás felé. Sebességükből eredő mozgási energiájuknak kell fedezni a Coulomb-gát megmásásához szükséges energiát. Makroszkopikus
méretű, láncreakcióra alkalmas módon történő mozgás csak hőmozgással tűnik megvalósíthatónak. Ezért a fúziós energiatermelés egyedüli járható útja jelenleg a termikus
aktiválás.

Megemlítjük, hogy laboratóriumban, magfizikai kísérletek céljaira kevés számú részecskét elektrosztatikusan is fel lehet gyorsítani elegendően nagy sebességre, és így a
fúziós reakció kis méretekben létrehozható. Ezen az elven működnek az ún. neutrongenerátorok. Nevük onnan ered, hogy segítségükkel a kísérletező fizikusok a fúziós folyamatok
során kilépő, nagy energiájú neutronokat állítják elő. A neutrongenerátorokban alkalmazott elektrosztatikus gyorsítás azonban makroszkopikus méretű, energia-felszabadító
láncreakció létrehozására alkalmatlan.

Ennek az az oka, hogy amikor a nagy sebességre felgyorsított részecskék behatolnak
az anyagba, a fúzió létrejöttének a valószínűsége sokkal kisebb, mint más egyéb folyamatoké (pl. rugalmas szóródás, ionizálás stb.). Az ”egyéb” folyamatokban viszont a bejövő,
nagy energiára felgyorsított részecskék gyorsan elveszítik az energiájukat, és alkalmatlanná válnak a fúzió létrehozására.
Ezeknek a folyamatoknak a következtében a bejövő részecskék mozgási energiája addig csökken, amíg termikus egyensúlyba nem kerülnek a
közeggel. Ezért a termikus mozgás energiáját kell olyan magas értéken tartani, hogy az már elég legyen a fúzió aktiválásához.

A termikus aktiváláshoz mintegy $10^7 - 10^8 K$ körüli hőmérséklet szükséges. Ezen a
hőmérsékleten az anyag már teljesen ionizált állapotba kerül. Ezt az állapotot nevezzük
plazmának. A plazma kifelé elektromosan teljesen semleges, mivel egyenlő számban tartalmaz elektronokat és pozitív ionokat. A plazma sok szempontból a közönséges gázokhoz
hasonlítható, lényeges különbség azonban az, hogy a plazma jól vezeti az elektromosságot, míg a semleges gázok nem. Hasonlítanak viszont abban, hogy a részecskék sebessége
nem egy meghatározott érték, hanem jellegzetes elosztást - ún. Maxwell-elosztást - követ.

\subsubsection*{Fúziós folyamatok}
Emlékezve a $^4_2He$ atommag kiemelkedően nagy kötési energiájára, első pillanatban a következő fúziós folyamatra gondolhatunk:
\begin{equation}
    ^2_1H + ^2_1H \rightarrow ^4_2He
\end{equation}
Közelebbi vizsgálat azonban megmutatja, hogy két összeütköző deutérium atommag nem
egyesülhet egyetlen újabb atommagga ($^4_2He$), mert ilyen ütközésnél nem teljesülhet egyszerre az energia-megmaradás és a lendület-megmaradás.

A fúziós energiatermelés szempontjából legértékesebb reakciók a következő:
\begin{equation}
    \begin{aligned}
        &^2_1H + ^2_1H &\rightarrow ^3_2He + n + 3.25 MeV \\
        &^2_1H + ^3_1He &\rightarrow ^4_2He + n + 17.6 MeV \\
    \end{aligned}
\end{equation}

\subsubsection*{Fúzió a csillagokban}
A termikusan aktivált fúzió létrehozásához magas hőmérséklet ($10^7 - 10^8 K$) szükséges.
A csillagok belsejében az összegyűlt, hatalmas mennyiségű anyag gravitációs nyomása
ellensúlyozni képes a sűrű, magas hőmérsékletű plazma nyomását. A csillag egyensúlyi
módon, folyamatosan termeli az energiát a belsejében, hidrogénatomok fúziójával. A hid-
rogénre alapozott fúzióval energiát termelő csillagok a csillagfejlődés kezdeti szakaszában
vannak. Ezeket H-égő csillagoknak nevezzük.
A csillagok (pl. a Nap) színképéből azonban kiderül, hogy bennük alig van deutérium
és trícium. A H-égő csillagok tehát NEM a az előbb látott reakciókkal működnek. Bennük
érdekes módon a könnyű hidrogénből (proton) keletkezik a hélium.

\textbf{proton-proton (pp) ciklus.}
Első pillanatban lehetetlennek látszik, hogy könnyű hidrogénnel fúziót lehessen megvalósítani, hiszen a $^1_1H+^1_1H\rightarrow ^2_2He$ reakció nem mehet végbe, mivel - magfizikai okoknál
fogva - két proton nem tud kötött állapotot létrehozni. A gyenge kölcsönhatás azonban
itt is segít! Roppant kis valószínűséggel az alatt a rövid idő alatt, amíg a két proton
egymás közelében, a Coulomb-taszítás által létrehozott magasabb energiájú állapotban
van, végbemehet az egyik proton neutronná alakulása, és a proton-neutron párból egy
deutérium atommag létrejötte:
\begin{equation}
    ^1_1H + ^1_1H \rightarrow ^2_1H + e^+ + \nu_e \qquad (Q = 0.42 MeV)
\end{equation}
A keletkezett pozitron szinte azonnal annihilál egy elektronnal:
\begin{equation}
    e^+ + e^- \rightarrow 2 \gamma \qquad (Q = 1.022 MeV)
\end{equation}
A keletkezett deuteron már gyorsan találkozik egy protonnal, és fuzionál:
\begin{equation}
    ^2_1H + ^1_1H \rightarrow ^3_2He + \gamma \qquad (Q = 5.49 MeV)
\end{equation}
Végül két $^3_2He$ atommag fejezi be a $^4_2He$ felépítését:
\begin{equation}
    ^3_2He + ^3_2He \rightarrow ^4_2He + 2 ^1_1H \qquad (Q = 12.86 MeV)
\end{equation}
Összegezve ezeket a folyamatokat végeredményben négy protonból lesz egy hélium atommag (közben két pozitron és két neutrínó is keletkezik):
\begin{equation}
    4(^1_1H) \rightarrow ^4_2He + 2 e^+ + 2 \nu_e \qquad (Q = 24.69 MeV)
\end{equation}
A teljes energiamérlegbe azonban a pozitronok annihilációjakor felszabaduló energiát is
bele kell számítani, végeredményben tehát egyetlen $^4_2He$ atommag létrejöttekor összesen
$24.69+2.044=26.73 MeV$ energia szabadul fel.

\textit{Megjegyzés:} a pp-ciklusnak vannak egyéb ”oldalágai” is, ezekkel azonban itt most nem foglalkozunk.

Mint említettük, a fúziós lánc ”indulása”, az első reakció roppant kis valószínűségű. Annyira kicsi, hogy egy kiszemelt proton másodpercenként többmilliárdszor ütközve
másik protonnal évmilliárdokig bolyonghat a Nap roppant sűrű és sokmillió fokos középpontjában, amíg egyszer véletlenül létrejön ez a reakció. Gondoljunk bele, hogy ez
a roppant vékony szálon csordogáló magreakció a fék, amely lehetővé teszi, hogy a Nap hidrogénjében rejlő hatalmas fúziós energia évmilliárdok alatt szabaduljon fel csaknem
egyenletes ütemben, elegendő időt hagyva a Naprendszer egyik — megfelelő távolságban lévő — bolygóján Élet kialakulására és az evolúcióra. Ennek köszönhetjük a létünket! Ez
az a fék, amely megakadályozza, hogy a Nap (és a többi hasonló csillag) nagyon rövid idő alatt iszonyú égi tűzijáték formájában, gigantikus hidrogénbombaként szétrobbanjon.
\begin{figure}[H]
    \centering
    \includegraphics[width=0.6\textwidth]{imgs/13-tetel/pp_cycle.png}
    \caption{A proton-proton ciklus vázlata.}
\end{figure}

\textbf{CNO-ciklus.} 
A csillagokban lejátszódó folyamatokat vizsgálva a közvetlen fúziós reakciókon túl érdekes
katalizálási lehetőségre mutatott rá H. BETHE (1906-2005, Nobel-díj 1967) német fizikus.
A csillagban levő, nagyobb rendszámú elemek — például a $^{12}_6 C$ — a hidrogénfúzióban katalizátor szerepet kaphatnak.

A körfolyamat lépései:
\begin{equation}
    \begin{aligned}
        ^{12}_6 C + ^1_1 H \quad &\rightarrow \quad ^{13}_7 N \\
        ^{13}_7 N \quad &\rightarrow \quad ^{13}_6 C + e^+ + \nu_e \\
        ^{13}_6 C + ^1_1 H \quad &\rightarrow \quad ^{14}_7 N \\
        ^{14}_7 N + ^1_1 H \quad &\rightarrow \quad ^{15}_8 O \\
        ^{15}_8 O \quad &\rightarrow \quad ^{15}_7 N + e^+ + \nu_e \\
        ^{15}_7 N + ^1_1 H \quad &\rightarrow \quad ^{12}_6 C + ^4_2 He \\
    \end{aligned}
\end{equation}
\begin{figure}[H]
    \centering
    \includegraphics[width=0.4\textwidth]{imgs/13-tetel/cno_cycle.png}
    \caption{A CNO-ciklus vázlata.}
\end{figure}

A körfolyamat végén a $^{12}_6 C$ magot ismét visszakapjuk, de közben négy $^{1}_1 H$-ból egy $^{4}_2 He$, két pozitron és két neutrínó keletkezett (ugyanúgy, mint a pp-ciklusban).
A CNO-ciklusban a béta-bomlások, valamint a reakcióban résztvevő atommagok nagyobb rendszáma miatt magasabb Coulomb-gátak, ill. nagyobb aktiválási energiák játsszák a ”fék”
szerepét. A CNO ciklusnak is vannak mellékágai, azonban ezekkel itt most nem foglalkozhatunk.

Részletek a fúziós folyamatokról: \href{https://mek.oszk.hu/16000/16052/16052.pdf#section.10.3}{Kisérleti magfizika - Fúziós folyamatok}

\subsubsection{Radioaktivitás}
 XIX. század végére a természetről alkotott ismeretek sokak számára teljesnek, befejezettnek tűntek. A kémia megállapította, hogy vannak olyan anyagok, amelyeket semmilyen módon nem lehetett átalakítani, szétbontani vagy más anyagokból létrehozni, ezek
voltak az elemek. Minden további anyag ezen elemek egyszerű szabályok szerinti kapcsolódásából állt össze (egyszeres és többszörös súlyviszonyok törvénye, John DALTON,
1766-1844). Az elemek tovább nem bontható, átalakítható apró részeit a kémikusok atomoknak nevezték, az összetett anyagok legapróbb részeit pedig molekuláknak. Amadeo
Lorenzo AVOGADRO (1776-1856) 1811-ben megfogalmazta a róla elnevezett tételt: azonos nyomás, térfogat, hőmérséklet mellett a gázok azonos számú részecskét (molekulát,
ill. atomot) tartalmaznak, anyagi minőségtől függetlenül. Az Avogadro-féle szám fontos természeti állandó, amely valamely kémiai elem vagy vegyület egy molekulasúlynyi
(mólnyi) mennyiségében levő részecskék számát jelenti. $N_A = 6.02 \cdot 10^{23}$. Jóllehet az atomokat és a molekulákat - közvetlen bizonyítékok hiányában - sokan még nem fogadták el ténylegesen létezőknek, az anyagok atomos/molekuláris szerkezetének feltételezése
nagyon sikeres modell volt. A szén égése a szénatomok és az levegő oxigén-molekuláinak széndioxiddá történő egyesülésével magyarázható: $C+O_2 \rightarrow CO_2$. Az Avogadro szám
ismeretében azt is meg lehetett határozni, hogy egyetlen szénatom elégetésekor mekkora energia szabadul fel:
\begin{equation}
    E_C = \frac{Q}{N_A} \approx 0.7 \cdot 10^{-18} J
\end{equation}
ahol $Q \approx 400 kJ$ a mólnyi szén égéshője. A kémiai átalakulások során történő energia-
változások valamennyien $10^{-18} J$ nagyságrendjébe esnek.

A XIX. század utolsó éveire a fizika már lefektette az energia-megmaradás törvényét,
valamint a hőtan törvényeit. A mechanika, a fénytani törvényei szintén kikristályosodtak,
és James Clark MAXWELL (1831-1879) egységbe foglalta az elektromosságtan és a mágnesség törvényeit is.

Ebbe a világképbe robbant bele 1896-ban Henri Antoine BECQUEREL (1852-1908)
francia kutató felfedezése, hogy egyes elemek - külső behatás és külső energiaforrás nélkül
- energiát hordozó sugárzást bocsátanak ki. Ez egyszerre mondott ellent az addigi kémiai
ismereteknek (mert az átalakulás során a korábban változhatatlannak hitt kémiai elemek
alakulnak át egymásba), és az energia-megmaradásnak, hiszen ennek az átalakulásnak a
során látszólag a semmiből keletkezik energia, amelyet az anyag nagy energiájú részecsk-
ék (alfa- béta- gamma-sugárzások) formájában kisugároz. Ezt a különleges jelenséget nevezték el radioaktivitásnak.

\subsubsection*{A Radioaktivitás felfedezése}
1895-ben Wilhelm Conrad RÖNTGEN (1845-1923, Nobel-díj 1901) felfedezte, hogy
léteznek olyan sugárzások, amelyek anyagokon áthatolnak (Röntgen-sugarak). Ezt követően a kutatók azt kezdték keresni, hogy vajon vannak-e olyan foszforeszkáló ill. fluoreszkáló anyagok, amelyek ilyen nagy áthatolóképességű sugarakat bocsátanak ki magukból?
Becquerel uránsók foszforeszkáló tulajdonságait vizsgálta 1896-ban, és nagy áthatolóképességű sugárzásokat keresett. Azt a meglepő felfedezést tette, hogy az urán külső
behatás és energiaforrás nélkül is nagy áthatolóképességű sugárzást bocsát ki, amely a
fekete papírba csomagolt fényképezőlemezt is megfeketíti. Becquerel felfedezésének ma
több közvetlen gyakorlati felhasználása is van: ezen az elven működnek a sugárveszélyes helyen dolgozókat ért sugárdózis ellenőrzésére szolgáló filmdoziméterek, valamint
ezt használja az autoradiográfia módszere is.

\subsection{Sugárzás és anyag kölcsönhatása: Bethe–Bloch-formula}
A radioaktív sugárzások energiája általában sokkal nagyobb, mint az anyagban lévő atomok vagy molekulák ionizálásához szükséges energia. Ezért ezek a sugárzások az
anyagban megtett útjuk során ionokat keltenek, vagy legalábbis kelthetnek. Ionizáló sugárzásoknak mindazokat a sugárzásokat nevezzük, amelyek - közvetve vagy közvetlenül
- az anyagban ionokat hozhatnak létre. Az elektromosan töltött részecskék (protonok,
alfa-sugárzás, béta-sugárzás, stb.) az anyagban haladva közvetlenül ionizálnak. Az elektromos töltés nélküli részecskék (pl. röntgen- vagy gamma-sugárzás, ill. neutronok) csak
közvetett módon ionizálnak: először valamilyen kölcsönhatás révén elektromosan töltött
részecskéket keltenek, illetve ilyeneknek adják át energiájukat, és ezek ionizálják tovább
az anyagot. Ezért először az elektromosan töltött részecskék anyaggal történő kölcsönhatásával foglalkozunk.

\subsubsection{Elektromosan töltött részecskék kölcsönhatása az anyaggal}
Az anyagok szerkezete elektromosan töltött részecskékből áll: pozitív elektromos töltésű,
kis kiterjedésű atommagok, amelyeket negatív töltésű elektronok vesznek körül (atomok),
ill. tartanak össze (molekulák, fémek, szilárd anyagok). Az atommagok térfogata az
anyag teljes térfogatának mindössze $10^{-15}$ része. Ebből azonnal következik, hogy bár a
töltött részecskék az atommagokkal is kölcsönhatásba lépnek (pl. Rutherford-szórás),
ez a kölcsönhatás azonban általában elhanyagolható valószínűségű az elektronokkal való
kölcsönhatáshoz képest. Elektromos töltések között a jól ismert Coulomb-kölcsönhatás
hat, azonos előjelű töltések taszítják, különböző előjelűek pedig vonzzák egymást.
Ha tehát egy nagy energiájú elektromosan töltött részecske behatol az anyagba, a Coulombkölcsönhatás révén energiát ad át az anyagban lévő elektronoknak.

Az energia megmaradása miatt a bejövő részecske energiája csökken, miközben energiát ad át az anyagnak.
Minél messzebb haladt egy részecske az anyagban, annál kisebb és kisebb lesz a részecske energiája.
Ez az energiacsökkenés addig folytatódik, amíg a részecske végül meg nem áll, ill. amíg hőmérsékleti egyensúlyba nem kerül a közeggel (átlagos mozgási energiája egyenlő nem lesz a hőmozgás energiájával).
Az energiaátadásnak a következő formái lehetnek:
\begin{itemize}
    \item gerjesztés (az atomokban, molekulákban egyes elektronok magasabb energiájú állapotba kerülnek)
    \item ionizáció (elektronokat leszakít, és ionok jönnek létre)
    \item rugalmas ütközés az atomokkal (az atomok mozgási energiát nyernek a bejövő részecskétől, a közeg felmelegszik)
    \item fékezési sugárzás (nehéz atommagok terében a könnyű töltött részecskék nagy gyorsulásra tesznek szert, és ennek következtében elektromágneses sugárzást bocsátanak ki)
\end{itemize}

\subsubsection{Lineáris energiaátadás}
Minden E energiájú, elektromosan töltött részecske (pl. alfa-részecske, proton, elektron,
stb.) az anyagon való áthaladásakor energiát ad át a közegnek. A hosszegységre jutó
átadott energiát $-dE dx$ lineáris energiaátadásnak (angolul: linear energy transfer, LET)
nevezzük. Természetesen, – az energia-megmaradás törvénye miatt – amekkora ener-
gia átadódik a közegnek, ugyanannyi energiát veszít a bejövő részecske. Emiatt ezt a
mennyiséget lineáris (vagy fajlagos) energiaveszteségnek is szokás nevezni. (A definíció
óban azért szerepel negatív előjel, mert a dE a részecske energiájának a megváltozása,
ami negatív. Így a LET — ami az anyagnak átadott energia – végülis pozitív lesz).

Először az elektronnal jóval nehezebb töltött részecskékkel (alfa-részecske, proton, hasadvány, nehézion stb.) foglalkozunk. A legnagyobb energiaátadás centrális egyenes ütközéskor történik,
ezért tegyük fel, hogy egy $E_{kin}$ mozgási energiájú, $M$ tömegű részecske
tökéletesen rugalmasan, egyenesen ütközik egy $m$ tömegű elektronnal. Az energia- és
lendület-megmaradás:
\begin{equation}
    \begin{aligned}
        &p(M) = p'(M) + p'(m) \\
        &E_{kin}(M) = \frac{p(M)^2}{2M} = \frac{p'(M)^2}{2M} + \frac{p(m)^2}{2m} = E'_{kin}(M) + E'_{kin}(m) \\
    \end{aligned}
\end{equation}
Ebből rövid számolással kapjuk, hogy a részecske mozgási energiájának megváltozása:
\begin{equation}
    \Delta E_{kin}(M) = E'_{kin}(M) - E_{kin}(M) = -\frac{4 m M}{(m+M)^2} E_{kin}(M)
\end{equation}
Mivel $M \gg m$, ezért $\Delta E_{kin}(M) \approx \frac{4m}{M} E_{kin}(M)$. Egy kb. 5 MeV energiájú alfa-részecske ezek
szerint egyetlen ütközéskor legfeljebb $2.7 \text{ keV}$ energiát képes veszíteni. Ha állandóan csak
egyenesen ütközne az elektronokkal, akkor is majdnem húszezer ütközés kellene a teljes
lelassulásig (és ennél a becslésnél még nem vettük figyelembe azt, hogy a későbbi ütközésekor
már csökken a képletben szereplő $E_{kin}$, tehát még egyenes (frontális) ütközéskor
sem tud annyi energiát veszíteni, mint az első ütközéskor). Ebből néhány következtetést
azonnal levonhatunk:
\begin{itemize}
    \item Sok tízezer ütközés kell, amíg egy nehéz részecske átadja az összes energiáját az anyagnak. Más szóval: egy-egy ütközéskor átadott energia kicsi.
    \item Az elektronnal való ütközéskor a nehéz részecske lendületváltozása igen kicsiny, ezért a részecske pályája az anyagban többé-kevésbé egyenes.
    \item Mivel a Coulomb-erő hossz hatótávolságú, ezért a részecske egyszerre sok elektronnal áll kölcsönhatásban, ennek következtében az energiáját folytonos ütemben adja át. Egy bizonyos távolság megtétele után, amikor a teljes energiáját átadta, a részecske ”megáll” (pontosabban a közeg hőmérsékletének megfelelő hőmozgást végez). Azt a távolságot, amikor ez bekövetkezik, a részecske behatolási mélységének nevezzük.
    \item Az atomok, molekulák ionizációs energiája a 10-100 eV tartományba esik. Ezért pl. a fenti példában a 2,7 keV energiájú kilökött elektron további atomokat tud ionizálni, és ezáltal másodlagos ionizáció jön létre az anyagban.
    \item Sok olyan kölcsönhatás is van, amelyben az átadott energia kisebb, mint az atom ionizációs energiája. Ennek következtében az anyagban gerjesztett atomok keletkeznek (amelyek azután rövidebb-hosszabb idő alatt visszabomlanak az alapállapotba).
\end{itemize}

Az anyagnak átadott energia leírásánál tehát ezeket a folyamatokat mind figyelembe
kell venni.

Az anyagban bekövetkező szóródási folyamatok kvantummechanikai tárgyalásával a
közegben haladó, nehéz töltött részecskékre (proton, alfa-részecske, stb.) a lineáris ener-
giaátadást közelítőleg a következő képlet alapján lehet kiszámítani:
\begin{equation}
    \boxed{-\frac{dE}{dx} = \text{const.} \cdot \frac{z^2}{v^2} \cdot \left(\frac{Z \cdot \rho}{A}\right) \cdot K}
\end{equation}
Ezt \textbf{Bethe-Bloch formulának} nevezik kidolgozóik, H. BETHE (1906-2005, Nobel-díj 1967),
F. BLOCH (1905-1983, Nobel-díj 1952) után. Az egyenlet jobb oldalán az első zárójel-
ben a bejövő részecske jellemző adatai vannak (z a töltése elemi töltés egységekben, $v$
a sebessége), a második zárójel a közeg jellemzőit tartalmazza ($\rho$ a sűrűség $[kg/m^3]$, $Z$ az
anyag átlagos rendszáma, $A$ pedig az átlagos móltömege
$[ kg/mol ]$), $K$ pedig egy, a nem-relativisztikus tartományban lassan változó függvény. Az egyenletben szereplő konstans
természeti állandókat, és egységválasztástól függő konstansokat tartalmaz:
\begin{equation}
    \text{const.} = \left(\frac{e^2}{4 \pi \epsilon_0}\right)^2 \cdot \frac{4 \pi N_0}{m_e}
\end{equation}
Itt $e$ az elemi töltés, $N_0$ az Avogadro-állandó, $m_e$ pedig az elektron nyugalmi tömege. A $K$ függvény függ a közeg atomjainak átlagos gerjesztési potenciáljától,
tartalmaz anyagsűrűség-korrekciót, héjkorrekciót és relativisztikus korrekciót:
\begin{equation}
    K = \ln \left(\frac{2 m_e c^2 \beta^2}{I}\right) - \ln(1 - \beta^2) - \beta^2
\end{equation}
Ebben a képletben $\beta = \frac{v}{c}$, a részecske sebességének viszonya a fénysebességhez, $I$ pedig a közeg ionizációjára jellemző energia, amelyet empirikusan szoktak meghatározni.
Általában $I \approx Z \cdot 10 \text{ eV}$. Levegőben például $I = 86 \text{ eV}$, alumíniumban pedig $I = 163 \text{ eV}$.

A Bethe-Bloch formulát vizsgálva meglepő, hogy a lineáris energiaátadás a bejövő
részecske sebességének négyzetével fordítottan arányos, azaz kétszer nagyobb sebességű
részecskék egységnyi kis úton négyszer kevesebb energiát adnak le. Ennek eredményeképpen egy nagy sebességű részecske az anyag felszíne közelében hosszegységenként csak
kis energiát ad le, azaz az energiája is lassan csökken. Ahogy azonban csökken a sebessége, egyre jobban és jobban ionizál, és végül szinte hirtelen megáll. A legtöbb energiát
a megállási pont környezetében adja le.

\subsubsection{Behatolási mélység}
Azt a távolságot, ameddig a sugárzás részecskéi egyáltalán eljutnak az anyagban, a sugárzás behatolási mélységének nevezzük. Ezt matematikailag a Bethe-Bloch formula inverzének integrálásával kaphatjuk:
\begin{equation}
    R = \int_{E_0}^{0} \left(-\frac{dE}{dx}\right)^{-1} dE
\end{equation}
Essenek részecskék egy alfa-sugárforrásból egy anyag felületére. Az ábra azt mutatja
meg, hogy hány részecske jut el az anyag felszínétől adott távolságra. Az ábra alapján az alfa-részecskék legnagyobb része eljut a behatolási mélység környékére, és ott áll meg. A
görbe azonban nem hirtelen esik le nullára. Ennek az az oka, hogy statisztikus folyamatok
miatt nem minden alfa-részecske halad pontosan ugyanolyan mélyre. Szigorúan véve azt szokás behatolási mélységnek nevezni, amelybe éppen a részecskéknek a fele jut el.
\begin{figure}[H]
    \centering
    \includegraphics[width=0.6\textwidth]{imgs/13-tetel/behatolasi_melyseg.png}
    \caption{Alfa-részecskék behatolási mélysége egy anyagban.}
\end{figure}
\textit{Megjegyzés:} Vannak olyan sugárzások (pl. gamma-sugárzás), ahol a behatolási mélység nem értelmes fogalom, mert nincs olyan ”keskeny” réteg, amelyben a sugárzás részecskéi megállnának. Ilyen esetben a felezési rétegvastagság fogalmát szokás használni.
A behatolási mélység fogalma jól használható az erősen ionizáló részecskéknél (proton,
alfa-részecske, nehézion), a felezési rétegvastagság pedig a gyengén ionizáló, és ezért nagy
áthathatolóképességű sugárzásoknál.

Egyes részecskék behatolási mélységéről bővebben: \href{https://mek.oszk.hu/16000/16052/16052.pdf#subsection.5.1.3}{Kisérleti magfizika - Sugárzás és anyag kölcsönhatása: Bethe–Bloch-formula}

\subsection{Radioaktív bomlások, magátalakulások}
\subsubsection{Alfa- béta- gamma-sugárzások}
A radioaktivitás felfedezését követően sok kutató bekapcsolódott a radioaktív sugarak
vizsgálatába. Sir Ernest RUTHERFORD (1871-1937, kémiai Nobel-díj 1908) angol fizikus 1898-ban megállapította, hogy a sugárzás nem homogén, hanem van egy erősen
ionizáló összetevője, amelyet azonban már egy nagyon vékony anyagréteg (pl. papír) is
elnyel, és amelynek a levegőben a hatótávolsága csak néhány centiméternyi. Van azonban
egy másik összetevő is, amelyik sokkal kevésbé ionizál, de sokkal nagyobb áthatolóképességgel rendelkezik. Rutherford ezeket az összetevőket a görög ábécé első két betűje
nyomán alfa- ($\alpha$) ill. béta- ($\beta$) sugaraknak nevezte el. Két évvel később Paul Ulrich
VILLARD (1860-1934) francia fizikus felfedezte a béta-sugaraknál is nagyobb áthatolóképességű sugárzást, ami a gamma-sugárzás ($\gamma$) nevet kapta. A vizsgálatok
\begin{figure}[H]
    \centering
    \includegraphics[width=0.6\textwidth]{imgs/13-tetel/radioaktiv.png}
    \caption{Sugárzások áthatolóképessége}
\end{figure}
ez után arra irányultak, hogy megállapítsák e sugarak természetét. Becquerel megállapította, hogy a béta-sugarak negatív elektromos töltést hordoznak, és a tömegük és
töltésük viszonya megegyezik az elektronokéval. A béta-sugárzás tehát nagysebességű
elektronokból áll, ezért mind az elektromos, mind a mágneses mező eltéríti.
1909-ben Rutherford állapította meg, hogy az alfa-részek tulajdonképpen nagy energiájú
héliumionok. Mivel ezek is elektromos töltésű részecskék, ezeket is eltéríti mind az elektromos, mind a mágneses mező, csak éppen ellenkező irányba, mint a béta-sugarakat. A
gamma-sugárzásról pedig kiderült, hogy nagyon kis hullámhosszúságú elektromágneses
sugárzás, ezért ezeket sem elektromos, sem mágneses mezővel nem lehet eltéríteni.

Az elektromos és mágneses mezőben való eltérülésből az elektromos töltésű részecskék energiája is meghatározható. Nagy megdöbbenést keltett, amikor kiderült, hogy ezek
az energiák is meghatározható. Nagy megdöbbenést keltett, amikor kiderült, hogy ezek
az energiák több nagyságrenddel – néha milliószor – nagyobbak, mint a kémiai átalakulásoknál fellépő energiák ($E \approx 10^{-18} J$). Ma már tudjuk, hogy
radioaktív bomláskor az atomok középpontjában lévő atommagok alakulnak át. Az őket
alkotó protonok és neutronok átrendeződnek, és erősebben kötött állapotba kerülnek.
\begin{figure}[H]
    \centering
    \includegraphics[width=0.3\textwidth]{imgs/13-tetel/radioaktiv2.png}
    \caption{$\alpha$, $\beta$ és $\gamma$ sugárzások eltérülése elektromos térben.}
\end{figure}
A radioaktív sugárzás részecskéi az átalakulás során felszabaduló kötési energiából nyerik hatalmas energiájukat. A radioaktivitás nem más, mint az anyag ”vándorlása” az
rik hatalmas energiájukat. A radioaktivitás nem más, mint az anyag ”vándorlása” az atommagok energiavölgyében.
Ezekben az átalakulásokban – bármilyenek legyenek is azok, – bizonyos fizikai mennyiségekre megmaradási törvények állnak fenn. Egy fizikai mennyiségre vonatkozó megmaradási törvény azt jelenti,
hogy a végállapotban (pl. a bomlás után) a részecskék megfelelő fizikai mennyiségeit összeadva ugyanannyit kell kapjunk, mint amennyi a kezdőállapotban (a bomlás előtt) volt. A fizikában sok megmaradási törvény van, ezek közül
most csak hármat sorolunk fel, amelyeknek a radioaktív bomlások leírásánál jelentős szerepe van.
\begin{itemize}
    \item az elektromos töltés megmaradása
    \item a nehéz-részecske szám (az ún. bariontöltés) megmaradása és a
    \item könnyű-részecske szám (leptontöltés) megmaradása.
\end{itemize}

Egy protonokból és neutronokból álló atommag paramétereit a benne található neutronok és protonok megfelelő paramétereinek összegzésével kapjuk meg. Mivel mind a
protonok mind a neutronok leptontöltése nulla, ezért az atommagoké is nulla, akárhány
protont vagy neutront is tartalmaznak. Egy Z protonból és N neutronból álló atommag
elektromos töltése $+Z$ (elemi töltés), bariontöltése pedig $A = Z + N$.
\begin{table}[H]
    \centering
    \begin{tabular}{|c|p{3cm}|p{3cm}|p{3cm}|}
        \hline
        Részecske & Elektromos töltés (elemi töltés egység) & Barion töltés (nehézrészecske-szám) & Lepton töltés (könnyűrészecske-szám) \\
        \hline
        Proton & +1 & +1 & 0 \\
        \hline
        Neutron & 0 & +1 & 0 \\
        \hline
        Elektron & -1 & 0 & +1 \\
        \hline
        Pozitron & +1 & 0 & -1 \\
        \hline
        Neutrínó & 0 & 0 & +1 \\
        \hline
        Antineutrínó & 0 & 0 & -1 \\
        \hline
    \end{tabular}
    \caption{Elektromos töltés, bariontöltés és leptontöltés különböző részecskéknél.}
\end{table}

\subsubsection{Alfa-részecske, alfa-sugárzás}
Az alfa-bomláskor kibocsátott részecskét alfa-részecskének nevezzük. Ez tulajdonképpen
hélium-atommag, amelyben 2 proton és 2 neutron van. Elektromos töltése tehát +2,
és összesen 4 nehéz részecske van benne (bariontöltése 4). Ezért az alfa-bomlás során
egy $Z$ protont (rendszám) és $A$ nehéz részecskét (tömegszám) tartalmazó atommagból
$Z - 2$ protont (töltésmegmaradás) és $A - 4$ nehéz részecskét (barionszám-megmaradás)
tartalmazó atommag jön létre. Képletben:
\begin{equation}
    ^A_Z X \quad \rightarrow \quad ^{A-4}_{Z-2} Y + ^4_2 He
\end{equation}
Ezek valamennyien ”nehéz” részecskék, a leptontöltést (könnyűrészecske-számot) tehát
nem változtatják meg. Az elektromos-töltés megmaradását az ”alsó indexek”, a barionszám-
megmaradását a ”felső indexek” alapján lehet ellenőrizni.

\subsubsection{Beta-részecske, beta-sugárzás}
A béta-bomlásoknak három fajtája van: negatív és pozitív béta-bomlás, valamint elektron-
befogás.

\textbf{Negatív béta-bomlás:} A legkorábban felfedezett béta-bomlás során elektron keletkezik. Az elektron negatív elektromos töltésű, könnyű részecske: $e^- = ^0_{-1} e$. Ezért sokáig a következőképpen
képzelték el a béta-bomlás lefolyását:
\begin{equation}
    ^A_Z X \quad \rightarrow \quad ^A_{Z+1} Y + ^0_{-1} e
\end{equation}
Látható, hogy az elektromos töltés (alsó indexek) és a bariontöltés (felső indexek) megmarad. A könnyűrészecske-szám megmaradási törvény azonban nem teljesül, hiszen a
jobb oldalon az elektron miatt a könnyűrészecske-szám eggyel nagyobb, mint a baloldalon. Ezért ki kell bocsátódjon egy olyan könnyű részecske, amely
elektromosan semleges (tehát nem ”rontja el” az elektromos töltésmegmaradást), viszont
a könnyűrészecske száma $-1$. Az Táblázat szerint az antineutrínónak vannak ilyen
tulajdonságai. A negatív béta-bomlás tehát helyesen:
\begin{equation}
    ^A_Z X \quad \rightarrow \quad ^A_{Z+1} Y(0) + ^0_{-1} e(1) + \tilde{\nu} (-1)
\end{equation}
A zárójelben lévő szám az illető részecske leptontöltését (könnyűrészecske-számát) jelzi.
A leptonszámot általában nem szokás jelölni, mi is csak most egyszer, az érthetőség
kedvéért jeleztük. A negatív béta-bomlás szokásos jelölése:
\begin{equation}
    ^A_Z X \quad \rightarrow \quad ^A_{Z+1} Y + e^- + \tilde{\nu}
\end{equation}

\textbf{Pozitív béta-bomlás:} 1932-ben felfedezték az elektron antirészecskájét, a pozitront.
Ezt követően a mesterségesen előállított radioaktív anyagok bomlásai között hamarosan találtak pozitív béta-bomlási folyamatokat is. 
\begin{equation}
    ^A_Z X \quad \rightarrow \quad ^A_{Z-1} Y + e^+ + \nu
\end{equation}
Az itt szerepl˝o $e^+$ könnyű részecske egyszeres pozitív elektromos töltésű, és az elektronnak az antirészecskéje. A neve: pozitron. A $\nu$ részecske elektromosan semleges
könnyű részecske, a neve: \textit{neutrínó}.

\textbf{Elektron-befogás:} Földi világunkban az atommagot sűrű elektronfelhő veszi körül.
Ezért – energetikailag kedvező esetben – az atommag be tud fogni egy elektront az atomi
elektronhéjból. A megmaradási törvények miatt ez a következőképpen mehet csak végbe:
\begin{equation}
    ^A_Z X + e^- \quad \rightarrow \quad ^A_{Z-1} Y + \nu
\end{equation}
Ezt a folyamatot elektronbefogásnak nevezzük. Vegyük észre, hogy ebben a folyamatban csak egy neutrínó lép ki (a megmaradt atommag mellett, természetesen0). Mivel a
neutrínót roppant nehéz érzékelni, tehát ezt a folyamatot közvetlenül nem tudjuk megfigyelni. Közvetett megfigyelésre azonban nyílik mód, hiszen az atommag által elnyelt
elektron helye ”üresen” marad az atomi elektronhéjban. Ebbe egy külső (magasabb
energiájú) pályáról beugrik egy elektron, és közben az atom a többlet-energiát karakterisztikus röntgen-sugárzás formájában kisugározza. Ekkor azonban az atommag töltése
már $Z - 1$, tehát a karakterisztikus röntgensugárzás hullámhossza is ilyen rendszámú
atomra lesz jellemző. Amikor tehát egy $Z$ rendszámú anyagdarabból $Z - 1$ rendszámú
atomokra jellemző karakterisztikus röntgensugárzást észlelünk, akkor tulajdonképpen az
elektronbefogásokat detektáljuk közvetett módon.

A béta-bomlások valamennyien az energiavölgy ”oldalán” lezajló folyamatok, az $A = \text{konstans}$ izobárokat leíró energia-parabolák mentén. Ezért a
negatív béta-bomlások azokra az atommagokra jellemzők, amelyekben túl kevés proton
van az adott tömegszámnak megfelelő $Z_{min}$-hoz képest, a pozitív béta-bomlások és az
elektronbefogások pedig azokra az atommagokra, amelyekben $Z_{min}$-nál több proton van.

Természetesen, ugyanezt a neutronok számával is megfogalmazhatjuk. Adott A tömegszámhoz egy optimális $N_{opt} = A - Z_{min}$ neutronszám tartozik. Ilyen neutronszámú
atommagok a legerősebben kötöttek az izobárok közül. Negatív béta-bomlások az ehhez képest túl sok neutront tartalmazó atommagokra, a pozitív béta-bomlások és az
elektron-befogások pedig a túl kevés neutront tartalmazó atommagokra jellemzőek.

\subsubsection{Gamma-bomlás}
A gamma-bomlás során az atommag elektromágneses sugárzást (fotont) bocsát ki. Mivel a kisugárzott fotonnak mind az elektromos-, mind a barion-, mind pedig leptontöltése
nulla, ezért a gamma-bomlás során az atommag semmilyen összetevője nem változik meg. Azaz 
\begin{equation}
    ^A_Z X^* \quad \rightarrow \quad ^A_Z X + \gamma
\end{equation}
Az $X^*$ atommag egy magasabb energiájú (gerjesztett) állapotából (ezt jelzi a csillag) egy
alacsonyabb energiájú állapotába megy át, és az energiakülönbséget elektromágneses sugárzás formájában kibocsátja. A sugárzás frekvenciájára érvényes a Planck- összefüggés:
\begin{equation}
    h\nu = E_{\text{magasabb}} - E_{\text{alacsonyabb}}
\end{equation}
Itt $\nu$ a kibocsátott elektromágneses sugárzás frekvenciája, $E_{\text{magasabb}}$ és $E_{\text{alacsonyabb}}$ az
atommag kezdeti ill. végállapotának energiája, $h$ pedig a Planck-állandó.
Ha a gamma-bomlást az energiavölgyön szeretnénk ábrázolni, akkor úgy képzelhetjük,
hogy a gerjesztett atommag a völgy ”fölött” helyezkedik el, hiszen a völgyön az alapállapotú atommagok vannak! A gamma-bomlás során tehát a völgy fölött lévő, gerjesztett
állapotban lévő atommag az összetételének változtatása nélkül lejjebb kerül, közelebb a
völgyhöz, de akár rá is kerülhet a völgyre, amikor alapállapotba kerül.
Magasan gerjesztett atommag általában valamilyen más magátalakulásban (pl. béta-bomlásban) keletkezik. Ezért a gamma-bomlás általában valamilyen más magátalakulást
követő bomlási mód.

\subsubsection{A radioaktivitás jellemző mennyiségei}
\subsubsection*{Aktivitás}
A radioaktív anyagok mennyisége időben fokozatosan csökken, hiszen a bomlásra képes atommagok egy része elbomlik. Az időegység alatt bekövetkező bomlások számát aktivitásnak nevezzük.
Jelöljük $N(t)$-vel a bomlásra képes atommagok számát a $t$ időpillanatban, ekkor
\begin{equation}
    a = - \frac{dN}{dt}
\end{equation}
A bomlásra képes atommagok száma a bomlás miatt csökken, ezért $\frac{dN}{dt}$ negatív. A fenti definícióval az \(a\) aktivitás pozitív értékű. Az aktivitás egysége a becquerel (ejtsd:
bekerel).
\begin{equation}
    1 \text{ Bq} = 1 \frac{\text{bomlás}}{\text{másodperc}} = 1 \text{ s}^{-1}
\end{equation}
Az aktivitás régi egysége a curie (ejtsd: küri). $1 Ci = 3, 7 \cdot 10^{10} \text{ Bq}$.

\textit{Megjegyzés:} Maria SKLODOWSKA-CURIE (1867-1934, Fizikai Nobel-díj 1903, Kémiai Nobel-díj 1911) lengyel származású fizikusnő, aki férjével, Pierre CURIE-vel (1859-
1906, Fizikai Nobel-díj 1903) együtt fedezte fel és állította elő a rádium és polónium
elemeket. Az aktivitás régi egységét róluk nevezték el. 1 g rádiumnek éppen 1 Ci
($=3, 7 \cdot 10^{10} \text{ Bq}$) aktivitása van.

Radioaktív forrás aktivitása egyenesen arányos az anyagban lévő, bomlásra képes
(még el nem bomlott) atommagok számával (N):
\begin{equation}
    a = \lambda \cdot N
\end{equation}
$\lambda$ neve: bomlásállandó, mértékegysége: 1/s.
A mag- és részecskefizikában gyakran használják a bomlásállandó reciprokát:
\begin{equation}
    \tau = \frac{1}{\lambda}
\end{equation}
$\tau$ idő [s] dimenziójú, és a a bomló magok átlagos élettartamát adja meg (ld. 3.1.
Feladat)

\textit{Megjegyzés:} A kvantummechanika szerint véges élettartamú állapot energiája nem
lehet teljesen határozott: $\tau \cdot \Delta E \geq \frac{\hbar}{2}$. Egyes könyvek ezt az energiára és időre vonatkozó
Heisenberg-féle határozatlansági relációnak hívják - helytelenül.

Az aktivitás az élettartammal is kifejezhető, mint az egyszerűen látható:
\begin{equation}
    a = \frac{N}{\tau}
\end{equation}

Mivel a bomlásállandó kapcsolatban van a $T$ felezési idővel, az aktivitás még következőképpen is írható:
\begin{equation}
    a = \frac{N \cdot \ln(2)}{T}
\end{equation}

\subsubsection*{Exponenciális bomlástörvény}
A radioaktív bomlás statisztikus jelenség, ezért törvényei gyakorlatilag csak nagyszámú
atommag bomlására alkalmazhatók, egyes atommagok bomlására csak valószínűségi ki-
jelentéseket tehetünk. Annak a valószínűsége például, hogy egy kiszemelt atommag a $T$
felezési idő alatt elbomlik éppen 1/2. Újabb $T$ idő alatt bekövetkező bomlás valószínűsége ismét 1/2, attól függetlenül, hogy az atommag mióta ”vár” már a bomlásra.
Az el nem bomlott atommagok $N$ számára elsőrendű differenciálegyenletet kapunk:
\begin{equation}
    -\frac{dN}{dt} = \lambda N(t)
\end{equation}
Ennek a differenciálegyenletnek a megoldása az \textbf{exponenciális bomlástörvény}:
\begin{equation}
    N(t) = N_0 \cdot e^{-\lambda t}
\end{equation}
Az $N_0$ konstans értelmezése: a függvény értéke a $t = 0$ kezdeti időpontban (a mi
esetünkben a kezdetben meglévő radioaktív atommagok száma).
Ennek az időfüggvénynek a menetét mutatja az ábra:
\begin{figure}[H]
    \centering
    \includegraphics[width=0.4\textwidth]{imgs/13-tetel/exponencialis_bomlas.png}
    \caption{Exponenciális bomlástörvény.}
\end{figure}

\textbf{Felezési idő:}
Az exponenciális bomlástörvény egy másik, gyakran használt alakját úgy kapjuk meg,
ha az $e$ alapról ( $= 2.718281828459...$, a természetes logaritmus alapszáma) 2-es alapra
térünk át:
\begin{equation}
    N(t) = N_0 \cdot 2^{-\frac{t}{T}}
\end{equation}
Az ebben az alakban szereplő T felezési idő és az előző képletben szereplő $\lambda$ bomlásállandó
nem függetlenek egymástól:
\begin{equation}
    T = \frac{\ln(2)}{\lambda}
\end{equation}
mint ahogyan azt könnyen igazolhatjuk.

Az aktivitásnál bevezetett bomlásállandónak ($\lambda$) érdekes értelmezést adhatunk. Legyen $\mu$ annak a valószínűségsűrűsége, hogy egyetlen atommag időegység alatt elbomoljon!
Ekkor az illető atommag dt idő alatt $\mu\cdot dt$ valószínűséggel bomlik el. Azt, hogy egy kiszemelt atommag pontosan mikor fog elbomlani, nem lehet megmondani, de ha nagyszámú
atommagot vizsgálunk, akkor meg lehet mondani, hogy nagyon sok ($N$) közül átlagosan
hány bomlik el dt idő alatt: $N \cdot \mu \cdot dt$. Ennyi idő alatt tehát ennyivel csökken a meglévő
atommagok száma, azaz az atommagok számának dN megváltozása: $dN =-N  \mu \cdot dt$.
(A negatív előjel azt fejezi ki, hogy az atommagok száma csökken.) Ez az egyenlet
összefüggést ad az atommagok számának megváltozása és a meglévő $N$ szám között:
\begin{equation}
    -\frac{dN}{dt} = \mu \cdot N(t)
\end{equation}
Ezt az egyenletet összehasonlítva az exponenciális bomlástörvény differenciálegyenletével azonnal látható, hogy $\mu = \lambda$, azaz a bomlásállandó éppen az időegységre eső bomlási valószínűség.

\subsubsection*{Poisson eloszlás}
A radioaktív bomlások statisztikus folyamatok, ezért csak valószínűségi törvényekkel
írhatók le. Legyen $p$ annak a valószínűsége, hogy egy atom $t$ idő alatt elbomlik, ekkor
nyilván annak a valószínűsége, hogy nem bomlik el, $1-p$. A fent levezetett exponenciális
bomlástörvény alapján $1 - p = e - \lambda \cdot t$, azaz
\begin{equation}
    p = 1 - e^{-\lambda \cdot t}
\end{equation}

Annak a valószínűsége, hogy kezdetben meglévő $N$ atomból éppen $k$ atom bomlik el
a fenti $t$ idő alatt a binomiális eloszlással számítható ki.
\begin{equation}
    P(k, N, p) = \binom{N}{k} p^k (1 - p)^{N-k}
\end{equation}
hiszen ehhez az kell, hogy éppen $k$ részecske bomoljon el, és $N - k$ ne bomoljon el. A $k$
darab elbomló részecskét viszont tetszőlegesen kiválaszthatjuk a meglévő $N$ részecskéből,
ezért jelenik meg az $\binom{N}{k}$ tényező. Megmutatható, hogy $k$ várható értéke: $\langle k \rangle = N \cdot p$.
Levezethető, hogy nagy $N$ részecskeszámok és kis $p$ valószínűségek esetén a binomiális
eloszlás átmegy egy úgynevezett Poisson-eloszlásba, ha a határátmenet úgy történik,
hogy a várható érték $N \cdot p$ konstans marad. Ennek következtében annak valószínűsége
hogy $t$ idő alatt egy $a$ aktivitású radioaktív forrásban pontosan k bomlás következzen
be, Poisson-eloszlást követ, ha $a$ várható értéke ($a t$) elegendően nagy ($a t > 10$):
\begin{equation}
    P(k, a t) = \frac{(a t)^k}{k!} e^{-a t}
\end{equation}
Ez az összefüggés addig használható, amíg az aktivitás időben állandónak tekinthető,
azaz $t \ll T$ (itt $T$ a felezési idő). Hasonlóképpen, Poisson-eloszlást követ az is, hogy
egy detektorral (pl. GM-cső) $t$ idő alatt pontosan $k$ beütést érzékeljünk, ha az időegység
alatt várható beütésszám $a$. 

További részletekért bomlási sorokról és természetes radioaktivitásról lásd: \href{https://mek.oszk.hu/16000/16052/16052.pdf#section.3.3}{Kisérleti magfizika - Radioaktív bomlások}

\subsection{Elemi részecskék és alapvető kölcsönhatások \cite{kiserleti_atomfizika}}
Az elemi részecskéket és kölcsönhatásait a mai modern fizikában a Standard Modell írja le. A részecskefizika Standard Modellje az elektromágneses, a gyenge és erős kölcsönhatást,
valamint az alapvető elemi részecskéket leíró kvantumtérelmélet. Összhangban van a kvantummechanikával és a speciális relativitáselmélettel.
Majdnem minden kísérleti teszt igazolja jóslatait, a korábbi kivételek legjelentősebbikét, a Higgs-bozon létezését már igazolták. A modell közvetlen előzményei és részei az elektrogyenge kölcsönhatások modellje
(Glashow–Weinberg–Salam-modell) és az erős kölcsönhatások elmélete, a kvantum-színdinamika (QCD).

A modell alapvető részecskéi között vannak egyrészt az anyagot felépítő ún. anyagi részecskék. Ezek mind feles spinűek, azaz fermionok. A \textbf{kvarkok} és \textbf{leptonok} tartoznak ide.
A fermionokra érvényes Pauli-elv miatt nem omlik össze egy anyaghalmaz, hanem kénytelen egy bizonyos térfogatot kitölteni.
A kvarkok az erős és elektrogyenge kölcsönhatásban is részt vesznek, míg a leptonok csak az elektrogyengében.

A másik típusú részecskék a – kölcsönhatásokat – közvetítő részecskék. Az elektrogyenge kölcsönhatást közvetítő \textbf{foton}, két \textbf{W-bozon}, és a \textbf{Z-bozon} valamint az erős kölcsönhatást közvetítő \textbf{gluonok} tartoznak ide.
Ezek mind egyes spinűek, azaz bozonok, spinjük miatt vektorbozonoknak is mondjuk őket. Bozonokra nincs Pauli-kizárás, ezek közül akárhány lehet azonos kvantumállapotban,
ami nagyon sok bozonnal nagyon kifinomult kölcsönhatást tesz lehetővé.
\begin{figure}[H]
    \centering
    \includegraphics[width=0.5\textwidth]{imgs/13-tetel/standard_model.png}
    \caption{A standard modell részecskéi.}
\end{figure}

A harmadik típusú részecske a már kísérletileg is igazolt nulla spinű \textbf{Higgs-bozon}, amiből a standard modellben egy van, de a kiterjesztett elméletekben több.
A Higgs-bozon az önmagával és a többi részecskével való kölcsönhatásaival tömeget „kölcsönöz” az anyagi részecskéknek és a közvetítő részecskék egy részének. 

\subsubsection{Történelmi előzmény}
A részecskefizika múltja a kozmikus sugárzás tanulmányozásához nyúlik vissza. (Ugyanakkor felhasználta és felhasználja a magfizika fejlődése során
kialakult fogalmakat, az elméleti és technikai eszközöket is.) A XX. század elején többen megfigyelték , hogy a legjobban elszigetelt elektroszkópok
is idővel elveszítik töltésüket . Akik ezt nem egyszerű technikai effektusnak
akarták betudni, hanem mélyebben elgondolkoztak rajta , azok két lehetséges magyarázatot láttak. Az egyik az, hogy a földkéreg radioaktivitása
az oka az elektroszkópok kisülésének. A másik elképzelés merészebb volt:
azt tételezte fel, hogy a világűrből érkezik hozzánk ionizáló sugárzás, és ez
az oka az elektroszkópok töltésveszteségének. A két elképzelés között kísérletileg különbséget tudunk tenni, ti. az utóbbi elképzelés esetében a Föld
felszínétől magasabbra emelkedve a sugárzás intenzitásának nőnie kell, hiszen külső, Földön kívüli eredetű. 1912-ben V. F. Hess 5000 m magasságig
végzett ballonkísérleteket és megállapította, hogy valóban van egy sugárzás, ami Földünket kívülről éri, és amelynek intenzitása ennek megfelelően
a Föld felszínétől számított magassággal együtt jelentősen (exponenciálisan) nő. Ezt a sugárzást kozmikus sugárzásnak nevezték el.

Az egyik fontos kérdés, hogy mi a kozmikus sugárzás primer komponense, azaz, milyen részecskék érkeznek a kozmoszból Földünkhöz? Erre már régebben is választ tudtak adni, ma pedig az űrszondák jóvoltából egész
pontosan megmondhatjuk, hogy a kozmikus sugárzás fő primer komponensét protonok alkotják, mellettük azonban megtalálhatók - bár sokkal kisebb
számban - nehéz elemek atommagjai is, egészen a vasig és még azon túl is. Hogyan keletkezik a kozmikus sugárzás: erre még ma sem tudunk pontosan felelni. 1949-ben Fermi kidolgozott egy elképzelést, amelynek alapján
a kozmoszban található mágneses terek és azok változása mint gigantikus gyorsítók működnek és kölcsönöznek óriási energiát az egyes részecskéknek.
Az atmoszférát elérő részecskék kölcsönhatnak az atmoszférában található oxigén- és nitrogénmolekulákban levő atomok magjaival, és így szekunder komponenseket (elektronok,$\gamma$ - sugarak, $\mu$- és $\pi$-mezonok, neutrínók
($\nu$) hoznak létre ("zápor")).

Visszatérve a részecskékhez, úgy nézett ki egészen az 1930-as évekig, hogy a világot protonokból, neutronokból és elektronokból lehet ,,összerakni". Ha ezekhez még hozzávesszük a fotont, mint az elektromágneses
kölcsönhatások közvetítőjét, akkor mindössze négy elemi részecskéről beszélhetünk. Később az atommagszerkezetének megismerése során felmerült
az a kérdés, hogy mi az az erő, amely összetartja a protont és a neutront az atommagban? H. Yukawa, japán elméleti fizikus volt az, aki 1925-ben elméletileg megjósolta, hogy ugyanolyan közvetítő részecskét kell keresni a
nukleonok közötti erőhatások leírására, mint amilyen a foton az elektromágneses hatás közvetítésében. A nukleonok közötti erőhatás intenzitásának és
hatótávolságának empirikus ismeretében Yukawa kiszámította, hogy a közvetítő részecske nyugalmi tömege nem nulla, hanem az elektron tömegének durván kétszázszorosa.

\subsubsection*{A pozitron felfedezése}
1932-ben C. D. Anderson kaliforniai professzor méréseket végzett mágneses térbe helyezett Wilson-kamrával. A Wilson-kamrán áthaladó
kozmikus sugárzási nyomok között talált olyanokat, amelyek töltése egyértelműen pozitív volt. Ez önmagában nem lett volna meglepetés, a felfedezés
idején azonban a legkisebb tömegű egységnyi pozitív töltéssel rendelkező részecskeként a protont ismerték. Anderson vizsgálatai során kiderült, hogy
az 1300 felvétel közül az a 15 nyom, ami pozitív töltésűnek mutatkozott, olyan részecskéktől származik, amelyeknek a tömege biztosan jóval kisebb,
mint a protoné. Gondos analízis után végül is maradt az a magyarázat, hogy itt az elektronnal azonos tömegű, de ellenkező töltésű részecskéről van szó.
Ezt a részecskét a későbbiek során \textit{pozitronnak} ($e^+$) nevezték el. Ez a felfedezés fényes igazolását jelentette az antirészecskék létezésének, melyeket a
Dirac-féle relativisztikus kvantumelmélet jósolt meg néhány évvel Anderson felfedezése előtt.

\subsubsection*{Müon felfedezése}
1938- ban S. H. Neddermeyer és C. D. Anderson Wilson-kamrával kozmikus
sugárzás méréseket végeztek. A mérési technikát annyiban fej lesztették tovább, hogy a kamra fölött és a kamra alatt egy-egy Geiger-Müller-számlálót
helyeztek el, amelyeket koincidenciába kapcsoltak, a kamra belsejébe pedig abszorbenst tettek. Ezzel lényegesen megnövelték az olyan kozmikus részecskék megtalálásának a lehetőségét, ame lyek nagy áthatolóképességűek
(kemény, áthatoló komponens). Méréseik során egy olyan felvételt találtak, amely az ábrán látható.
\begin{figure}[H]
    \centering
    \includegraphics[width=0.3\textwidth]{imgs/13-tetel/muon_felfedezese.png}
    \caption{
        Az ábrán egy müonnyomot láthatunk, amely a Wilson-kamra
        gáztérfogatában lefékeződött és az ábrán látható módon e/bomlott egy elektronra,
        és két nyomot nem hagyó részecskékre
    }
\end{figure}
A Wilson-kamrát 0.79 T indukciójú mágneses térben helyezték el, így
mérhetővé vált a tömeg. Azt. találták, hogy az észlelt részecske tömege
mintegy 240-szerese az elektron tömegének, világos, hogy a részecske nem
lehet sem elektron, sem proton. Az analízis alapján kijelenthették, hogy
egy eddig még nem ismert részecskét találtak, amely se nem elektron,
se nem proton. Ez a részecske pozitív töltésű, tömege nagyobb az elektron tömegénél, de kisebb a protonénál, kb. 200-szorosa az elektron tömegének. Ezt a részecskét először mezotronnak, illetve $\mu$-mezonnak nevezték el.
Később azonban, amikor a tulajdonságait jobban megismerték, kiderült,
hogy nem tartozik bele az ún. mezonok családjába, emiatt áttértek a ma is használatos müon (jele: $\mu$) elnevezésre.

Hamarosan más kutatók is találtak hasonló részecskéket a kozmikus
sugárzásban, és ezzel egyértelművé vált a müonok felfedezése. A müont kezdetben érthető módon a Yukawa-féle ún. közbenső mezonnal azonosították.
Mint említettük (s még részletezni fogjuk), a későbbiekben kiderült, hogy ez a feltételezés tévedés volt. A felvételen látható, hogy a müon, miután
lefékeződött, elbomlott, és a bomlásnál egy elektron keletkezett. 

\subsubsection*{V-részecskék felfedezése}
1947-ben G. D. Rochester és C. C. Butler manchesteri fizikusok Wilson-kamrával végeztek kozmikus sugárzási méréseket. Vizsgálták a kozmikus
sugárzás ún. áthatoló komponensét, amely a kozmikus sugárzás ólmon való áthaladása után megmarad. Ekkor 5000 felvételből (ez kb. 1500 óra Wilson- kamra működésnek felel meg) 2(!) olyat találtak, amelyen villaszerű nyomok
voltak láthatók. A kamrába 3 cm-es ólomlemezt helyezve, a villa alak tí. részecskék száma megnövekedett. Ha ezek a részecskék - melyeket egyébként a
nyom alakjáról V-részecskéknek neveztek el - ütközésből származtak volna
(ami nagyon természetes feltevésnek látszott), akkor néhány 100-szor nagyobb valószínűséggel kellett volna keletkezzenek ólomban, mint gázban.

A kapott eredmény tehát kizárja azt, hogy ütközési folyamatról legyen szó.
Az első ilyen felvételek egyikét láthatjuk az 18.4. ábrán. Felmerült az a
feltevés, hogy V-részecskék valamilyen semleges, tehát a Wilson-kamrában
nyomot nem hagyó részecskének a bomlásából származnak. E feltételezés
alapján a bomló láthatatlan részecske tömegére 1000 elektrontömeget, közepes élettartamára pedig 5 · 10-8 s értéket kaptak. A fentiek alapján a
felvételt úgy magyarázták, hogy egy semleges részecske bomlásáról van szó,
melynek a tömege 770 és 1600 elektrontömeg közé esik. Megemlítjük, hogy
ezekben a vizsgálatokban részt vett egy magyar t udós, Jánossy Lajos is.
Később a V-részecskéket a semleges kaonnal ($K^0$) és $\Lambda$-hiperonnal ($\Lambda^0$) azonosították . A megfelelő bomlási folyamatok:
\begin{equation}
    K^0 \quad \rightarrow \quad \pi^+ + \pi^- \qquad \Lambda^0 \quad \rightarrow \quad p + \pi^-
\end{equation}
\begin{figure}[H]
    \centering
    \includegraphics[width=0.4\textwidth]{imgs/13-tetel/v_reszecske.png}
    \caption{V-részecske felvétele Wilson-kamrában.}
\end{figure}
A részecskefizikában a bomlási folyamatokat kétféleképpen szokás jelölni: vagy + jelet tesznek a bomlási termékek közé , vagy csak egyszerűen
egymás mellé írják öket. A fenti példánkban a két lehetséges jelölési mód:
\begin{equation}
    K^0 \rightarrow \pi^+ + \pi^-, \quad \Lambda^0 \rightarrow p + \pi^- \quad \text{vagy} \quad K^0 \rightarrow \pi^+ \pi^-, \quad \Lambda^0 \rightarrow p \pi^-
\end{equation}

\subsubsection*{A $\pi$-mezon felfedezése}
C. F. Powell és munkatársai br istoli kutatók fotoemulziós méréseket végeztek 1947-ben Bolíviában az Andokban 5500 m, és a P ic du Midi hegycsú-
cson (Pireneusok) 2800 m magasságban. Az emulzió kiértékelésénél olyan
esetekre bukkantak, amelyekben a lefékeződő primer részecskéből egy szekunder, másodiagos részecske is keletkezett. Összesen 644 olyan mezonnyo-
mot figyeltek meg, amely az emulzióban végződött , és így teljes mértékben
kiértékelhe tő volt.

A másodiagos részecskék kilépésének az iránya teljesen véletlenszerű
volt, a tömegük pedig azonos. Úgy tűnik, hogy a kozmikus sugárzásban lévő
ismeretlen részecskék nem hatoltak be az atomokba, hanem elbomlottak.
A két nyom közel volt egymáshoz, ezért a fotoemulzió előhívása sor án fellépő
torzító effektusok mind a kettőre egyform án hatottak. Megszámlálták a
fotoemulzióban az 1 cm-re eső szemcséknek a számát, és ebből következtetni
tudtak a részecskék sebességére. A hatótávolságból és a sebességből pedig ki
lehetett számítani a tömeget. Azt találták, hogy a pr imer részecske tömege
majdnem másfélszer nagyobb, mi nt a szekunderé. Az észlelésekből és az ezt
követő interpretálásból egyértelműen kiderült, hogy itt egy újfajta, eddig
ismeretlen részecskének a bomlási folyamatáról van szó. Ezt a mezont $\pi$-
mezonnak ($\pi$) nevezték el (újabb nevén egyszerűen csak pion), és később az
erős kölcsönhatást közvetítő Yukawa-féle mezonnal azonosították.
A $\pi$-mezon elbomlott egy kisebb tömegű részecskére, amelyet azonosítani tudtak a már ismert müonnal. A kozmikus sugárzásban egyaránt találtak pozitív ($\pi^+$) és negatív ($\pi^-$) pionokat. Ma már tudjuk, hogy a teljes
bomlási folyamat: $\pi^+ \rightarrow \mu^+ + \nu_\mu$, azaz müon-neutrínó is keletkezik. A müon tovább bomlik: $\mu^+ \rightarrow e^+ + \nu_e + \bar{\nu}_\mu$.

\subsubsection{Részecskék osztályozása}
Közben 1930-ban feltételezték egy új részecskének, a tömeg nélküli, semle-
ges, csak gyengén ható részecskének, a neutrínónak a létezését. (Feltétele-
zésére - mint azt a 8. fejezetben láttuk , és a 20. fejezetben látni fogjuk -
azért volt szükség, hogy meg tudják magyarázni a $\beta$-bomlásnál az elektron
energiaspektrumának folytonos voltát. A neutrínó kísérleti kimutatására
csak az ötvenes évek közepén került sor.)

Ezek alapján tehát az 50-es évekig csak néhány elemi részecske volt
ismert. A fejlődés azonban felgyorsult, egyre újabb és újabb részecskéket
fedeztek fel, és ma már a részecskék száma mintegy 300-ra(!) tehető. Ez
túl sok, ha azt várjuk tőlük, hogy ezekkel a részecskékkel építsük fel a
világot. Érthető módon új elképzelések szü lett ek, és folynak ma is a kísérleti
vizsgálatok, hogy nincs-e ennél egy mélyebb szint, amin jóval kevesebb
számú elemi részec s kéből lehet fe lépíteni a világot. Úgy néz ki , hogy van
ilyen szint. (Ezzel a kérdéssel e fejezet utolsó pontjában fogunk foglalkozni.)

Az egyes részecskék egymással kölcsönhatásba lépnele A részecskék
egyik legfonto sa bb jellemző tu lajdonsága, hogy milyen kölcsönhatásban vagy
kölcsönhatásokban tudnak részt venni. Éppen ezért a részecskék rendszerezésének, csoportokba foglalásának alapját a kölcsönhatási típusok képezik.
A következő pontban a részecskék osztályozásával ismerkedünk meg.

\subsubsection*{Stabil részecskék}
A természetben a közelmúltban négy alapvető kölcsönhatástípust ismertünk, ezek száma az utóbbi időben- kettejük közötti "rokonság"
felismerése révén - háromra csökkent. A táblázatban még a 4 kölcsönhatás néhány jellemző tulajdonságát foglaljuk össze. A táblázatban a
kölcsönhatások csökkenő erősség szerint szerepelnek.
\begin{figure}[H]
    \centering
    \includegraphics[width=0.8\textwidth]{imgs/13-tetel/kolcsonhatasok.png}
    \caption{Az alapvető kölcsönhatások jellemzői.}
\end{figure}
Lehetséges, hogy a gravitációnál nem is beszélhetünk a hagyományos
értelemben vett kölcsönhatásról, hanem csupán a térnek és időnek valamilyen geometriai jellegű tulajdonságáróL Ez esetben természetesen közvetítő
részecske sem létezhet. A fentiek ellenére a táblázat nagyságrendileg helyes képet ad a
különböző kölcsönhatások erősségéről , amelyek - mint látjuk - sok nagyságrendet fognak át. Hasonlóképpen sok nagyságrendet fognak át az átlagos
kölcsönhatási idők és az átlagos hatótávolságok, kezdve a rendkívül gyorstól a nagyon hosszú ideig tartókig, a nukleáris távolságoktól az Univerzum
méretével összemérhctőkig. A kölcsönhatásokat mai elképzeléscink szerint
részecskék közvetítik; ezeket tüntettük fel a harmadik oszlopban.

A ma ismert részecskék közül a legfontosabbakat a következő táblázatban
foglaltuk össze, a teljességre való törekvés nélkül. A részecskéket három alapvető csoportra osztjuk. Az első csoportban teljcsen különállóan, magányosan
foglal helyet az atomfizikából már jól ismert \textit{foton} ($\gamma$) ($\varphi \tilde{\omega} \varsigma$, $\varphi \omega \tau \text{ó} \varsigma$  = fény), ,
a következő nagyobb csoport a leptonok családja, a harmadik a hadronoké.
A Ieptonok a gyenge kölcsönhatás ban vesznek részt (ha elektromos töltéslik
van, akkor az elektromágnesesben is). A hadronok legfontosabb közös vonása az, hogy erős ($\acute{\alpha} \delta \varrho \text{ó} \varsigma$ = erős) kölcsönhatásban vesznek részt (ez nem
zárja ki, hogy az erős mellett részt vegyenek a gyenge és az elektromágneses kölcsönhatásban is). A leptonok tömege általában kisebb, mint a hadronoké. (Történelmileg innen nyerték elnevezésüket a görög $\lambda \varepsilon \pi \tau o \varsigma$ = könnyű szóból.)
\begin{figure}[H]
    \centering
    \includegraphics[width=0.8\textwidth]{imgs/13-tetel/stabil_reszecskek.png}
    \caption{"Stabil" részecskék}
\end{figure}

A táblázat el ső oszlopában az előbb em lí tett családok elnevezése található. A következő oszlopban a csoportokon belüli alcsoportok elnevezését
tüntettük fel, ezután következik a. részecske szimbolikus jele. A következő a nyugalmi energiát ($E_o = mc^2$ , MeV-ban), majd a spineket tartalmazó
oszlop ($s$). Az ezután következő számsor a részecskék bomlásának közepes
élettartamát ($T$) tartalmazza s-ban. Itt meg kell jegyeznünk, hogy a táblázat a "stabil" részecskék elnevezést viseli, holott a részecskék köztil stabilnak
csak kivételesen nevezh ető néhány. Különbséget szokás tenni azonban a viszonylagosan stabilis részecskék, amelyek közepes élettartama nem túl rövid
($\sim 10^{-23} s$) és az ennél rövidebb élettartamúak között. Ez utóbbiakat rezo-
nanciáknak is szokás nevezni, tipikus adataikat a mezon- ($\mu  \varepsilon \sigma \text{o} \varsigma $ = középső),
a barion- ($\beta \alpha \rho \iota \varsigma$ = "nehéz"), illetve a hiperontáblázat ($\upsilon \pi \varepsilon \varrho$ = felett , rendkívül) végén tüntettük fel.

A következő négy oszlop külön magyarázatot igényel.

a) Ritkaság (strangeness, S). Vannak részecskék, amelyek atommag-
kölcsönhatásban - a többi részecskékhez képest - aránytalanul ritkán jelennek meg, ezért ritka részecskéknek nevezzük őket. A ritka részecskék erős
kölcsönhatásban mindig csak párosával keletkeznek. Az ilyen páros keletkezésnél valamilyen "töltés" megmaradására gyanakodhatunk. (Itt természetesen a töltés szót általános értelemben kell érteni, és nem csak az elektromos
töltésre gondolunk.) Mivel erős kölcsönhatással könnyen előállíthatóak , és
ugyanakkor olyan részekre bomlanak, amelyek részt vesznek erős kölcsönhatásban, azt várták, hogy az erős kölcsönhatásokra jellemző gyorsasággal
bomlanak. Ez azonban nem így van. Ennek a jelenségnek az értelmezésére
kellett bevezetni egy új megmaradási törvényt, ami kb. 10-13 nagyságrenddel lassítja a bomlást. Ezt az új megmaradási tételt a ritkaság megmaradásának nevezik. A $K^+$-hoz és $K^0$-hoz az $S = +1$ ritkaságot rendelték, a
"nem ritka" részek ritkasága értelemszerűen $0$. Vannak $S = 2, 3, ...$ ritkaságú és negatív ritkaságú részecskék is. A ritkaság megmaradásának törvénye kimondja, hogy zárt rendszerben a ritkaság~kvantumszámok - vagy
ritkaságtöltések - összegének állandónak kell maradnia. E törvényt pontosan kielégíti az erős és az elektromágneses kölcsönhatás , a gyenge azonban
nem. A részecskepárok keletkezésekor (erős kölcsönhatás) mindig keletkezik egy másik ritka részecske is, amelynek ritkasági száma helyrehozza a
ritkaság megmaradásának törvényét. Pl.:
\begin{equation}
    p + \pi^- \rightarrow p + \Lambda^0 + K^0
\end{equation}
A ritkaságuk pedig:
\begin{equation}
    S = 0 + 0 = (-1) + (+1)
\end{equation}
Mivel a gyenge kölcsönhatásban a ritkaság megmaradási tétel sérül, ritka
részecske gyenge kölcsönhatással zérus ritkaságú részecskére is bomolhat.
Emiatt a ritka részek sokkal hosszabb ideig élnek, mint a nem ritkák (ahol
$S = 0$, tehát a megmaradási tétel nem sérül, vagyis a bomlás erős kölcsönhatás révén, azaz nagyon gyorsan is végbemehet). A ritka részek gyenge
bomlására példa a kaon ($K$) egyik lehetséges bomlása:
\begin{equation}
    K^+ \rightarrow \pi^+ + \pi^+ + \pi^-
\end{equation}
Mivel a $K^+$-mezon ritkasága +1, a keletkező pionoké pedig nulla, nyilvánvalóan sérül a ritkaság megmaradásának törvénye. Ezért lassúak ezek a folyamatok.

b) Leptonszám ($L$). Megfigyelték, hogy valahányszor leptonok keletkeznek, mindig párosával jelentkeznek. Ez - úgy, mint ritka részecskéknél - meg-
maradási tételre utal, amit a leptonszám megmaradási tételének hívnak.
Tulajdonképpen három független leptonmegmaradási törvény van: egy az
elektronokra, egy a müonokra és egy a tau-részecskékre. Az e- és neutrínója, $\nu_e$, a $+1$ (elektron)leptonszámot, az antirészskéi pedig a $-1$ (elektron)leptonszámot kapták. A tapasztalat szerint bármilyen zárt rendszerben
az elektron-leptonszámok összeg állandó marad. PL a neutron bomlása:
\begin{equation}
    n \rightarrow p + e^- + \bar{\nu}_e
\end{equation}
Ahol a leptonszámok:
\begin{equation}
    0 = 0 + (+1) + (-1)
\end{equation}
Mivel a neutron leptonszáma nulla, a bomlástermékek leptonszámainak
összege is zérus kell legyen. Láthatjuk , hogy ez valóban teljesül. A leptonszámok megmaradásának törvényét Marx György és J. B. Zeldavies ismerte fel.

c) Barionszám ($B$). A következő oszlopban a barionszámot tüntettük
fel. A leptonoknál megisme rt mintára minden barionnak van egy $+1$ barionszáma, az antibarionoknak pedig $-1$. Zárt rendszerben a barionszámok
összegének állandónak kell maradnia (Wigner Jenő). A nem barionrészecs-
kék (tehát pl. a foton, leptonok, a mezonok és más közvetítő részecskék)
barionszáma természetesen $0$.

d) Izospin ($T$). A magerők szempontjából a proton és a neutron teljesen
egyformán viselkednek Ami kis különbség van köztük, az egyszerűen az
elektromos töltéstik különbözőségének számlájára írható. A két részecskét
egyetlen részecske, a nukleon két állapotának tekinthetjük. Világunk tehát
- legalábbis az erős kölcsönhatás szempontjából - szimmetrikus a proton és
a neutron vonatkozásában. Ennek a ténynek a kifejezésére- W. Heisenberg
javaslatára bevezettek egy új fogalmat, az izospint, amelyet $T$ betűvel
jelölünk. (Az elnevezésnek nincs köze az izotópokhoz, csupán a "rendes"
spintől való megkülönböztetést szolgálja).

A $T$ kvantumszám a részecske ízo spinjének abszolút értékét (típusát) jelöli, és a közönséges spinhez hasonlóan 0, 1/2, 1 , 3/2, ... stb. értéket vehet fel.

A $T$ mellett még egy másik mennyiség (kvantumszám) is használatos, az izospin harmadik komponense, amelyet $T_3$-mal jelölünk. Ennek különböző
értékével jelöljük azokat az állapotokat, amelyeket a szimmetrikus világ egy részeeskéje elfoglalhat. Pl. a nukleonnak két állapota van: a proton és a neutron.

A nukleonok esetében, ahol $T = 1/2$ , a protonra $T_3 = 1/2$ , a neutronra $T_3 = -1/2$ jelölést használunk. A $T_3$ mennyiség tehát a nukleonok
családjába tartozó egyes családtagok jellemzésére szolgál. Az egész család a $T$ értékével határozható meg.

Az előbbieket úgy is mondhatjuk, hogy a nukleonok clublettet alkotnak. Vannak tripiettek is, tehát olyan családok, amelyek három tagból állnak. Ilyen pl. a pion. A pioncsaládot a $T = 1$ izospin abszolút érték jellemzi, $T_3$ pedig háromféle értéket vehet fel: $-1 , O, +1$. Általában egy $T$
izospinű részecskecsalád 2T + l tagból állhat, T3 lehetséges értékei pedig:
$- T, -T+ l, ... , T - l , T$.

A hadronok jellemző kvantumszámai között fennáll egy általános összefüggés:
\begin{equation}
    Q = T_3 + \frac{B + S}{2}
\end{equation}
Itt $Q$ az elektromos töltés, $T_3$ az izospin harmadik komponense, $B$ a bari-
onszám, $S$ a ritkaság. A ritkaság és a barionszám összegét hipertöltésnek
nevezik:
\begin{equation}
    Y = B + S
\end{equation}
A fenti összefüggés segítségével, ha egy kivételével ismerjük a töltéseket, akkor a hiányzót kiszámíthatjuk.

\subsubsection*{Antirészecskék}
Az antirészecs ke mindenben azonos az eredetivel (névadó), csak töltésben különbözik tőle. A töltést itt tágabb értelemben kell felfognunk: nemcsak
elektromos, hanem barion-, leptontöltésről, helicitásról is szó van. Dirac relativisztikus kvantumelméleténekfényes bizonyítá-
sát jelentette, amikor Wilson-kamrában sikerült felfedezni a pozitront, amely
- mint láttuk - az elektronnal teljesen azonos, azonban a töltése nem negatív, hanem pozitív. Célszerűnek látszott e részecskének a Dirac-elméletben
szereplő "lyukkal" való azonosítása.

Később felfedezték a proton antirészét, az antiprotont, amely mindenben megegyezik a protonnal, csak negatív töltésű , és barionszáma is ellentettje a protonénak. Ma már kísérletileg ismerjük többek
között az antineutront, sőt még összetett képződmények , atommagok "antirészecskéit" (pl. antihéliummagot) is sikerült előállítani. Nem kétséges,
hogy valamennyi részecske rendelkezik antirészecskéveL Ennek megfelelően
a korábban látott részecskék száma a valóságban (közel) kétszer
annyi, mert mindegyikhez rendelhető egy antirészecske. Van olyan eset, amikor a részecske és az antirészecskéje ugyanaz. Ilyen pl. a foton és a $\pi^0$. Egy
adott típusú neutrínó (mondjuk elektronneutrínó) antirészecskéje abban különbözik tőle , hogy a helicitása és a leptontöltése ellentétes.

Antiatomok is elképzelhetőek, ahol a protonok és neutronok helyett antiprotonok és antineutronok vannak, körülöttük elektronok helyett pozitronok keringenek. Kísérletileg azonban még nem mutatták ki ezeket. Érdekes
kérdés, hogy a természetben - legalábbis az Univerzum általunk eddig ismert részében - miért találkozunk elsősorban részecskékkel, és miért nulla
(vagy rendkívül kevés) az antirészecskék száma . Vannak elméletek, amelyek
választ próbálnak adni erre. Minden részecske, ha találkozik az antirészecskéjével, azonnal "megsemmisül" (annihilálódik), és a teljes energia- beleértve a nyugalmi tömeget is - sugárzási energiává alakul át.

\subsubsection*{Megmaradási törvények}
Mai tudásunk szerint a világon a tér és idő bizonyos szimmetriatulajdonságokkal rendelkezik. Így pl. a tér homogén és izotróp, az idő homogén.
Ugyanakkor ismerünk a természetben alapvető megmaradási törvényeket.
A szimmetriatörvények és megmaradási törvények egymástól függetlenül fogalmazódtak meg a fizi,ka története során. A későbbiekben az emberi gondolkodás nagy előrelépését jel entette annak feli smerése,
hogy a szimmetriák és a megmaradási törvények között egyé rtelmű és szoros kapcsolat van.
Aszimmetria tulajdonságokból egy-egy alapvető fizikai mennyiség megmaradásának a törvénye követ kezik (Emmy Noether, Göttingen, 1918). Sikerült kimutatni, hogy
\begin{itemize}
    \item a tér homogenitásából következik az impulzus megmaradás törvénye,
    \item az impulzusmomentum-megmaradás törvényének a mélyén a tér izotróp volta rejtőzik,
    \item az idő homogenitása az oka az energiamegmaradás törvényének.
\end{itemize}

Így tulajdonképpen minden szimmetriatulajdonsághoz egy-egy megmaradási tétel tartozik . Az előzőekben ismertetett ún. folytonos szimmetriák
mellett vannak ún. diszkrét szimmetriák is, amilyen a tükrözési szimmetria, amelynek teljesülése a paritás megmaradásához vezet. Azonban kiderült,
hogy a gyenge kölcsönhatások meghökkentő módon nem tükörszimmetrikusak. Erre először a kaon bomlásakor bukkantak rá.
Az ötvenes évek vége felé két érdekes mezoncsoportot találtak. Az egyiket $\tau$-mezonnak (nem azonos az 1975-ben felfedezett nehéz leptonnal, a $\tau$-val), a másikat $\Theta$-mezonnak nevezték el.
A kétféle mezon mindegyikének többféle bomlási típusa van, a bomlási típusok azonban egymástól eltérnek. Így pl.:
\begin{equation}
    \tau^+ \rightarrow \pi^+ \pi^+ \pi^- \qquad \Theta^+ \rightarrow \pi^+ \pi^0
\end{equation}
stb. A különböző bomlási folyamatok vizsgálata során meglepetésként hatott az a felismerés, hogy bármire is bomlik el az egyik vagy a másik részecske, tömege és az átlagos élet tartama mind a kettőnek ugyanaz. Sőt
az egyes bomlás folyamatok termékének, azok szög- és sebességeloszlásának
analíziséből az impulzusmomentum-megmaradás figyelembevételével az is
kiderült, hogy mind a két mezon (a $\tau$ és a $\Theta$) spinje nulla. A két mezon azonos egymással, a különböző bomlások csak "radioaktív elágazások", vagy
valami másról van szó? A $\tau$ és a $\Theta$-mezon azonosítása ellen komoly érv szól.
A $\tau$-mezon csak két pioura bomlik. A páratlan számú piont leíró hullámfüggvény a tükörállapot hullámfüggvényétől előjelében különbözik. A hullámfüggvény paritásának minden tükörszimmetrikus kölcsönhatásnál meg
kell maradnia, ezért a $\tau$-mezon csak +1 paritású, a $\Theta$-mezon csak $-1$ paritású lehet, hiszen belőlük ilyen állapotok keletkeznek bomlás révén. A két
mezon nem lehet azonos. A kísérletek során azonban kiderült, hogy a két
típusú mezon között sem tömegben, sem töltésben, sem spinben, sem élettartamban, sem magkölcsönhatásban nincs különbség. Úgy tűnik tehát , hogy
a két mezon azonos, amely a g+ nevet kapta. Ekkor merült fel T. D. Lee
és C. N. Yang merész hipotézise, hogy hátha a két mezon valójában teljesen
azonos, azonban a természetben - legalábbis a gyenge kölcsönhatásoknál -
a paritás megmaradásának törvénye sérül. Ezért az elméletüké rt 1957-ben
Nobel-díjat kaptak.

Később direkt kísérletekkel sikerült kimutatni a tükörszimmetria-sérülést gyenge kölcsönhatásoknál, és egy sokkal kisebb mértékű sérülést az
elektrogyenge kölcsönhatásnál.

A következő táblázatban összefoglaltunk néhány megmaradási törvényt.
Kérdés, hogy ezek valóban ugyanolyan ál talánosak-e, mint pl. az elektromos
töltés megmaradása. A válasz az, hogy nem teljesen, ti. a gyenge kölcsönhatás pl. sérti a ritkaság megmaradásával kapcsolatos szimmetriát. Az összes
ritka részecske (a $\Xi^0$ kívül) gyenge kölcsönh atás útján bomlik el
úgy. hogy ennek során a ritkaság értéke megváltozik Előbbi összefüggéseinkből következik, hogy ennek következményeként a $T$ és $T 3$ sem megmaradó
mennyiség. A gyenge kölcsönhatás tehát igencsak "engedetlen " kölcsönhatás, egy sor szimmetriával szembeszegül.
\begin{figure}[H]
    \centering
    \includegraphics[width=0.6\textwidth]{imgs/13-tetel/megm_torvenyek.png}
    \caption{Megmaradási törvények érvényessége}
\end{figure}

\subsubsection{Leptonok}
A Ieptonok tu lajdonképpen két csoportra osztható k: az elektronok és a neutrínők csoportjára . Mind a két csoportban három- három komponenst ismerünk. Az elektronok esetében maga az elektron ($e$), ezután a "nehéz elekt-
ron": a müon ($\mu$), és az 1975-ben felfedezett nagy tömegű tau-részecske (semmi köze sincs a paritás megnemmaradásánál említett $\tau$- illetve $\Theta$-
mezonhoz), vagy más néven tauon ($\tau$). Mindegyik elektron típushoz tartoz ik egy-egy neutrínó: az elektronhoz elektron-neutrínó ($\nu_e$),
a müonhoz a müonneutrínó ($\nu_\mu$) és a tau-részecskéhez a tau-neutrínó ($\nu_\tau$).
(A tau-neutrínó létezésében nem kételkedünk, bár közvetlenül kísérletileg
még nem sikerült kimutatni.) Nagyon szemléletes képet kapunk , ha a leptonokat dublettekben rendezzük el:
\begin{equation}
    \begin{pmatrix}
        e \\ \nu
    \end{pmatrix} , \quad
    \begin{pmatrix}
        \mu \\ \nu_\mu
    \end{pmatrix} , \quad
    \begin{pmatrix}
        \tau \\ \nu_\tau
    \end{pmatrix}
\end{equation}
Az elektronnal már alaposan megismerkedtünk a héjfizikában. A müon - a tömegét kivéve - minden tulajdonságában megegyezik az elektronnal
(A müon mintegy kétszázszor nehezebb az elektronnál.)

\subsubsection{Hadronok}
A hadronokat két nagy csoportra osztjuk: a mezonokra és a barionokra.
(A mezon és a barion itt ismertetett szétválasztása történelmi) A szétválasztás alapja a tömeg és a spin. A barionok spinje feles, amezonoké egész.
A mezonok tömege az elnevezés kialakulásának idején a Ieptonok és a barionok között foglalt helyet. (Az újabban felfedezett mezonokra már nem áll
fenn , hogy tömegük a Ieptonok és a bar ionok közé esik.) A mezonok családja az utóbbi időben két új (nagy tömegű) taggal gyarapodott $J/\psi$ (ejtsd:
\textit{dzsi-pszi}) és az $\Upsilon$ (ejtsd: \textit{üpszilon}).
A barionok közül már jól ismerjük a nukleonokat: a protont és a neutront. A többiek az ún. hiperonok családját alkotják.

\subsubsection{Kvarkok}
1911-ben E. Rutherford a-részecskékkel bombázott anyagokat, és a szórás iképből arra következtetett, hogy az atomnak szerkezete van: központi kemény magból és az ezt körülvevő elektronokból áll. Filozófiájában és koncepciójában hasonló kísérleteket végeztek R. Hofstadter és munkatársai a 60-as
évek végén Stanfordban (USA), a lineáris elektrongyorsítón, ahol 15 GeV-os elektronokkal bombáztak nukleonokat. A rendkívül nagy energiájú elektronok egészen közel tudtak férkőzni a protonokhoz és a neutronokhoz, és
mintegy "letapogatták" belsejüket. A kísérletek eredményeképpen kialakult
szóráskép úgy magyarázható, hogy a nukleonok nem "elemi" részek, hanem szerkezettel rendelkeznek: kis kemény magok vannak bennük, amiket
partonoknak (part of proton) neveztek el. Később kiderült, hogy a partonok azonosak az elmélet által már régebben megjósolt ún. \textit{kvarkokkal}.
Később hasonló eredményre jutottak nagyenergiájú müonok szóródása során a CERN-ben. (A kvark elnevezés M. Gell-Mann nevéhez fűződik. Aszó
val ószínűleg JJames Joyce: "FFinnegan's Wake" c. regényének egy részletéből
származik, és a könyvben tulajdonképpen nincs meghatározott jelentése.
E rész ben az alkoholista szereplő - felébredve delíriumos álmából - újabb
italt kér ilyenképpen: "Three quarks for muster Mark". Németül egyébként Quark túrót jelent.)

\subsubsection*{Kvark fajták (ízek)}
Gell-Mann és Zweig háromféle kvark létezését tételezte fel, ame lyeket az
up (= fel), down (= le) és strange (= ritka, furcsa, idegen) angol szavak
kezdőbetűiből $u$, $d$, $s$ betűkkelje löltek. Tulajdonságaik meglepőek: nem
egész számú elemi töltéssel, hanem annak tört részével rendelkeznek. Meg
kell azonban jegyezni, hogy az egységnyi töltés létezését senunilyen természeti törvény nem írja elő, megállapodás, aminek a kísér leti tapaszta lat eddig
nem mondott ellen. Nem sérti meg egyetlen eddigi elvünket sem, ha léteznek
tört elektromos töltések. Az antikvarkok minden töltése abszolút értékben
ugyanaz, csak ellentett előjelű . Úgy néz ki, hogy az $u$- és $d$-kvarkból, továbbá
az elektronból felépíthető közvetlen környezetünk. Pl. a proton kvarkszerkezete:
\begin{equation}
    p = uud
\end{equation}
\begin{figure}[H]
    \centering
    \includegraphics[width=0.3\textwidth]{imgs/13-tetel/kvarkok.png}
    \caption{A proton kvarkszerkezete}
\end{figure}
Hasonlóan a neutroné:
\begin{equation}
    n = udd
\end{equation}
A valóságban a helyzet bonyolultabb , és a hadronok (köztük a proton és neutron) kvarkszerkezete sokkal komplikáltabb. Leptonszórásokból ui. tudjuk, hogy pl. a nukleonok az említett
három, ún. "vegyérték''-kvarkon (és az őket összetartó ún. gluonokon) kívül még kvark- antikvark párokból álló "kvarktengerrel" is rendelkeznek.
Pl. a pozitív pion ($\pi^+$) kvarkszerkezete:
\begin{equation}
    \pi^+ = u \bar{d}
\end{equation}
A mezonok kvarkszerkezete tehát egy kvarkból és egy antikvarkból áll.
A ritka részecskék felép ítésében szerepet játszik a $s$-kvark is. Pl. a $K^+$-mezon kvarkszerkezete:
\begin{equation}
    K^+ = u \bar{s}
\end{equation}
A $J/\psi$-részecske úgy magyarázható, mint egy újfajta (az $u$, $d$ és $s$ kvarkoktól különböző) kvarknak és az antirészecskéjének a kötött állapota. Az új kvarkot charmnak (bájos, c) nevezték el,
tehát a kvarkok családja új taggal bővült. Ezek szerint a $J/\psi$-részecske kvarkszerkezete
\begin{equation}
    J/\psi = c \bar{c}
\end{equation}
A később felfedezett $\Upsilon$-részecske egy még nehezebb kvarknak és az antirészecskéjének a kötött állapota. Ezt a kvarkot bottomnak (alsó, b) nevezték el, így a kvarkok családja ismét új taggal bővült. Az $\Upsilon$-részecske kvarkszerkezete tehát
\begin{equation}
    \Upsilon = b \bar{b}
\end{equation}
Így a kvarkoknak jelenleg hat ismert fajtája (íze) van: up ($u$), down ($d$), strange ($s$), charm ($c$), bottom ($b$) és top ($t$). A top kvarkot 1995-ben fedezték fel a Fermilabban (USA). A kvarkok tömege
nagyon különböző: az $u$- és $d$-kvark tömege néhány MeV/c², míg a $t$-kvark tömege mintegy 175 GeV/c².

A kvarkokat érdemes párokba rendezni:
\begin{equation}
    \begin{pmatrix}
        u \\ d
    \end{pmatrix} , \quad
    \begin{pmatrix}
        c \\ s
    \end{pmatrix} , \quad \quad
    \begin{pmatrix}
        t \\ b
    \end{pmatrix}
\end{equation}
Hasonlóan a Mengyelejev-rendszerhez, itt is sejthető volt a szimmetria
miatt még egy kvarknak, a $t$-kvarknak (top = felső vagy truth = igazság)
a létezése. Ezt csak a közelmúltban sikerült kísérletileg kimutatni. Érdemes
megjegyezni, hogy úgy látszik, mélyen fekvő szimmetria áll fenn a kvarkok
és leptonok között. A leptonokat is - mint láttuk - párokba
(doublettekbe) rendezhetjük el, a fenti sémához hasonlóan.

\subsubsection*{Szín}
A kvarkoknak azt a tulajdonságát, hogy $u , d , s , c, t$ vagy $b$ típusúak-e,
"iznek" nevezzük. A későbbiek során kiderül, hogy a helyzet bonyolultabb,
mert a kísérleti tényekkel csak akkor tudunk összhangba jutni, ha feltételezzük,
hogy a kvarkok mindegyikének az íze mellett van még egy tulajdonsága,
azaz van még egy kvantumszáma is. Ezt utóbbit "színnek" nevezték el. Természetesen mindkét szó játékos fantáziára utal, semmiképpen nincs köze a
gyakorlati értelemben vett ízhez vagy színhez. (Az antikvarkok színe "ellen-
tétes".) Az egyszerűség kedvéért arról beszélünk, hogy a kvarkok színe lehet
piros, zöld és kék, az antikvarkoké "antipiros", "antizöld" és "antikék".

A ténylegesen megjelenő részecskék - a hadronok - mindegyike a tapasztalat szerint "színtelen ": azaz a kvarkokl1ak és antikvarkoknak úgy kell kombinálódniuk,
hogy színeik eredőként eltűnjenek. A kvarkok kétféle módon adhatnak színtelen (fehér) eredőt, azaz hadront. (A színtelen az optikában
megszakott szóhasználathoz képest nem azonos a fehérrel, a részecskefizikában azonban e két jelzőt felváltva, ugyanabban az értelemben használják.)
Az egyik lehetőség, hogy három kvark kapcsolódik össze, mindegyik kvarknak más és más színe van, és a háromszín (éppen úgy, mint az optikában)
eredőként fehéret (színtelent) ad: ezek a barionok. A másik lehetőség, hogy
egy adott színű kva.rk az antikvarkjával kapcsolódik össze, amikor is a szín
és az "antiszín" ugyan csak kompenzálja "kifehéríti" egymást, és eredőként
az objektumnak nincs színe: ily módon épülnek fel kvarkokból a mezonok.

\subsubsection*{Glüonok}
A többi kölcsönhatás mintájára fe ltételezhetjük, hogy az erős kölcsönhatást
kvarkok között valami részecske közvetíti, éppúgy, ahogy az elektromágneses kölcsönhatást a foton és a gyenge kölcsönhatást a megfelelő vektorbozonok.
Valóban, a kvantum-színdinamika (kvantum-kromodinamika, QCD) elmélete szerint a közvetitök tömeg nélküli részecskék, amelyeket "ragasztóknak" , gluonoknak neveztek el. A gluonok szintén rendelkeznek színnel,
éppúgy, mint a kvarkok. Összesen 8 féle gl uon van. Sajátos vonás, hogy színük következtében a gluonok is alkothatnak kötöttállapotokat, ezeket gluonlabdáknak nevezték.
Létüket ez ideig kísérletileg még nem sikerült megerősíteni .

Az erős kölcsönhatás tehát alapjait tekintve színkölcsönhatás, amelyet
színes kvarkok között színes gluonok közvetítenek. (A kvarkok közti kölcsönhatás anaJógiába állítható az elektromágneses kölcsönhatás Coulomb-törvényével. A színtelen hadronqk közti magerőnek ebben a hasonlatban a
semleges atomok között ébredő Van der Waals-erők feleltethe t ők meg.) A
leptonoknak nincs színük, tehát az erős kölcsönhatásban nem vesznek részt.

Ha a kvarkok kis távolságban hatnak kölcsön, akkor a kölcsönhatási
energia nagy, a csatolás viszont kicsi, így a kvarkok majdnem szabad részecskeként viselkednek. Ez az ún. "aszimptotikus szabadság". Ebben az esetben
QCD és a QED (kvantum-elektrodinamika) nagyon hasonlóan viselkednek
A helyzet azonban te lj esen e l térő, ha nagy távolságra vannak egymástól a
kvarkok: ebben az esetben ugyanis a csatolás erőssége növekszik. Az, hogy
a kvarkok között ható erő $10^{-12}$ cm távolságban lényegében ugyanakkora,
mint pl. 1 cm távolságban - meghökkentő. (Ez, ha naiv hasonlattal akarunk
élni , akkor azt mondhatjuk, hogy olyan, mint am ikor a rabszolgák össze voltak láncolva, és amíg senki nem távolodott el messze a szomszédjaitól, addig
- az adott kereteken belül - szabadon mozogbattak. Mihelyt azonban egy
is a rabszolgák közül megpróbált távolabbra eljutni, ezzel megfeszítette a
láncot, amely visszatartotta. Hasonlatunkban a láncnak a színerők felelnek meg.)

\subsubsection*{Szubelemi részecskék rendezése}
Önkéntelenül felhívja magára a figyelmet az, hogy mind a kvarkokat, mind
a leptonokat három dublettcsoportba lehet elrendezni. Ezek szerint ésszerű
egy közös elrendezés az alábbiak mintájára:
\begin{equation}
    \begin{pmatrix}
        \nu_e & u_p & u_z & u_k \\
        e^- & d_p & d_z & d_k
    \end{pmatrix} 
\end{equation}
ahol a kvarkok indexei a megfelelő színekre (p = piros, z = zöld, k = kék)
utalnak. Tudjuk azonban, hogy a fentiekben feltüntetetteknél több lepton
és több kvark létezik, ez csak az első lepton-, illetve kvarkdublettnek az
összefoglalása. 
\begin{figure}[H]
    \centering
    \includegraphics[width=0.8\textwidth]{imgs/13-tetel/rendszerezes.png}
    \caption{A szubelemi részecskék rendszerezése}
\end{figure}
Alul helyezkedik el a fentiekb en ismertetett legegyszertíbb, a közvetlen
környezetünket alapjában meghatározó ún. első generáció. Feljebb következik a második, majd a harmadik generáció. A kvarkok rovatában az első
oszlop a piros ($p$ index), a második a zöld ($z$ index), a harmadik a kék($k$ index) színnek felel meg. A részecskék mellett feltünt ettük mindenütt az
elektromos töltést is. Mint. már említettük , valamennyi atom és az összes
közönséges anyag nyolc részecskéből , azaz az első generáció tagjaiból állítható össze. Figyeljük meg, hogy minden egyes generációnál az elektromos
töltéseknek az összege 0-t ad . Pl. az első generációnál: $-1 + 3(1/3) = 0$. Lényeges pont, hogy a töltések összege csak akkor ad 0-t, ha a kvarkokat három
különböző színben, azaz három különböző kvantumszámmal rendelkezőnek
tekintjük Ellenkező esetben nem jön ki a zérus, és ezt rendkívül erős érv
amellett, hogy valóban núnden egyes kvarkfajta három különböző kvantumszámmal, színnel rendelkezik. Ha csak a kvarkíz létezne és szín nem, akkor
az első generáció tölt.ésére a következőt kapnánk: $- 1 + (2/3)- (1/3) = -2/3$.

Az ún. Híggs-részecskék ($H$) (skaláris meunyiségek) rendkívül fontos
szerepet játszanak az elméletben, ők képezik annak a mechanizmusnak az
alapját , ame ly a leptonok , kvarkok és közbenső bozonok tömegének a megjelenéséhez vezet. E részecskék kísérleti kimutatása a részecskefizika előtt
álló rendkívül fontos, alapvető jelentőségíí feladat.
Az alacsony energiákon előforduló 8 részecske ($e$, $\nu$, $u$ és $d$, az utóbbiak
3- 3 színben) mellett vannak még további részecskék is, amelyek a fizikában
fellépnek. Figyelembe kell vennünk először is a részecskék antirészecskéit,
ami számukat 16-ra emeli. Fel kell még tételeznünk, hogy vannak bal és
jobb helicitású részecskéi->, ami a számot ismét megduplázza: tehát 32-re
emelkedne a részecskék száma. A valóságban - eddigi tapasztalataink szerint - nincsen jobb hclicitású neutrínó és bal helicitású ant ineutrínó , ezért
azt kell mondanunk, hogy jelenleg a "legelemibb" részeknek a száma 30 körül mozog. Ez meglehetősen sok, de mindenesetre a több száznál lényegesen
kevesebb. A 24 .4. ábrán tréfás rajzban tüntettük fel a 6 különböző "ízű" kvarkot.

Figyelembe véve az ismert leptonok és kvarkok viszonylag nagy számát
és azt a tényt , hogy a leptonoknak egész számú töltései vannak, felmerülhet
egy olyan elképzelés, amely szerint lehet, hogy a leptonok és talán a kvarkok
is összetett részek, és néhány "szubkvarkból" állnak. Van egy olyan elképzelés. amely szerint kvarkok és leptonok két fundamentális részecske kötött
állapotai . Az egyik töltése $+ 1/3$ lenne, a másiké pedig 0.


\subsubsection{Kölcsönhatások}

\subsubsection*{A gravitációs kölcsönhatás}
Mint ahogy láthattuk, ez egy rendkívül gyenge kölcsönhatás,
amelynek a hatótávolsága végtelen: azaz két tömeg akármilyen nagy távolságból is kölcsönhat egymással a gravitációs erők közvetítéséveL A távolság
növekedésével a kölcsönhatás erőssége (négyzetesen) csökken és a végtelenben nulla felé tart. Kivétel nélkül minden testre hat, a kölcsönhatásban részt
vevő testek tömegének szorzatával arányos. Ennek megfelelően a gravitációs kölcsönhatás jelentősége kozmikus méretekben óriási, ez kormányozza
ugyanis az Univerzum egészét, a Naprendszert, és ily módon meghatározza
életfeltételeinket, atomi méretekben azonban elhanyagolható, ezért e helyütt
nem foglalkozunk vele. Lehet az is, hogy a gravitáció nem tartozik a hagyományos értelemben vett kölcsönhatások közé, hanem a téridő tulajdonságának a következménye.
\begin{figure}[H]
    \centering
    \includegraphics[width=0.6\textwidth]{imgs/13-tetel/grav_kolcsonhatas.png}
    \caption{
        Szemléltető rajz a kicserélődési erőkre. A két csónakban labdázó
        lányok között a kapcsolatot a "közvetítő részecske", a labda hozza létre.
    }
\end{figure}

\subsubsection*{Az elektromágneses kölcsönhatás}
Az elektromágneses kölcsönhatás körébe tartozó kémiai, biológiai stb. jelenségek szabják meg mindennapi életünket. Az elektromágneses kölcsönhatás volt az első, amelynek elméletét pontosan kidolgozták,
és a kvantum-elektrodinamika ( QED) egyike a fizika legpontosabb és legidőtállóbb területeinek.
rületeinek. Eddig nem találkoztunk olyan jelenséggel, amely a QED-val letve mágneses nyomatékkal rendelkező testek között hat. Lényeges különbség a gravitációval szemben, hogy
míg a testek gravitáció révén mindig csak vonzzák egymást, add ig az elektromágneses kölcsönhatás lehet vonzó és taszító is.
\begin{figure}[H]
    \centering
    \includegraphics[width=0.4\textwidth]{imgs/13-tetel/em_kolcsonhatas.png}
    \caption{
        Az erők kicserélődési képének szemléltetése Feynman-gráffal. Jelen esetben elektron
        szóródik eelektronon és a kölcsönhatást egy $\gamma$ kvantum közvetíti.
    }
\end{figure}
A kvantumelektrodinamikában két elektromosan töltött részecske közötti kölcsönhatást egy harmadik részecske közvetíti. Ez a részecske a foton,
az elektromágneses sugárzás kvantuma. A foton tömege nulla, nincs saját
elektromos töltése és fénysebességgel mozog. Az elektromágneses erőket úgy
írhatjuk le , mint a fotonoknak a cseréjét (kicserélődési erők): az egyik töltött
rész kibocsát egy fotont, a másik elnyeli és így "labdáznak" egymás között a fotonokkal.

\subsubsection*{A gyenge kölcsönhatás}
E kölcsönhatás gyengesége miatt semmilyen rendszert nem képes összetartani. Ez a kölcsönhatás felelős részint az atommagok radioaktív bomlásáért,
s ez szabályazza a csillagok (köztük a Nap) energiatermeléséf a termonukleáris fúziót - közvetve tehát létfontosságú az élet számára. Hatótávolsága
rövid ($\sim 10^{-15} cm)$: elsősorban az atommagok belsejében uralkodó kölcsönhatás.
A gyenge kölcsönhatások típusába tartozik valamennyi ,$\beta$-bomlás. Ennek néhány esete:
\begin{equation}
    n \rightarrow p + e^- + \tilde{\nu}_e, \quad \quad \mu^- \rightarrow e^- + \tilde{\nu}_e + \nu_\mu, \quad \quad \pi^- \rightarrow \mu^- + \tilde{\nu}_\mu
\end{equation}
A gyenge kölcsönhatásokkal kapcsolatban meg kell ismerkednünk egy új
fogalommal, és ez a helicitás ($H$). Az 1/2 spinnel rendelkező részecskéknek
kétféle orientációja lehetséges, a spin ($s$) vagy a részecske mozgásának irányába (P) mutat
($H= +1$), vagy vele ellentétes irányban ($H = -1$). Az előbb i esetet nevezzük "jobbkezűségnek", mivel ha a
jobb kéznek az ujjai gondolatban körülfogják a részecskét, akkor a hüvelykujj a mozgás irányát mutatja.
Ellenkező esetben a bal kéznek a hüvelykujja mutatja a részecskének a
mozgását, ekkor "balkezűségről" beszélünk. Általában az irányítottság, azaz a helicitás megfordítható: egyszerűen lefékezzük a részecskét, nyugalomba hozzuk és ezután ellenkező irányba
rűen lefékezzük a részecskét, nyugalomba hozzuk és ezután ellenkező irányba
\begin{figure}[H]
    \centering
    \includegraphics[width=0.6\textwidth]{imgs/13-tetel/helicitas.png}
    \caption{
        A helicitás értelmezése. A z ábra bal oldalán egy "jobbkezes" rendszert látunk, melynél a spin
($s$) iránya és az impulzus ($p$) iránya egybeesik. Ebben az esetben konvencionálisan $+1$-es helicitásról beszélünk.
Az ábra jobb oldalán a "balkezes" helicitás látható
    }
\end{figure}
gyorsítjuk fel őket anélkül, hogy megzavarnánk a spint. Ennek megfelelően
a részecskék mindegyikének van egy bal- és egy jobbkezes komponense. Kivételt képeznek a neutrínók (ha a neutrínótömegre utaló jelektől itt most
eltekintünk) , mert ezek mindig fénysebességgel mozognak és soha nem hozhatók nyugalomba. Ennek megfelelően a neutrínók esetében csak egyféle
helicitás lehetséges:
\begin{equation}
    H_\nu = -1
\end{equation}
Ez a tény tükrözi leg szembetűnőbbe n a tértükrözési szimmetria sérülését,
azaz a paritás meg nem maradását a gyenge kölcsönhatásban. Kísérletileg
eddig csak "balkezes" neutrínókat figyeltek meg. (Az antineutrínók viszont
mindig "jobbkezesek".) Ugyanúgy, ahogy az elektromágneses kölcsönhatásokat a foton közvetíti, gyenge kölcsönhatásoknál is van közvetítő részecske,
pontosabban részecskék, mégpedig számszerint három: $w^+, w^- $ és $Z^0$. Az
első kettő hat az elektromos töltésre is, a harmadik a töltést nem változtatja
meg. Mindhárom közvetítőnek, ún. vektorbozannak a tömege nagyságrendileg $100 GeV/c^2$, a proton tömegének 100-szorosa!

\subsubsection*{Az Erős kölcsönhatás}
Az erős kölcsönhatásokkal ugyancsak a magfizikában találkoztunk, a magot felépítő részecskék között: a részecskék széles kategóriájában uralkodó,
rendkívül heves, vehemens kölcsönhatástípus. A magerő igen rövid hatótávolságú. Közvetítő részecske a $\pi$. (Abból a feltevésből , hogy a kölcsönhatást egy m tömegű részecske közvetíti,
meg lehet becsülni a hatótávolságot. A közvetítő kibocsátásához $E > mc^2$ energia szükséges, ez az állapot viszont a $\Delta E \cdot \Delta t \approx \hbar$ Heisenberg- féle határozatlansági összefüggésből
$\Delta t \approx \hbar/mc^2$ ideig maradhat fenn. Ezáltal a közvetítő részecske legfeljebb
$d = c \Delta t \approx \hbar/mc$ utat tehet meg. Ha $m$ helyébe a pion tömegét írjuk, $d$-re kb. $10^{-13}$ cm, a magerők hatótávolsága adódik.) A hatására végbemenő
folyamatok rendkívül gyorsan ( $\approx 10^{-23}$ s alatt) zajlanak le. (Ez egy történelmi, Yukawa japán elméleti fizikus által elgondolt kép.)

Az erős kölcsönhatás elmélete az ún. kvantum-színdinamika (kvantumkromodinamika, QCD). Egyelőre ezzel az elmélettel a kísérletek megegyeznek, azonban nagyon sok következtetését még nem sikerült igazolni,
továbbfejlesztése várható. Az erős kölcsönhatások részleteire a későbbiekben a kvarkokkal kapcsolatban még visszatérünk. Ma már t udjuk, hogy a magerők
nem elemi erők, mint pl. az elektromágneses erők, hanem indirekt következményei a kvarkok között ható valóban elemi, ún. "szín"-erőknek. (A magerők a Coulomb-kölcsönhatásbeli Van der Waals - erőknek
felelnek meg.) A színerőket az ún. gluonok közvetítik. A részecskék mindegyike részt vesz valamelyik kölcsönhatásba n. Vannak olyanok, amelyek csak
egyetlen kölcsönhatásban vesznek részt , így pl. a neutrínó egyetlen lehetséges kölc sö nha tása a gyenge kölcsönhatás (eltekintve a gravitációs kölcsönhatástól,
amelytől azonban atomi részecskék esetében eltekintünk). Mások több kölcsönhatásban vesznek részt, pl. az elektron, amely részese a gyenge,
elektromos töltése révén az elektromágneses és tömege következtében a gravitációs kölcsönhatásnak is (megjegyezzük, hogy ez utóbbi nagyságrendekkel kis ebb effektus az atomi mérettartományban, mint az első kettő).
A $\pi^\pm$ pedig még az erős kölcsönhatásban is részt vesz, az előbb említett három kölcsönhatáson túl.

További információ az egyesített elméletekre tett kísérletekrők: \href{http://fizweb.elte.hu/download/Fizika-BSc/Mag-es-reszecskefizika/Jegyzet/Horvath_Kiss-Kiserleti_atomfizika.pdf#%5B%7B%22num%22%3A1637%2C%22gen%22%3A0%7D%2C%7B%22name%22%3A%22FitH%22%7D%2C678%5D}{23. fejezet}



\subsection{Kísérleti eszközök}
A részecskefizika berkein belül kétféle fő komponenst különböztethetünk meg a kísérleti eszközök terén: a \textbf{részecskedetektorokat} és a \textbf{részecskegyorsítókat}.
A részecskedetektorok felelnek azért, hogy meghatározzuk a detektálni kívánt részecskék helyét, energiáját, mozgásának irányát, impulzusát, tömegét, töltését, stb. amiket mind különböző
konfigurációkkal tudunk megkapni.

A Gyorsítókra pedig azért van szükség, hogy a nagyobb és nehezebb részecskéket nagy energiával tudjuk ellátni, és azokat ütköztetve kisebb, elemi récsecskéket tudjunk előállítani. Ezzel kontrollált módon
tudjuk vizsgálni a részecskék kölcsönhatásait, és tulajdonságait. Ha nem áll rendelkezésünkre ilyen eszköz, akkor a kozmikus sugárzásból származó részecskéket tudjuk csak vizsgálni, amiknek az energiája és típusa véletlenszerű,
emiatt csak a szerencsén múlik, hogy képesek vagyunk-e egy adott részecskét detektálni. Előnye viszont a kozmikus sugárzásnak, hogy a gyorsítókban nem előállítható energiájú részecskék is meg tudnak jelenni
(egy részecskegyorsító kb. $10^{13}$ eV energiájú részecskéket tud előállítani, míg a kozmikus sugárzásban $10^{20}$ eV energiájú részecskék is megjelennek). Ezért az extrém magas energiájú részecskék vizsgálatára csak a kozmikus sugárzás a alkalmas jelenleg.

\subsubsection{Gáztöltésű számlálók}
A gáztöltésű számlálók alapját az \textbf{ionizációs kamra} képezi, amely az egyik legegyszerűbb részecskedetektor. A lényege, hogy van egy zárt terünk, amiben egy gázszikgeteléses kondenzátort helyezünk.

Ha a kondenzátor elektródjai között ionizáló részecskék haladnak át, a gázt ionizálják Az elektródokra kapcsolt $V$ feszültség hatására az ionok az elektródra vándorolnak, a $G$ műszer áramot jelez.
\begin{figure}[H]
    \centering
    \includegraphics[width=0.8\textwidth]{imgs/13-tetel/ionizacios_kamra.png}
    \caption{Ionizációs kamra felépítése és működése}
\end{figure}
A görbén három szakasz különböztethető meg. A középső feszültségtartományban ($B$) az elektromos tér minden, a sugárzás által keltett iont, és
csakis azokat, begyűjti: az áram független a feszültségtőL Ha a feszültség
túl kicsi, a töltéshordozók olyan lassan mozognak, hogy van idejük rekombinálódni, tehát nem mindegyik jut el az elektródokhoz ($A$). A jobb oldali
emelkedő szakaszról ($C$) a későbbiekben lesz szó. Az áramerősség igen kis érték. Ezért új módszereket, eszközöket kellett keresni.
Ma ionizációs kamrákat főleg nagy intenzitások mérésére (pl. reaktorok-ban) doziméterként és egyes kozmikus sugárzási mérésekben használnak.

\subsubsection*{Proporcionális számláló}
A proporcionális számláló egy ionizációs kamra, amelyben a gáz ionizációját mérjük.
Az elég nagy elektromos tér hatására mozgó szabadelektronok gerjesztik a gáz molekuláit. Ha a térerősség elég nagy, az elektronok a két ütközés
közötti szabad úthosszon annyi energiát nyerhetnek, mint a gáz ionizációs potenciálja. A következő ütközés alkalmával minden elektron ionizálhat egy
atomot, s ezután már két elektron megy tovább. A következő ütközésnél mindkettő kivált egy-egy újabb elektront és így tovább: elektronlavina alakul
ki. A kialakult lavinák nagysága és így a kapott jel arányos (proporcionális) a primer elektronok számával.

Ez a térerősség túl nagy feszültség használata nélkül könnyen megvalósítható. Ha a számláló katódja $r_1$ sugarú cső, amelynek belsejében koncetrikusan
helyezkedik el az $r_2$ sugarú anód, a tengelytől $r$ távolságban a térerősség
\begin{equation}
    E(r) = \frac{V}{r \ln(r_2/r_1)}
\end{equation}

Ha $r_1 = 100mm$, $r_2 = 0.05mm$ és $V = 2000V$, a szál felszíne felett $0.05mm$ távolságban a térerősség $1.3 \cdot 10^4 V/cm$, ami már kb. elég a lavina megindításához.
\begin{figure}[H]
    \centering
    \includegraphics[width=0.6\textwidth]{imgs/13-tetel/proporcionas_szamlalo.png}
    \caption{Proporcionális számláló felépítése és működése}
\end{figure}

szokásos töltőgázok ionizációs potenciálja 15 eV körül van. Rekombinációkor ilyen energiájú (ultraibolya) fotonok sugárzódnak ki. Mivel a szokásos katódok ki lépési munkája néhány (3- 5) eV, a nagy energiájú fotonok
nagy valószínűséggel váltanak ki fotoelektronokat, amelyek újabb lavinákat indíthatnak el, tehát állandó kisülés lép fel.

A fotoeffektus megszüntethető úgy, hogy a csövet olyan gázzal (pl. metán) töltjük meg, amyelnyek erős az abszorpciója az ultraibolya tartományban.
Így a fotonok nem juthatnak el a katódhoz, a gázerősítés megfelel az elméletileg vártnak. Amikor a menánmolekula fotont nyel el, széthasad, ezzel felhasználódik.
Ezért a cső élettartalma korlátozott, nagyságrendben $10^{12}$ impulzus.

Mint azt láttuk, a térerősség csak a szál közvetlen közelében elegendő a lavina megindításához. A cső térfogatának túlnyomó részében csak arra elég,
hogy a keletkezett elektronokat a szál felé továbbítsa. Ezért gázerősítés lényegében független lehet a részecske áthaladásának helyétől.

A síkgeometriájú proporcionális számlál ők mérete néhány cm-től l méterig (kozmikus sugárzási célokra) terjedhet, a síkok távolsága általában
néhány cm. Az anód kb. 50 J.Lm vastagságú volfrámszál, a katód valamilyen fém, pl. vörösréz vagy alumínium (1 - 2 mm vastagságú). A töltőgáz metán,
ha a gázt állandóan cseréljük, akkor a számláló élettartama nagyon hosszúra növelhető. Az alkalmazott feszültség 2000 V körül mozog.
A proporcionális számlálók egyik legfontosabb fe lhasználási ter ül ete a neutronfluxus mérése. Ha a a számlálót bór-trifluorid ($BF_3$) gázzal töltjük, a
lassú neutronok nagy valószínűséggel váltják ki a $^{10}B + n \rightarrow ^{11}B \rightarrow ^{7}Li + ^{4}He$ reakciót amelyben összesen kb. $2.5 MeV$ energia szabadul fel. A kapott jel
jól érzékelhető és könnyen elválasztható az esetleges $\gamma$-kvantum háttértől is
(egy 1-kvantum mindössze $10^5 eV$ nagyságrendű energiát veszíthet a számlálóban). Gyors neutronok fluxusa úgy mérhető, hogy a számlálót paraffinba
vagy más, nagy hidrogéntartalmú anyagba ágyazzuk, amely lelassítja a neutronokat. 

\subsubsection*{Proporcionális kamra}
Proporcionális számlálók úgy is kialakíthatók, hogy egy lapos tartályban egymással párhuzamosan sok szálat feszítünk ki, ezek lesznek az anódszálak.
A szálsík fölött és alatt fémlemezek helyezkednek el, amelyek a katódot alkotják. A gáztér közös, ennek ellenére az egyes szálak egymástól függetlenül mint önálló proporcionális számlálők képesek működni.
\begin{figure}[H]
    \centering
    \includegraphics[width=0.6\textwidth]{imgs/13-tetel/proporcionalis_kamra.png}
    \caption{Proporcionális kamra felépítése és működése}
\end{figure}

A töltőgáz nyomása lényegében azonos a légköri nyomással, ezért a beés a kilépési oldalon a kamra rendszerint vékony műanyagfólia . E detektortípus egyik fő előnye,
hogy nagyon kevéssé abszorbeálja és szórja a rajta áthaladó részecskéket. Mivel egy ilyen kamra vákuumtechnikai szempontból
"piszkos", állandó gázöblítésről kell gondoskodni. Ha minden szálhoz külön erősítő tartozik és két, egymásra merőleges szálrendszert alkalmazunk,
mm nagyságrendű pontossággal megkapjuk a részecske áthaladásának helyét (koordináta-detektor). A kamra felhasználható triggerjel-keltésére is, az
elérhető felbontási idő- elsősorban az elektronikától függően- 10 ns-100 ns.

A proporcionális kamrák az utóbbi években uralkodó szerephez jutottak
a nagyenergiájú fizikában. Méretük négyzetdecimétertől több négyzetméterig terjed. A szálsűrűség általában 1-2 szál három milliméterenként. Egy
mérésben $10^2 -10^5$ szálat használnak, mindegyiket a saját erősítőjéveL Mivel egy szál elektronikája néhány dollár körüli összegbe kerül, ezek a detektorok rendkívül költségesek.

\subsubsection*{Geiger-Müller-számláló (GM-cső)}
Ha egy proporcionális számláló tápfeszültségét a normális üzemfeszültségfölé növeljük, eljuthatunk egy olyan értékhez, amelynél a gázerősítés végtelen lesz: az elektron által megindított kisülés önfenntartóvá válik. Tiszta
nemesgáztöltés esetén ez akkor következik be, ha az első lavinában keletkezett fotonok számának és a katódon a fotoeffektus valószínűségének szorzata
1 lesz. Ha kívülről gondoskodunk arról, hogy a kisülést kioltsuk, nagyon értékes számlálótípust nyerünk: olyan számlálót, amelyben a részecske áthaladása csak kiváltja az elektromos impulzust, de annak nagysága már
független a részecske energiájától (kiváltó számláló). A kisülésből kapott jel lényegesen nagyobb a proporcionális számlálóénál, tehát egyszerűbb elektronika is elegendő.

Lehetséges olyan .töltőgázkeveréket (pl. argon + alkohol) találni, amelylyel a számláló önkioltóan működik: egy kisülési ciklus után visszaáll alapállapotába.
önkioltó számláló működése időben három szakaszra bontható:
\begin{enumerate}
    \item A részecske által keltett szabadelektronok a szál felé (anód), az ionok a "ház" felé (katód) mozognak.
    \item Az anód közelében kialakult nagy térerősség hatására az elektronok lavinát indítanak el, amely a szálon nagy áramimpulzust hoz létre. 
          a gázkeverékben a lavinákból származó ultraibolya fotonok abszorpciós hossza kb. 0.1 mm nagyságrendű, ezért a katódra nem jutnak el, de a lavinák
          közelében újabb szabadelektronokat hozhatnak létre, amik újabb lavinákat indítanak el. A kisülés így szétterjed a szál teljes hosszában, és a végén kialszik.
          A lavinákban keletkező elektronok belépnek a katódba, a nagy tömegű pozitív ionokból pedig egy tértöltésfelhő marad vissza a szál közelében.
    \item A tértöltésfelhő az elektromos tér hatására a ház felé mozog. Ha a nemesgázionok megérkeznének a katódra, azzal újabb kisülési szakaszt indítanának el.
          De ha jól választjuk meg a kisérő gázt (pl. alkohol), hogy az ionizációs potenciálja kisebb legyen (kb. 11.3 eV), akkor töltéscsre történik a két gáz között 
          és semleges argonatomokat és ionizált alkoholt kapunk. Az alkoholionoknak viszont disszociálnak, a maradék energiájuk már kevés, hogy újabb kisülést indítsanak be,
          ezért a rendszer visszaáll alapállapotába. Viszont amíg a pozitív ionfelhő el nem megy a szál körül, addig a GM cső érzéketlen lesz (holtidő). Ez olyan 100-300 $\mu s$ ideig tart.
\end{enumerate}

Kísérőgáznak lehet még halogéngázt is választani (pl. brómgőz), ez hosszabb élettartamú (az alkohol kb. $10^8-10^9$ impulzus után elhasználódik),
de a holtidő is hosszabb lesz (néhány ms) ezért koincidencia mérésekre nem alkalmas.

\subsubsection{Szcintillációs számlálók}
Az egyik legrégibb detektálási eljárás a sugárzások által bizonyos kristályokban (pl. $ZnS$) keletkező fényfelvillanások megfigyelésén alapul. Ezt alkalmazták a magfizikai kutatások kezdetén.
A keltett fe lvillanásokat vizuálisan észlelték~eszámlálták. Bár a mérések fáradságosak és szubjektívek
voltak , mégis ezen eszköz segítségével jutottunk az atom alapvető szerkezetének megismeréséhez (Rutherford-szórás). Az 1930-as évekig ez az eszköz
volt az egyetlen sugárzásdetektor.
A szcintillációs számláló továbbfejlődését a fotoelektron-sokszorozók
megjelenése jelentette: a vizuális megfigyelést helyettesítették fotoelektronsokszorozóval.

Felépítése:
\begin{itemize}
    \item Szcintillátor: anyaga lehet szervetlen kristály (pl. $NaI(Tl)$, $CsI(Tl)$, $BGO$) vagy szerves anyag (pl. folyadék: toluol, szilárd: polisztirol).
          A szcintillátor feladata az ionizáló sugárzás energiájának fényenergiává alakítása.
    \item fotoelektron-sokszorozó: a szcintillátorban keletkezett fényt elektromos jellé alakítja át és felerősíti.
    \item Fotócső: a fényt elektromos jellé alakítja át.
\end{itemize}

\subsubsection*{A szcintillátor}
A szcintillátor anyagokban a számlálóba belépő részecskék, illetve $\gamma$-kvantumok
energiájának egy része fényenergiává alakul át. Ez töltött részecskék esetén
közvetlenül valósul meg: semleges részecskék és $\gamma$-kvantumok csak másodlagos folyamatok segítségével detektálhatók.
Univerzális szcintillátorok nincsenek. ,$\beta$-részecskék, nehézionizáló részecskék és $\gamma$-sugárzás detektálására más-más szcintillációs anyagokat alkalmazunk,
amelyekben a fényfelvillanás kialakulásának és lefolyásának mechanizmusa különböző.

A primer kölcsönhatás általánosságban a következő két lépésre bontható:
\begin{enumerate}
    \item A beérkező részecske vagy $\gamma$-foton energiát ad át a szcintillátornak, amely gerjesztett állapotba kerül.
    \item A szcintillátor gerjesztett állapotából fotonok, ún. fluoreszcenciafény kibocsátásával tér vissza alapállapotába.
\end{enumerate}

A fényfelvillanásole száma a szcintillátorba jutó $\gamma$-fotonok számával, az
egyes felvillanásban a fotonok száma a $\gamma$-kvantum energiájával arányos.
A szcintillációs detektorok egyik fő előnye a nagy fényhozam és a kitűnő időfelbontás.

\subsubsection{Félvezető dsetektorok}
Az 1960- as évek e lejétő l a félvezető detektorok gyors fejlőd ésé nek és térhódításának leh et ünk tanú i. E detektorok fő e l őnye , hogy kis méretük ellenére igen jó energiafelbontó képességük van,
s ennek révén alkalmazásuk mindazokon a területeken előnyös, ahol töltött részecskék vagy $\gamma$-sugárzás
energiáját kell mérni, de nincsen szükség néhány $cm^2$-t meghaladó érzékeny
felületre vagy 100 $cm^3$-nél nagyobb érzékeny térfogatra. Ennek alapján a
félvezető detektorok elsősorban a magfizikai méréseknél terjedtek el széles-
körben, de sikerrel alkalmazzák más területeken is (aktivációs analízis, több-
szörösen nyomjelzett vegyületek mérése, ipari és orvosi izotóptechnika).
A különböző félvezető anyagok közül félvezető magsugárzás-detektorok
előállítására széleskörűen csak a szilíciumot és a germániumot alkalmazzák
A sziliciumból készült detektorok készítése egyszerűbb, nem csak alacsony
hőmérséklettel használhatók, elsősorban töltött részecskék spektrometriájára szolgálnak. A germániumdetektorok előá llítása költséges, csak a folyékony nitrogén hőmérsékletén használhatók, de mivel a germániumnak a szilíciuménállényegesen nagyobb a rendszáma,
sokkal e lőnyösebbek $\gamma$-sugárzás spektrometriájára, mint a szi lí ciumdetektorok. A germániumot vagy Li-mal
aktiválva, vagy nagyon tiszta (high purity, HPGe) elemi állapotban használják detektornak.

\subsubsection{Cserenkov-számlálók}
Nagysebességű töltött részek detektálására fellehet használni a Cserenkov-sugárzást. Az egyik leggyakoribb, fókuszáló típusú számláló az ábrán látható. A radiátorban Cserenkov-sugárzás jön létre, ezt tükrökkel
fotoelektron-sokszorozóra ejtjük , ahol elektromos jellé alakul.
Mivel az elektronsokszorozóra csak egy bizonyos szűk szögtartományban kilépett fotonok juthatnak, ez a számláló alkalmas valamilyen szűk sebességintervallumba eső sebességű részecskék kiválasztására.
\begin{figure}[H]
    \centering
    \includegraphics[width=0.4\textwidth]{imgs/13-tetel/cserenkov_szamlalo.png}
    \caption{Cserenkov-számláló felépítése és működése}
\end{figure}

Igen nagy energiájú részecskéket lehet detektálni úgy, hogy a közeg, az ún. radiátor anyaga gáz. A levegő törésmutatója 1,000292, levegőben a
minimális $\beta = 0,999708$, ilyen sebességű részecskére $\frac{m}{m_0} \approx 3400$. A légkört
használják radiátornak egyes kozmikus sugárzási mérésekben. A Cserenkov-sugárzás frekvenciaspektruma szűk: a látható és a közeli
ibolya tartományra szorítkozik. Minél jobban megközelítjük a küszöbsebességet, annál kisebb az intenzitás, a küszöbsebességnél zérus. A Cserenkov-sugárzás látványos módon
jelentkezik az atomreaktorok vízmoderátorában. A hasadási termékekből nagy intenzitású, nagy energiájú radioaktív sugarak lépnek ki, amelyek a moderátor vízében Cserenkov-sugárzást keltenek,
és emiatt a víz jellegzetes, kékeslila fényben világít.

\subsubsection{Részecskenyom-detektorok (vizuális detektorok)}
A létező detektorokat durván két csoportba sorolhatjuk: számlálők és
részecskenyom-detektorok Gyakran az előbbit mondjuk röviden detektornak. Az első csoportba azokat az eszközöket soroljuk, amelyek egy részecske
adott időben, adott helyen való megjelenéséről adnak számot (a pontosság $\sim 0,1 mm$ és $\sim 10^{-9} s$ lehet). A nyomdetektorokban a részecske pályája
mentén nyomot hagy, amelyet gyakran trek-nek (track) neveznek. A pályát
valamilyen módszerrel, pl. fényképezéssel rögzítik. A pálya adataiból (ionizációs sűrűség, görbület mágneses térben stb.) sokkal több információt lehet
nyerni, mint a számlálók esetében. 

Nyomot csak az ionizáló részecske hagyhat. Mégis, semleges részecskék
esetében is alkalmaznak nyomdetektort, ilyenkor a treket (vagy trekeket)
ionizáló szekunder termékek hozzák létre.
A részecskenyom-detektorok jelentősége növekedett az automatikus,
számítógép segítésével megvalósított nyomkiértékelő rendszerek terjedésével. A nyomok kézi kiértékelése mikroszkóppal rendkívül munkaigényes, és
ez a körülmény erősen fékezte a részecskenyom-detektorok alkalmazását. Ma már alkalmazásuk erősen korlátozott.

\subsubsection*{Ködkamra}
A legrégibb nyomdetektor a ködkamra (Wilson-kamra, 1912). Egy gáz- és gőzkeverékkel töltött edényben túltelítettséget hozunk létre, a gőz kicsapódik a jelenlevő gáz ionokra, majd a kicsapódott ködcseppek tovább növekednek és láthatóvá válnak. Így ki lehet mutatni az elemi részecskék pályáit, mert a töltött részecskék a haladásuk nyomában ionizálják a gázatomokat.
A ködkamrában a túltelítettséget úgy hozzák létre, hogy a gáz-gőz
térfogatot adiabatikusan kitágítják. Ennélfogva a ködkamra működtetése
három fázisból áll:
\begin{enumerate}
    \item A kamrát telítjük gáz- és gőzkeverékkel (pl. levegő és vízgőz).
    \item A térfogatot gyorsan megnöveljük, így túltelítettséget hozunk létre.
    \item A részecskék áthaladnak a kamrán, ionizációt hoznak létre.
\end{enumerate}
A kamraalján levő dugattyú hirtelen lefelé mozgatása biztosítja az
adiabatikus expanziót. A nyomokat a
felső üveglap felett elhelyezett fényképezőgép regisztrálja. A kialakult ködcseppeket oldalról történő megvilágítással teszik láthatóvá.

A ködkamra a kozmikus sugárzási kutatásokban kapott jelentős szerepet, számos elemi részecskét, így pl. a pozitront , a $\mu$-t , a $\pi$-t a ködkamrafelvételek segítségével fedezték fel.

A kis anyagsűrííség, másrészt a jelentős ho l tidő miatt, amit jórészt a
működés harmadik (előkészítő) fázisa igényel, a ködkamrával csak kevés
eseményt lehetett. összegyűjteni, vagyis nehéz volt megfelelő statisztikát biztosítani. Jelentős előny volt viszont a kamra vezérelhetősége: a kamra köré
elhelyezett számlálódetektorokkal, illetve ezek koincidenciáival indított kamráknál csak meghatározott típusú események esetén működtették a kamrát.
Bár a külső vezérelhetőség jelentős előny, a ködkamrát fokozatosan kiszorították az egyéb nyomdetektorok.

\subsubsection*{Buborékkamra}
Az utóbbi két évtizedben a buborékkamra kiemelkedő sze repet töltött be
a kísérleti részecskefizikában, a komplex, több részecskés kölcsönhatások
vizsgálatában . Mind nagyobb méretfí buborékkamráka-t hoztak létre. Az
első néhány cm nagyságú kamrák mai utódai néhány méter kiterjedésűek.
A buborékkamrában a nyomok (a forráspont fölé) túlfűtött folyadékban alakulnak ki. A működési ciklus elindítása hasonlóan történik , mint a
ködkamráknál, a nyomás hirtelen csökkentésével.
Egy folyadékban normál is körűlmények között véletlenszerűen állandóan keletkeznek és eltfínnek buborékole Tú lfi.ítött folyadé kban azonban
a buborékok nem semmisülnek meg, hanem növekedni kezdenek. A kamrafolyadékban kölcsönható, regisztrálandó, ionizáló részecske általlétr ehozott
ionpárok buborékokat hoznak létre, és igy megfelelő kritikus hőmérséklet.
esetén a folyadék nyomásának hirtelen csökkentése után a részecske pályája
mentén fénykép ezhető méretfí ($\approx 200 \mu m$) buboréksorozat alakul ki.

A buborékkamra térbeli felbontóképességét a bubo~ékok mérete és a
fényképek alapján történő visszaállítás pontossága szalhja meg: kb. mm
nagyságrendű. Az időbeli fe lbontóképességet az érzékenységi idő szabja meg,
ez ms nagyságrendű. A holtidő nagy kiterjedésű kamráknál néhány száz $\mu s$.
Speciális kamrákat is kidolgoztak, amelyekben a következő regisztrálásig
szOkségcs idő rövidebb, $\mu s$ rendű.

A buborékkamra jelentős előnye, hogy a tanulmányozandó fizikai folyamathoz legmegfelelőbb folyadékkal tö l thető. Így használnak töltőanyag
ként folyékony hidrogént, mód van tehát elemi proton céltárgyat használni.
ami egyben detektorként is szolgá l. Továbbá a folyadékatom súlyának növelése jól ehetőséget nyújt foton detektálására, elsősorban ez a szerepe az
ún. nehéz folyadék töltésű (propán, freon, xenon) kamráknak. A folyadék jó
kompromisszum a gáz és a szilárd anyag között - a buborékkamra egyesíti
magában a ködkamTa és magemulzió előnyeit, azok hátrányai nélkül: megfelelően nagy számú eseményt detektál és a kamrában a lkalmazoLt mágneses
tér segítségével nem csak a részecskék töltése határozható meg, de impulzusuk is kielégítő pontossággal mérhető. Tipikus példája a jó kompromiszszumnak,
hogy a nehéztöltésű buborékkamrák ig en jól alkalmazhatók a neutrínókölcsönhatások tanulmányozására, ahol egyszerre van szükség sűrű abszorbensre és jó mérési felt ételekre.

A buborékkamrák igen terjedelmesek, drágák, bonyolult a készítésük és
üzemben tartásuk. Ezért il yen berendezéseket csak olyan kísérletekben alkalmaznak, ahol ez elkerülhetetlen: elsősorban bonyolult, sokrészecskés,
nagyenergiájú folyamatok, valamint ritka jelenségek, pl. neutrínó-kölcsönhatások
vizsgálatánáL Ma már csakne'm teljesen kiszorították őket a modern sokszálas (proporcionális-, drift- stb.) kamrák.

\begin{figure}[H]
    \centering
    \includegraphics[width=0.6\textwidth]{imgs/13-tetel/buborekkamra.png}
    \caption{Tipikus buborékkamra felvétel}
\end{figure}

\subsubsection{elektrosztatikus részecskegyorsítók - Van de Graaff-gyorsító}
\begin{figure}[H]
    \centering
    \includegraphics[width=0.6\textwidth]{imgs/13-tetel/VanDerGraaf.png}
    \caption{Van de Graaff-gyorsító felépítése és működése}
\end{figure}
A harmincas évek elején a részecskegyorsítók fejlődése az elektrosztatikus Van de Graaff
generátor kifejlesztésével folytatódott. VAN DE GRAAFF (1901 - 1967) amerikai fizikus
volt. Az általa kifejlesztett generátor belsejében két henger közé kifeszített, szigetelő-anyagból készült heveder (szalag) fut. A szalag alsó részén
feszültséggenerátorral összekötött fémcsúcsokról töltések áramlanak át a szalag felmenő
ágára. Ezek a töltések a szalag felső részén töltésleszedő csúcsokon keresztül kerülnek
a nagyfeszültségű elektródára. Az így előállított nagyfeszültséget üveg vagy porcelán
gyorsítócsőben lévő elektródarendszerre vezetik, amely egyenletes feszültségeloszlást és
nyalábfokuszálást tesz lehetővé.

A Van de Graaff generátorok feszültségének felső határa ma 10 millió volt körül van.
Ezt azonban csak úgy lehet elérni, hogy az egész nagyfeszültségű elektródát nagy
átütési szilárdságú, nagy nyomású gázzal (pl. SF6, kén-hexafluorid) töltött tartályba
helyezik. A létrejött nagy térerősség ugyanis polarizálja - rosszabb esetben ionizálja - a
gázban lévő molekulákat (különösen, ha vannak benne dipólus-molekulák, pl. víz), és ez
olyan sok töltést tud elvinni az elektródáról, amelynek utánpótlására a szalag már nem
lesz képes. Ezért a szigetelő gázban vízpára legkisebb nyomokban sem lehet jelen. A Van
de Graaff generátorokkal előállított feszültséget nagyon pontosan lehet szabályozni, és
ezért a generátorból nyerhető részecskenyaláb energiája is nagyon pontosan meghatározott.
Nem túl nagy energiájú részecskéket igénylő, precíziós magfizikai mérések végzésére még ma is használják.


\subsubsection{Lineáris Rezonancia Gyorsító (Linac)}
A sztatikus (időben állandó) elektromos mezővel történő részecskegyorsításnál jóval hatékonyabb módszer az, ha ugyanazt a (nem túl nagy) gyorsító-feszültséget többször is felhasználhatjuk ugyanazon részecske
gyorsítására. Ha éppen megfelelő időben ”kapjuk el” a gyorsítandó részecskét, akkor hatalmas energiákig fel lehet gyorsítani viszonylag kis feszültségekkel is.
Mivel ennél a módszernél valamilyen fizikai folyamat által pontosan meghatározott időközönként kell alkalmazzuk a gyorsító feszültséget, a módszert rezonancia-módszernek nevezik.
A rezonancia-módszert alkalmazni lehet egyenes pályájú gyorsítók esetén (lineáris
rezonancia-gyorsító) éppúgy, mint zárt pályájú gyorsítók esetén (ciklotron, szinkrotron).
Szinte valamennyi rezonancia-gyorsító kihasználja azt a tényt, hogy egy fémelektróda
belsejében a részecskékre nem hat az elektromos mezőből származó erő, bármekkora is
a fémelektróda potenciálja. Ezért amíg a részecske az elektróda belsejében tartózkodik,
addig az elektróda potenciálja megváltozhat. Így elérhető, hogy az elektróda elektromos
tere a belépés előtt is gyorsítsa a részecskét, és a kilépés után is.

A lineáris rezonanciagyorsítóban több, egymás utáni, henger alakú gyorsítóelektróda van
\begin{figure}[H]
    \centering
    \includegraphics[width=0.6\textwidth]{imgs/13-tetel/linac.png}
    \caption{Lineáris rezonancia gyorsító felépítése és működése}
\end{figure}

Amíg a részecskék a hengerek belsejében haladnak, az elektromos mező nem gyorsítja
(de nem is lassítja) őket. Gyorsítás az elektródák közötti térrészben történik. Amikor a
töltött részecskék elhagyják az egyiket és a következő felé haladnak, az előbbi taszítja, az
utóbbi vonzza őket. Mire a részecske a következő elektróda-pár közötti térrészhez ér, az
elektromos térerősség előjelet vált, és így ismét gyorsító hatás lép fel. Ahogy a részecskék
az elektródák között haladnak, a sebességük növekszik. Az elektródák hosszát folyamatosan növelni kell, hogy a sebességnövekedés ellenére a nagyfrekvenciás váltakozó feszültség
előjelváltása a kellő pillanatban következzen be. Nagyobb térerősséget - és ezáltal jobb
gyorsítási hatásfokot - lehet elérni, ha a gyorsítócső méreteit úgy választják meg, úgy
”hangolják”, hogy rezonáljon az elektromos mező váltakozásának frekvenciájára. Ilyenkor a gyorsítócső hullámvezetőként üzemel, és a benne kialakuló hullámok amplitúdója
jóval nagyobb lehet, mint rezonancia nélkül. Az elérhető energiát a gyorsítócső hossza
és a nagyfrekvenciás rendszer teljesítményigénye korlátozza.
További megjegyzés az ábrához : Vegyük észre, hogy a gyorsítási szakaszok közötti
távolság (az elektródák hossza) csak a gyorsítócső elején növekszik, utána többé-kevésbé
már állandó marad. Ennek a részecskék (pl. elektronok) relativisztikus viselkedése az
oka. Ha már a fénysebességhez közeli sebességgel mozog egy részecske, akkor az energi-
ájának növekedése elsősorban a tömeg növekedésében nyilvánul meg, a sebessége szinte
alig változik. Ezért a gyorsítási szakaszok közötti távolságnak sem kell már tovább növekedni.

Az egyik legnagyobb Linac gyorsító a Stanfordi Egyetemen (SLAC) működik. Ez a
3 km hosszú gyorsítócsőben 22 GeV energiájú elektronokat és pozitronokat állít elő.

\subsubsection{Ciklotron}
E. LAWRENCE (1901 - 1958, Nobel-díj 1939) találmányaként tartják számon (ezért is
kapott Nobel-díjat), de 1929. január 5-én Szilárd Leó nyújtotta be az első ilyen jellegű
szabadalmat, tőle függetlenül pedig 1929. május 6-án GA ÁL Sándor (1885 - 1972) küldött be egy hasonló tartalmú cikket a Zeitschrift für Physik című német folyóirathoz.
Az elv tényleges megvalósítása, az első működő ciklotron megépítése azonban valóban E.
Lawrence érdeme. Találmányára 1934-ben kapta meg a védelmet, és 1939-ben a Nobel-
díjat.
\begin{figure}[H]
    \centering
    \includegraphics[width=0.6\textwidth]{imgs/13-tetel/ciklotron.png}
    \caption{Ciklotron felépítése és működése}
\end{figure}
A ciklotron közepén lévő ionforrásból lépnek ki a gyorsítandó részecskék. Az üreges gyorsító elektródákat (kékkel, ill. pirossal az ábrán) duánsoknak, vagy
- alakjukra emlékezve - röviden csak D-nek nevezik. A részecskék a duánsok közötti
elektromos mezőben gyorsulnak, a duánsokra merőleges mágneses mező pedig körpályára kényszeríti őket. Könnyű belátni, hogy a körbefutás ideje a sebességtől és a pálya
sugarától független (legalábbis addig, amíg a részecskék nem érnek el relativisztikusan
nagy sebességeket, azaz a tömegnövekedésük elhanyagolható). A körpályán tartáshoz
szükséges centripetális erőt a mágneses Lorentz-erő biztosítja, azaz $\frac{mv^2}{r} = qvB$. Figyelembe véve, hogy $\omega = \frac{v}{r}$, ebből kapjuk
\begin{equation}
    \omega = \frac{qB}{m}
\end{equation}
Azaz a mozgás $\omega$ körfrekvenciája csak a részecske fajlagos töltésétől (q/m) és az alkalmazott mágneses indukciótól függ, tehát nem függ sem a pálya sugarától, sem a részecske
sebességétől, energiájától stb. Mivel $T = \frac{2\pi}{\omega}$ , ezért a körülfutás ideje sem függ attól,
hogy a részecske éppen milyen pályasugáron mozog. A gyorsított részecske energiája
azonban függ a pályasugártól, hiszen $v = r\omega = r \frac{q}{m}B$, és (nem-relativisztikus esetben)
\begin{equation}
    E = \frac{1}{2} mv^2 = \frac{1}{2} m r^2 \left( \frac{q}{m} B \right)^2 = \frac{1}{2} \frac{q^2 B^2 r^2}{m}
\end{equation}
Látható, hogy az elért energiát végeredményben a ciklotron sugara ($r$ maximális értéke),
és az alkalmazott mágneses indukció fogja meghatározni.
A ciklotron három fő része tehát: a gyorsítókamra, a nagyfrekvenciás berendezés és az
elektromágnes. Amikor a részecske elérte a kellő energiát, akkor egy ”kihúzófeszültségre”
kapcsolt elektróda letéríti a körpályáról és a céltárgyhoz vezeti a nyalábot.
A ciklotron a harmincas évek végére elterjedt laboratóriumi eszközzé vált, a gyorsított protonok energiája 20 MeV körüli volt. Ennél nagyobb energiáknál a relativisztikus
tömegnövekedés már nem elhanyagolható. A tömegnövekedés miatt
a nagyobb energiájú részecskék körülfutási ideje megnő, ”kiesnek” a szinkronból, és már
nem gyorsulnak tovább. Lawrence találmánya a frekvenciamodulált ciklotron (szinkrotron) is. Ennél a berendezésnél a gyorsító-frekvenciát a tömeg növekedésével szinkronban
csökkentik. Az ötvenes évek óta a fázisstabilitás és a fókuszálás módszereinek kidolgozása, a szupravezető mágnesek és a kaszkádgyorsítás módszerének alkalmazása segítségével
egyre nagyobb energiájú gyorsítók épülhettek.

\subsubsection{CERN}
A Föld egyik legnagyobb gyorsítóberendezés-komplexuma Svájcban van Genf mellett.
A neve CERN, amely a létesítmény francia elnevezésének kezdőbetűiből alkotott betűszó: Conseil Européen pour la Recherche Nucléaire (Európai Nukleáris Kutatási Tanács).
Bár a neve nukleáris kutatásokat sejtet, manapság már hagyományos magfizikai (nukleáris) kutatás alig folyik benne. A
CERN a nagyenergiájú fizikai és részecskefizikai kutatások legnagyobb központja. A
CERN nemzetközi státuszú kutatóközpont, működését a tagországok közösen biztosít-
ják. Magyarország is tagja a CERN-nek, ezért magyar kutatók is részt vehetnek az
ottani óriásgyorsítókkal folyó kísérletekben, használhatják mindazt az infrastruktúrát,
számítógép-kapacitást, programokat, amelyeket a CERN-ben fejlesztettek ki az évtizedek során.

    A Nagy Hadronütköztető átlagosan 100 m föld alatti mélységben futó, 27 km kerületű kör alakú alagútban épült Genf közelében, a svájci-francia határon. A
gyorsítóban igen nagy energiájú (7000 GeV) protonnyaláb rohan szembe egy ugyanolyan
energiájú másik protonnyalábbal. A két nyaláb egymással szemben, külön csövekben
kering, ám a kör mentén négy helyen keresztezik egymást. Ezeken a helyeken történnek
meg az ütközések, és ide építették a fizikusok és mérnökök a hatalmas méretű detekto-
rokat is. A gyorsítót úgy építették, hogy nemcsak proton-proton ütközéseket lehet vele
vizsgálni, hanem lehetővé teszi nehézionok (például ólomionok) gyorsítását és egymással
való ütköztetését is.
Az LHC-be bonyolult úton kerülnek a részecskék. Amikor protonnyalábot állítanak elő, akkor először egy lineáris gyorsító
gyorsítja a protoncsomagokat. A nehézionok a Linac 3-ból jönnek, amikor ezek gyorsítása folyik. A Linac 4-ből a protonok 250 MeV energiával kerülnek be a PS-Booster
\begin{figure}[H]
    \centering
    \includegraphics[width=0.8\textwidth]{imgs/13-tetel/CERN_LHC.png}
    \caption{A CERN gyorsítókomplexumának részlete}
\end{figure}
nevű gyorsítóba, amely már 1,4 GeV energiára gyorsítja őket. Innen kerülnek tovább a
proton-szinkrotronba (PS), ahonnan 25 GeV energiával kerülnek tovább a szuper-proton
szinkrotronba (SPS). Ezt a gyorsítót végül 450 GeV energiával hagyják el, és kezdik el
útjukat a Large Hadron Collider-ben (LHC). Az LHC-t először ”feltöltik” protoncsomagokkal, amelyek egymással szemben szaladnak két különálló nyalábvezetékben.
A csövekben egymással szemben 2808 protoncsomag szalad, csomagonként mintegy 100 milliárd
protonnal. Az LHC ”feltöltése” 4 perc 20 másodpercig tart. Ez alatt az idő alatt nem történik gyorsítás, csak ”tárolják” a beérkező protoncsomagokat (”tárológyűrű” üzemmód).
Ez után kezdődik meg a protonok gyorsítása. Mintegy 20 perc alatt a protonok energiáját
450 GeV-ról 7000 GeV-re (7 TeV) növelik (”gyorsító” üzemmód). A felgyorsított protonok még órákig keringenek az LHC nyalábvezetőiben (”tárológyűrű” üzemmód ismét),
másodpercenként 11246-szor futnak körbe a mintegy 27 km kerületű gyorsítóban. Útjuk
legnagyobb részét különálló csövekben teszik meg. A kör kerülete mentén azonban van
négy pont, ahol lehetőség van az egymással szemben rohanó protonnyalábok összeütköztetésére. Az ütközési pontokban a protonnyalábokat 16 mikron átmérőre fókuszálják. A
néhány cm hosszú, hajszálnál is vékonyabb protoncsomagok találkozását nagyobb technikai bravúr megvalósítani, mintha két hajszálat akarnánk ”hosszában” összeütköztetni.
Az LHC-ban másodpercenként 800 millió ütközés várható, és ezekből olyan mennyiségű adathalmaz keletkezik, amelynek a feldolgozásához a számítástechnikai módszereknek is
ugrásszerű fejlődésére van szükség.

Az ütközési pontok köré minden eddiginél nagyobb detektorrendszereket építettek a
fizikusok és mérnökök, hogy az összesen 14 TeV ütközési energiából keletkező részecskék
záporát érzékelni tudják. A négy legnagyobb detektor neve: ATLAS, ALICE, CMS,
LHCb. A CMS detektort mutatja építés közben a 11.21 kép. A detektorok óriási méretének oka kettős. Egyrészt a hatalmas energiájú nyalábok ütközésekor óriási energiájú
részecskék is keletkeznek. Ezeknek nagy energiájuk miatt nagy az áthatolóképességük,
tehát csak nagy detektorokkal lehet őket megfogni. Másrészt pedig a nagy ütközési energia miatt nagyon sok (esetleg több száz) részecske is keletkezik egyetlen ütközés során.
Ahhoz, hogy ezeket mind külön-külön érzékelni lehessen, ugyancsak sok elemből álló
detektorrendszerekre van szükség.

Az óriási energiájú ütközésekben keletkezett részecskék között vannak olyanok, amelyek természetes úton talán csak hatalmas csillagkatasztrófák (szupernóvák), vagy az
Univerzum létrejöttekor keletkeztek. Ezáltal a földi gyorsító-berendezések közelebb visznek bennünket a Világegyetem működésének a megértéséhez is.
A kutatók azt várják, hogy ezekből a kísérletekből a részecskefizika egyes olyan alapkérdéseire is választ kapnak (pl. a Higgs-bozon léte, vagy a kvark-gluon plazma tulajdonságai),
amelyek az Univerzum kezdeti pillanataiban irányították az akkor lezajló hatalmas energiájú folyamatokat, és amelyek végső soron megszabták, hogy a világunk
miért éppen olyan, mint amilyennek ma látjuk.

























\end{document}